\appendix

\section{Proof Details}\label{sec:proof-details}

\subsection{Quotient and Cartesian-product structure}\label{app:step-001}

\begin{lemma}[Evaluation-quotient invariance]\label{lem:step-001-quotient-invariance}
Under Assumptions~\ref{assump:vc-one-factors} and
\ref{assump:countably-coded-evaluation}, define
\[
\mathsf P_i:\mathcal H_i\longrightarrow\{0,1\}^{X_i},
\qquad \mathsf P_i(\bar h_i)=\bar h_i\circ\kappa_i.
\]
Then \(\mathsf P_i:\bar C_i\to C_i\) is a bijection.  For every
finite tuple \(x_1,\ldots,x_m\in X_i\),
\begin{equation}\label{eq:step-001-patterns}
\begin{split}
&\{(\bar c_i(\kappa_i(x_1)),\ldots,
        \bar c_i(\kappa_i(x_m))):\bar c_i\in\bar C_i\}\\
&\hspace{25mm}=
\{(c_i(x_1),\ldots,c_i(x_m)):c_i\in C_i\}.
\end{split}
\end{equation}
The analogous equality holds for quotient tuples after choosing raw
representatives.  Raw and quotient classes shatter the same finite VC
patterns and the same finite depths of Littlestone trees.  Consequently,
\[
\operatorname{VC}(\bar C_i)=\operatorname{VC}(C_i)=1,
\qquad
\operatorname{LD}(\bar C_i)=\operatorname{LD}(C_i)=d_i.
\]
\end{lemma}

\begin{proof}
By the definition of \(\equiv_i\), every \(c_i\in C_i\) is constant on
each fiber of \(\kappa_i\).  It therefore has a unique representative
\(\bar c_i:Q_i\to\{0,1\}\) satisfying
\(c_i=\bar c_i\circ\kappa_i\), which proves surjectivity of the displayed
pullback on \(\bar C_i\).  If \(\bar c_i\ne\bar c_i'\), the two functions
differ at some \(q\in Q_i\).  Since \(\kappa_i\) is onto, a point
\(x\in\kappa_i^{-1}(\{q\})\) witnesses that their pullbacks differ.
Thus the pullback is injective.

Equation~\eqref{eq:step-001-patterns} follows from this bijection.  For a
finite quotient tuple, choose one raw representative for each occurrence
and apply the same equality.  This is only finite set-theoretic choice and
does not construct a measurable section.  Quotient collisions cause no
loss: if \(\kappa_i(x_j)=\kappa_i(x_\ell)\), every concept gives the same
label at \(x_j\) and \(x_\ell\).  A pattern assigning different labels to
those occurrences is absent on both sides, and every consistent pattern is
present on one side precisely when it is present on the other.  Hence the
VC dimensions agree.

For Littlestone dimension, replace every raw node label \(x\) in a finite
complete binary instance tree by \(\kappa_i(x)\).  Every root-to-leaf bit
sequence remains realized by the quotient representative of the raw
concept that realized it.  A repeated quotient label with conflicting bits
on one path cannot occur in a shattered raw tree because equivalent points
have identical labels under every concept.  Conversely, a finite quotient
tree has finitely many nodes; choose a raw representative for each node
label.  Pullback then realizes every representative path.  Thus the two
classes shatter the same finite tree depths, and Assumption~\ref{assump:vc-one-factors}
gives the stated dimensions.
\end{proof}

\begin{lemma}[VC dimension and finite cardinality of the full product]
\label{lem:step-001-vc-cardinality}
Under Assumptions~\ref{assump:canonical-product} and
\ref{assump:vc-one-factors},
\[
\operatorname{VC}(C)=k,
\qquad |C|=\prod_{i=1}^k|C_i|.
\]
If every factor is finite, then
\(\log|C|=\sum_{i=1}^k\log|C_i|\).  If at least one factor is infinite,
then \(C\) is infinite and no finite-real logarithmic identity is asserted.
\end{lemma}

\begin{proof}
For each factor choose \(x_i\in X_i\) whose singleton is shattered.  Given
\(b=(b_1,\ldots,b_k)\in\{0,1\}^k\), choose a factor concept taking value
\(b_i\) at \(x_i\).  Assumption~\ref{assump:canonical-product} combines
these factor choices into a global concept, so
\(\{x_1,\ldots,x_k\}\) is shattered.  Conversely, any \(k+1\) distinct
points contain two in a common block.  Shattering the whole set would
shatter those two points in the corresponding factor, contradicting its VC
dimension one.  This proves \(\operatorname{VC}(C)=k\).

The restriction bijection in Assumption~\ref{assump:canonical-product}
gives the cardinal product.  For finite factors, taking the fixed natural
logarithm gives the sum identity.  If one factor is infinite, fixing one
concept in every other nonempty factor injects it into \(C\), so \(C\) is
infinite.  Constant-valued concepts are allowed.  In particular, when a
quotient has one cell and the factor contains both constant labels, that
cell is still shattered.  A genuinely singleton factor is excluded by
Assumption~\ref{assump:vc-one-factors}.
\end{proof}

\begin{lemma}[Rank predictor for a finite-Littlestone class]
\label{lem:step-001-rank-predictor}
Let \(H\subseteq\{0,1\}^Y\) be nonempty with
\(\operatorname{LD}(H)=d<\infty\).  There is a deterministic online
predictor making at most \(d\) mistakes on every sequence realizable by one
member of \(H\).  The construction requires neither an order nor a
cardinality assumption on \(Y\) or \(H\), and it selects no hypothesis.
\end{lemma}

\begin{proof}
Set \(\operatorname{LD}(\varnothing)=-1\).  At a query \(x\), let
\(V\ne\varnothing\) be the current version class and write
\[
V_b(x)=\{h\in V:h(x)=b\},\qquad b\in\{0,1\}.
\]
If \(a=\operatorname{LD}(V)\), both children cannot have Littlestone
dimension at least \(a\), since adjoining the query \(x\) above depth-
\(a\) trees for both children would give a depth-\(a+1\) tree shattered by
\(V\).  Hence
\begin{equation}\label{eq:step-001-rank-drop}
\min_{b\in\{0,1\}}\operatorname{LD}(V_b(x))\le a-1.
\end{equation}
Predict a bit whose child has larger rank, breaking ties in favor of zero.
On a mistake, the realized child is a smaller-ranked child; in a tie both
children satisfy the minimum bound.  Thus every mistake lowers the rank by
at least one.  Correct predictions pass to a subclass and cannot increase
rank.  Starting at rank \(d\), at most \(d\) mistakes occur.  At rank zero,
two nonempty children would shatter a depth-one tree, so the realizable next
label cannot cause another mistake.
\end{proof}

\begin{proposition}[Additivity of Littlestone dimension]
\label{prop:step-001-product-littlestone}
Under Assumptions~\ref{assump:canonical-product} and
\ref{assump:vc-one-factors},
\[
\operatorname{LD}(C)=\sum_{i=1}^k d_i.
\]
\end{proposition}

\begin{proof}
For each \(i\), choose a complete depth-\(d_i\) tree shattered by \(C_i\).
Attach a copy of the factor-two tree at every leaf of the factor-one tree,
then the factor-three tree at every new leaf, and continue through factor
\(k\).  The resulting depth is \(\sum_i d_i\).  Along a path, its segment
in factor \(i\) is realized by some \(c_i\in C_i\), and
Assumption~\ref{assump:canonical-product} combines these choices into one
global concept.  Therefore
\begin{equation}\label{eq:step-001-ld-lower}
\operatorname{LD}(C)\ge\sum_i d_i.
\end{equation}

For the reverse inequality, run one copy of the predictor in
Lemma~\ref{lem:step-001-rank-predictor} for each factor, updating only the
copy corresponding to the block containing the current instance.  A
sequence realized by \(c\in C\) induces a factor subsequence realized by
\(c|_{X_i}\), so copy \(i\) makes at most \(d_i\) mistakes and the global
predictor makes at most \(\sum_i d_i\).  If \(C\) shattered a tree of depth
\(\sum_i d_i+1\), traverse it by taking at each node the edge opposite the
global prediction.  Shattering supplies one concept realizing the resulting
path, on which the predictor makes a mistake at every level, a contradiction.
Thus
\begin{equation}\label{eq:step-001-ld-upper}
\operatorname{LD}(C)\le\sum_i d_i.
\end{equation}
The two bounds prove the claim.  A zero-dimensional singleton factor would
contribute neither tree depth nor mistakes, while the present VC-one factors
have \(d_i\ge1\).  At \(k=1\), both bounds are the one-factor identity.
\end{proof}

\begin{lemma}[Log-star and aggregate-size comparisons]
\label{lem:step-001-logstar}
Under Assumption~\ref{assump:vc-one-factors}, let
\[
r_i=\log_2^*d_i,\qquad
s_i=1+\log_2^*(d_i+1),\qquad
M=\sum_{i=1}^k s_i=M_{\oplus}(C).
\]
Then, for every \(i\),
\[
s_i\ge2,\qquad s_i\le r_i+2,\qquad M\ge2k.
\]
\end{lemma}

\begin{proof}
VC dimension one gives a depth-one Littlestone tree, so \(d_i\ge1\).
Consequently \(d_i+1\ge2\), at least one base-two logarithm is needed to
reach a value at most one, and \(s_i\ge2\).  The stopping time
\(\log_2^*t\) is nondecreasing on \([1,\infty)\).  For each integer
\(d\ge1\),
\begin{equation}\label{eq:step-001-log-elementary}
\log_2(d+1)\le d,
\end{equation}
because \(2^d\ge d+1\), with equality at \(d=1\) and induction thereafter.
The log-star recursion and monotonicity yield
\begin{equation}\label{eq:step-001-logstar-comparison}
\log_2^*(d_i+1)
=1+\log_2^*\!\bigl(\log_2(d_i+1)\bigr)
\le1+\log_2^*d_i=r_i+1.
\end{equation}
Adding one proves \(s_i\le r_i+2\), and summing \(s_i\ge2\) proves
\(M\ge2k\).  At the smallest allowed dimension,
\[
d_i=1,\qquad r_i=0,\qquad s_i=2,
\]
so the endpoint is retained exactly.
\end{proof}

\begin{lemma}[Standard-Borel output and measurable decoding]
\label{lem:step-001-output-measurability}
Under Assumptions~\ref{assump:canonical-product} and
\ref{assump:countably-coded-evaluation}, every
\((\mathcal H_i,\mathscr H_i)=(\{0,1\}^{Q_i},\mathscr H_i)\) and their
finite product \((\mathcal H^{\oplus},\mathscr H^{\oplus})\) are standard
Borel.  The map
\[
E:\mathcal H^{\oplus}\times X\to\{0,1\},
\qquad E(\bar h,x)=\bar h_i(\kappa_i(x))\quad(x\in X_i),
\]
is jointly measurable.  Hence every decoded \(h_{\bar h}=E(\bar h,\cdot)\)
is \(\Sigma\)-measurable, and decoding is measurable for every finite
evaluation cylinder specified in the learning model.
\end{lemma}

\begin{proof}
If \(Q_i\) is finite, \(\mathcal H_i\) is finite discrete.  If it is
countably infinite, enumerate it as \(q_{i,1},q_{i,2},\ldots\) and define
\begin{equation}\label{eq:step-001-product-metric}
\Delta_i(\bar h,\bar g)
=\sum_{m=1}^{\infty}2^{-m}
\mathbf1\{\bar h(q_{i,m})\ne\bar g(q_{i,m})\}.
\end{equation}
This metric induces the product topology: a finite coordinate cylinder is
clopen, while agreement on a sufficiently long initial segment makes the
remaining metric tail arbitrarily small.  It is complete because a Cauchy
sequence stabilizes coordinatewise and converges to its coordinatewise
limit.  It is separable because functions vanishing beyond a finite initial
segment form a countable dense subset.  Thus the topology is Polish, and
its Borel sigma-field is exactly the finite-evaluation-cylinder sigma-field
\(\mathscr H_i\).

The disjoint coordinate set
\(
\widetilde Q=\bigsqcup_{i=1}^k(\{i\}\times Q_i)
\)
is finite or countable.  Coordinate relabeling identifies
\(\mathcal H^{\oplus}\) with \(\{0,1\}^{\widetilde Q}\), and the two
sigma-fields have the same coordinate generators.  The preceding argument
therefore makes the finite product standard Borel.

For joint evaluation,
\begin{equation}\label{eq:step-001-joint-evaluation}
E^{-1}(\{1\})
=\bigcup_{i=1}^k\ \bigcup_{q\in Q_i}
\{\bar h:\bar h_i(q)=1\}
\times\kappa_i^{-1}(\{q\}).
\end{equation}
The output sets are cylinders, the raw cells are measurable, and the union
is countable.  For fixed \(\bar h\),
\[
h_{\bar h}^{-1}(\{1\})
=\bigcup_{i=1}^k\ \bigcup_{q:\bar h_i(q)=1}
\kappa_i^{-1}(\{q\})\in\Sigma.
\]
Finally, for fixed \(x_1,\ldots,x_m\) and \(b\in\{0,1\}^m\), the decoder
preimage of the corresponding evaluation cylinder is
\begin{equation}\label{eq:step-001-decoder-cylinder}
\bigcap_{j=1}^m
\{\bar h:\bar h_{i(j)}(\kappa_{i(j)}(x_j))=b_j\},
\end{equation}
where \(i(j)\) is the unique block containing \(x_j\).  This proves exactly
the required decoder measurability without imposing a sigma-field on the
entire raw hypothesis collection.
\end{proof}

\begin{lemma}[Exact measurable quotient-to-raw risk pullback]
\label{lem:step-001-risk-pullback}
Under Assumptions~\ref{assump:canonical-product} and
\ref{assump:countably-coded-evaluation} and
Lemma~\ref{lem:step-001-quotient-invariance}, fix \(c\in C\), a probability
measure \(D\) on \((X,\Sigma)\), and \(\bar h\in\mathcal H^{\oplus}\).
For \(\rho_i=D(X_i)>0\),
\begin{equation}\label{eq:step-001-factor-risk}
\begin{split}
\Pr_{x\sim D_i}[\mathsf P_i(\bar h_i)(x)\ne c_i(x)]
&=\Pr_{q\sim\bar D_i}[\bar h_i(q)\ne\bar c_i(q)]\\
&=\sum_{q\in Q_i}D_i(\kappa_i^{-1}(\{q\}))
\mathbf1\{\bar h_i(q)\ne\bar c_i(q)\}.
\end{split}
\end{equation}
Globally,
\begin{equation}\label{eq:step-001-global-risk}
\begin{split}
R_D(h_{\bar h},c)
&=\sum_{i:\rho_i>0}\rho_i
\Pr_{q\sim\bar D_i}[\bar h_i(q)\ne\bar c_i(q)]\\
&=\sum_{i=1}^k\sum_{q\in Q_i}
D(\kappa_i^{-1}(\{q\}))
\mathbf1\{\bar h_i(q)\ne\bar c_i(q)\}.
\end{split}
\end{equation}
These risks are measurable functions of the quotient output.  The identities
hold for arbitrary-support \(D\), and a block of mass zero contributes zero
without requiring \(D_i\) to be defined.
\end{lemma}

\begin{proof}
Lemma~\ref{lem:step-001-quotient-invariance} gives
\(c_i(x)=\bar c_i(\kappa_i(x))\) and
\(\mathsf P_i(\bar h_i)(x)=\bar h_i(\kappa_i(x))\).  The raw mistake
indicator is therefore the pullback of the quotient mistake indicator.
Pushforward by \(\kappa_i\) gives the first equality in
\eqref{eq:step-001-factor-risk}; discreteness and countable additivity give
the second.

The measurable blocks partition \(X\), so the global mistake set is the
disjoint union of its block intersections.  Countable additivity across the
finite blocks and then their quotient cells gives the second line of
\eqref{eq:step-001-global-risk}.  For positive block mass,
\[
D(\kappa_i^{-1}(\{q\}))
=\rho_iD_i(\kappa_i^{-1}(\{q\})),
\]
which gives the first line.  If \(\rho_i=0\), every cell in \(X_i\) has
zero mass, so the contribution vanishes without a conditional law.

Each summand is a fixed nonnegative coefficient times the indicator of an
evaluation cylinder \(\{\bar h:\bar h_i(q)\ne\bar c_i(q)\}\).  Finite
partial sums are measurable and increase pointwise to the displayed risks.
Their limits are measurable, as is the finite sum over factors.  Hence every
risk threshold event used below is measurable, and no approximation term is
introduced by quotient decoding.
\end{proof}

\begin{proposition}[Quotient-product structural certificate]
\label{prop:step-001-structural-certificate}
Under Assumptions~\ref{assump:canonical-product},
\ref{assump:vc-one-factors}, and
\ref{assump:countably-coded-evaluation}, quotient pullback preserves every
factor's VC and Littlestone dimensions; the global class satisfies
\[
\operatorname{VC}(C)=k,\qquad
\operatorname{LD}(C)=\sum_{i=1}^k d_i,\qquad
s_i\ge2,\quad s_i\le\log_2^*d_i+2,\quad M_{\oplus}(C)\ge2k,
\]
and, for finite factors, \(\log|C|=\sum_i\log|C_i|\).  Moreover
\(\mathcal H^{\oplus}\) is standard Borel, decoding yields measurable raw
hypotheses, and quotient, raw-factor, and global risks agree exactly as in
\eqref{eq:step-001-factor-risk}--\eqref{eq:step-001-global-risk}.
\end{proposition}

\begin{proof}
The quotient and factor-dimension statements are
Lemma~\ref{lem:step-001-quotient-invariance}.  Product VC and the qualified
cardinality formula are Lemma~\ref{lem:step-001-vc-cardinality}.  The rank
argument in Lemma~\ref{lem:step-001-rank-predictor} and
Proposition~\ref{prop:step-001-product-littlestone} gives product
Littlestone additivity.  Lemma~\ref{lem:step-001-logstar} supplies every
scalar comparison, including \(d_i=1\).  Output legality follows from
Lemma~\ref{lem:step-001-output-measurability}, and
Lemma~\ref{lem:step-001-risk-pullback} gives the exact measurable risk
identities, including countably infinite quotients, arbitrary-support
distributions, and zero-mass blocks.  At \(k=1\), each product, output, and
risk identity is precisely its sole-factor version.
\end{proof}

\subsection{A measurable totalized private learner for one factor}
\label{app:step-002}

\begin{lemma}[Reference-transformed quotient order geometry]
\label{lem:step-002-order-geometry}
Under Assumptions~\ref{assump:vc-one-factors} and
\ref{assump:countably-coded-evaluation} and
Lemma~\ref{lem:step-001-quotient-invariance}, fix a factor \(i\) and the
pre-data reference \(\bar f_i^\circ\in\bar C_i\).  There is a finite-height
partially ordered active transformed quotient \((U_i^+,\preceq_i)\) with
the following properties.  The transformed class has VC dimension one and
Littlestone dimension \(d_i\); every nonempty transformed positive set and
every nonempty intersection of such sets is a finite principal upset; if
\(H_i\) is the maximum strict-chain length, then
\begin{equation}\label{eq:step-002-height}
1\le H_i\le \operatorname{TD}(\bar C_i\mathbin\triangle\bar f_i^\circ)
\le2^{d_i+1};
\end{equation}
and each principal upset contains exactly one point of every integer depth
between one and the depth of its minimum.
\end{lemma}

\begin{proof}
For \(\bar c\in\bar C_i\), set
\begin{equation}\label{eq:step-002-xor-class}
g_{\bar c}(q)=\bar c(q)\oplus\bar f_i^\circ(q),
\qquad \mathcal G_i=\{g_{\bar c}:\bar c\in\bar C_i\}.
\end{equation}
XOR by a fixed function is a bijection on every finite label vector and on
every binary instance tree, with edge bits XORed nodewise.  Hence
\begin{equation}\label{eq:step-002-xor-dimensions}
\operatorname{VC}(\mathcal G_i)=1,\qquad
\operatorname{LD}(\mathcal G_i)=d_i,\qquad
0=g_{\bar f_i^\circ}\in\mathcal G_i.
\end{equation}
The transform may identify coordinates, so define
\[
q\sim_iq'\quad\Longleftrightarrow\quad
g(q)=g(q')\ \text{for every }g\in\mathcal G_i,
\]
let \(U_i=Q_i/{\sim_i}\), and write \(\tau_i:Q_i\to U_i\).  This quotient
is finite or countable.  The two-way finite-representative argument in
Lemma~\ref{lem:step-001-quotient-invariance} shows that quotienting preserves
all finite patterns and finite Littlestone trees, so
\eqref{eq:step-002-xor-dimensions} remains true on \(U_i\).

Let
\[
U_i^+=\{u\in U_i:\text{some }g\in\mathcal G_i\text{ has }g(u)=1\}.
\]
It is nonempty because the factor is nonconstant and the reference is one
of its concepts.  Put
\begin{equation}\label{eq:step-002-implication-order}
u\preceq_i v
\quad\Longleftrightarrow\quad
\bigl(g(u)=1\Longrightarrow g(v)=1\bigr)
\quad\text{for every }g\in\mathcal G_i.
\end{equation}
This relation is reflexive and transitive.  Antisymmetry follows from the
secondary evaluation quotient: mutual implication means all transformed
concepts agree at the two coordinates.

If \(u,v\) are incomparable, failures of the two implications provide the
patterns \((1,0)\) and \((0,1)\), while the zero concept provides
\((0,0)\).  A concept taking value one at both points would shatter the pair,
contrary to VC dimension one.  Therefore every positive set
\begin{equation}\label{eq:step-002-positive-set}
P_g=\{u\in U_i^+:g(u)=1\}
\end{equation}
is both a chain and an upset.

Every strict chain \(u_1\prec_i\cdots\prec_i u_r\) yields a threshold
pattern of length \(r\).  Active membership supplies the all-one suffix
starting at \(u_1\); for \(j\ge2\), strictness supplies a concept zero at
\(u_{j-1}\) and one at \(u_j\), and implication makes it zero before \(j\)
and one from \(j\) onward.  The threshold-dimension comparison
\(\operatorname{TD}(H)\le2^{\operatorname{LD}(H)+1}\) for a binary class
\citep{yan2025vc1} proves \eqref{eq:step-002-height}.  Thus every chain is
finite and the largest integer chain length exists.  No tree-construction
or countable-set selector is used.

For \(u\in U_i^+\), define
\begin{equation}\label{eq:step-002-upset-depth}
I_i(u)=\{v\in U_i^+:u\preceq_i v\},
\qquad \ell_i(u)=|I_i(u)|.
\end{equation}
A transformed concept positive at \(u\) is positive throughout \(I_i(u)\),
so the incomparability argument makes \(I_i(u)\) a chain of at most
\(H_i\) points.  It is the longest chain starting at \(u\), and successive
points in its finite total order have suffix sizes
\(\ell_i(u),\ell_i(u)-1,\ldots,1\).  This proves uniqueness at every depth.

Finally, for nonempty \(V\subseteq\mathcal G_i\), set
\begin{equation}\label{eq:step-002-core}
B(V)=\bigcap_{g\in V}P_g.
\end{equation}
It is an upset contained in the finite chain \(P_{g_0}\) for any
\(g_0\in V\).  If nonempty it has a unique minimum \(u_V\), and upward
closure gives \(B(V)=I_i(u_V)\).  This proves every asserted structural
property.
\end{proof}

\begin{lemma}[VC dimension of the improper version-core class]
\label{lem:step-002-core-vc}
Under Lemma~\ref{lem:step-002-order-geometry}, define, with every function
extended by zero on \(U_i\setminus U_i^+\),
\begin{equation}\label{eq:step-002-core-class}
\mathcal K_i=
\{\mathbf1_{B(V)}:\varnothing\ne V\subseteq\mathcal G_i\}\cup\{0\}.
\end{equation}
Then \(\operatorname{VC}(\mathcal K_i)\le1\), every transformed target is
in \(\mathcal K_i\), and the totalized version-core indicator of every
realizable labeled block belongs to \(\mathcal K_i\) and has zero empirical
error against its target.
\end{lemma}

\begin{proof}
For two distinct active points \(u,v\), if \(u\preceq_i v\), no upset
realizes \((1,0)\).  If they are incomparable, every nonempty core is
contained in one defining positive set, and the argument in
Lemma~\ref{lem:step-002-order-geometry} shows that no such set contains both
points; hence no core realizes \((1,1)\).  Adding zero cannot restore the
missing pattern.  A pair containing an inactive coordinate also cannot be
shattered, so the VC dimension is at most one.

For \(g_*\in\mathcal G_i\), its positive set is \(B(\{g_*\})\), unless
empty, in which case \(g_*=0\).  Thus every target is in the class.  For a
finite transformed block \(T\), let
\begin{equation}\label{eq:step-002-version-space}
V(T)=\{g\in\mathcal G_i:g\text{ agrees with every record of }T\}.
\end{equation}
Set its core to \(B_T=\varnothing\) when \(V(T)=\varnothing\), and otherwise
to \(B(V(T))\).  On a realizable block, \(g_*\in V(T)\).  A positive record
belongs to every consistent concept's positive set and therefore to
\(B_T\); a negative record is excluded by the consistent target.  Hence
\(\mathbf1_{B_T}\) agrees with every record.  The argument never asserts
that this core indicator is a proper concept.
\end{proof}

\begin{lemma}[Countable positive-support replacement choosing rule]
\label{lem:step-002-countable-choosing}
Let \(F\) be finite or countable, let
\(\mathsf O=F\cup\{\bot\}\) have the discrete sigma-field, and let
\(D=(v_1,\ldots,v_t)\in\mathsf O^t\).  For
\(0<b<1\), \(0<\varepsilon_c\le2\), and \(0<\delta_c<1\), define
\[
q(D,f)=\sum_{j=1}^t\mathbf1\{v_j=f\},\qquad
\operatorname{OPT}(D)=\sup_{f\in F}q(D,f),
\]
with optimum zero when \(F=\varnothing\).  Apply the thresholded choosing
kernel of \citet{bunNissimStemmerVadhan2015} with bounded-growth parameter
one: set
\(\widetilde{\operatorname{OPT}}=\operatorname{OPT}+\operatorname{Lap}(4/\varepsilon_c)\),
return \(\bot\) if
\[
\widetilde{\operatorname{OPT}}
<\frac8{\varepsilon_c}
\log\!\left(\frac4{b\varepsilon_c\delta_c}\right),
\]
and otherwise use its \(\varepsilon_c/2\)-exponential stage on
\(G(D)=\{f:q(D,f)>0\}\).  If \(G(D)=\varnothing\), return \(\bot\)
immediately.  The resulting rule is a replacement-
\((\varepsilon_c,\delta_c)\)-DP Markov kernel on the countable spaces.
Defining \(q(D,\bot)=0\), with probability at least \(1-b\),
\begin{equation}\label{eq:step-002-choosing-utility}
q(D,\hat f)\ge \operatorname{OPT}(D)
-\frac{16}{\varepsilon_c}
\log\!\left(\frac{4t}{b\varepsilon_c\delta_c}\right)
\end{equation}
whenever the right side is positive; then \(\hat f\ne\bot\).
\end{lemma}

\begin{proof}
Under addition adjacency, one added non-sentinel summary increments exactly
one score, while a sentinel changes none, so the score is one-bounded-growth.
For fixed-size replacement, neighboring databases share \(t-1\) entries and
have one leaving and one entering entry, whence
\begin{equation}\label{eq:step-002-histogram-replacement}
q(D',f)-q(D,f)
=\mathbf1\{f=v_{\mathrm{in}}\}-\mathbf1\{f=v_{\mathrm{out}}\}.
\end{equation}
Every score and the optimum have sensitivity one, and each directional set
difference between \(G(D)\) and \(G(D')\) has size at most one.  The privacy
proof for the cited choosing mechanism uses exactly these facts and sums
only over the positive supports, each of which contains at most \(t\) points
\citep{bunNissimStemmerVadhan2015}.  Its finite sums therefore remain valid
when the ambient solution set is countable.  The same observation preserves
the cited utility bound \eqref{eq:step-002-choosing-utility}.

For each fixed database, the thresholded and finite-support stages define a
probability measure.  Since the database and output spaces are countable
discrete, every transition-probability map is measurable.

It remains to check the newly totalized empty-support branch.  Suppose
\(G(D_0)=\varnothing\) and \(D_0,D_1\) are replacement neighbors.  If the
second support is also empty, both laws are \(\delta_\bot\).  Otherwise
\eqref{eq:step-002-histogram-replacement} gives
\(G(D_1)=\{f_0\}\) and \(q(D_1,f_0)=\operatorname{OPT}(D_1)=1\).  Since
\[
1<\frac4{\varepsilon_c}
\log\!\left(\frac4{b\varepsilon_c\delta_c}\right),
\]
passing the noisy threshold requires a Laplace fluctuation exceeding the
right side.  Thus, for \(p=\Pr[A(D_1)\ne\bot]\),
\begin{equation}\label{eq:step-002-empty-support-probability}
p\le\frac12\exp\!\left(-\log\frac4{b\varepsilon_c\delta_c}\right)
=\frac{b\varepsilon_c\delta_c}{8}\le\delta_c.
\end{equation}
If \(\bot\notin E\), the \(D_1\)-probability of \(E\) is at most
\(\delta_c\) and the \(D_0\)-probability is zero, proving both DP
directions.  If \(\bot\in E\), one direction follows from
\(1\le e^{\varepsilon_c}+\delta_c\); for the other,
\(\Pr[A(D_1)\in E]\ge1-\delta_c\) and
\begin{equation}\label{eq:step-002-empty-support-reverse}
1\le e^{\varepsilon_c}(1-\delta_c)+\delta_c
\le e^{\varepsilon_c}\Pr[A(D_1)\in E]+\delta_c.
\end{equation}
Relabeling the neighbors covers a transition into empty support.  Hence the
total rule has the stated privacy for every event and sentinel transition.
If the right side of \eqref{eq:step-002-choosing-utility} is positive,
output \(\bot\), whose score is zero, cannot satisfy it.
\end{proof}

\begin{lemma}[Explicit internal calibration and quota domination]
\label{lem:step-002-calibration}
Under Assumptions~\ref{assump:vc-one-factors} and
\ref{assump:global-privacy-range}, Lemma~\ref{lem:step-001-logstar}, and
Lemma~\ref{lem:step-002-order-geometry}, let \(C_{PM}>0\) be the universal
constant in the following fixed-accuracy private-median consequence of
\citet{yan2025vc1}: on a finite ordered domain of size \(m\ge2\), a
replacement-\((\varepsilon_m,\delta_m)\)-DP kernel returns a
\(1/3\)-median except with probability \(\eta\) whenever its input count is
strictly larger than
\[
C_{PM}\frac{\log_2^*m}{\varepsilon_m}
\log^2\!\left(\frac{e\log_2^*m}{\eta\delta_m}\right).
\]
Set
\begin{equation}\label{eq:step-002-fixed-calibration}
\eta=\frac1{16384},\quad \gamma=\frac1{48},\quad
\varepsilon_m=\varepsilon_c=\frac\varepsilon4,\quad
\delta_m=\delta_c=\frac\delta4,
\end{equation}
and
\begin{equation}\label{eq:step-002-block-size}
L_*=\left\lceil
\frac{48}{1/64}
\left(10\log\frac{48e}{1/64}+\log\frac5\gamma\right)
\right\rceil.
\end{equation}
There are universal constants \(C_T,K_*\) such that, for
\begin{equation}\label{eq:step-002-counts}
\Lambda_i=\log\frac{es_i}{\varepsilon\delta},\qquad
\Xi_i=\frac{s_i}{\varepsilon}\Lambda_i^2,
\qquad t_i=\lceil C_T\Xi_i\rceil,\qquad N_i=L_*t_i,
\end{equation}
the count \(t_i\) exceeds that private-median threshold on
\(\{0,\ldots,H_i\}\), and
\begin{equation}\label{eq:step-002-calibration-conclusions}
e^{-t_i/128}\le\eta,\qquad
\frac{16}{\varepsilon_c}
\log\!\left(\frac{4t_i}{\eta\varepsilon_c\delta_c}\right)
\le\frac{t_i}{12}.
\end{equation}
If the setting constant satisfies \(K_Y\ge K_*\), then
\begin{equation}\label{eq:step-002-quota-domination}
N_i\le
\left\lceil K_Y\frac{s_i}{\varepsilon}
\log^2\!\left(\frac{es_i}{\varepsilon\delta}\right)\right\rceil=q_i.
\end{equation}
All statements include \(d_i=1\).
\end{lemma}

\begin{proof}
Let
\[
\Lambda_0=\log(20e),\qquad
c_1=1+\frac{\log((3/5)/\eta)}{\Lambda_0}.
\]
Choose \(C_T\) as the least positive integer satisfying
\begin{equation}\label{eq:step-002-CT-conditions}
\begin{split}
C_T&>6C_{PM}c_1^2,\\
20C_T\Lambda_0^2&\ge128\log(1/\eta),\\
C_T&\ge\frac{384}{\Lambda_0}
\left(3+\frac{\log(64(C_T+1)/(\eta e))}{\Lambda_0}\right).
\end{split}
\end{equation}
Such an integer exists because the last right side grows logarithmically
while the left side grows linearly.  Set
\begin{equation}\label{eq:step-002-Kstar}
K_*=L_*(C_T+1).
\end{equation}

Writing \(u_i=\log_2^*(H_i+1)\), the height bound gives
\(H_i+1\le2^{d_i+1}+1\).  Log-star recursion and monotonicity,
\(d_i+2\le2(d_i+1)\), and
\(\log_2^*(2x)\le1+\log_2^*x\) for \(x\ge2\) imply
\begin{equation}\label{eq:step-002-height-logstar}
u_i\le1+\log_2^*(d_i+2)
\le2+\log_2^*(d_i+1)=s_i+1\le\frac32s_i.
\end{equation}
At \(d_i=1\), this says \(u_i\le3\), while \(s_i=2\) and
\(H_i+1\le5\); no zero-domain endpoint is used.

The privacy range and \(s_i\ge2\) give
\begin{equation}\label{eq:step-002-Lambda-lower}
\Lambda_i\ge\Lambda_0,\qquad
\Xi_i\ge20\Lambda_0^2>1.
\end{equation}
Moreover,
\begin{equation}\label{eq:step-002-median-log}
\begin{split}
\log\frac{eu_i}{\eta\delta_m}
&=\log\frac{4eu_i}{\eta\delta}
=\Lambda_i+\log\frac{4u_i\varepsilon}{\eta s_i}\\
&\le\Lambda_i+\log\frac{3/5}{\eta}
\le c_1\Lambda_i.
\end{split}
\end{equation}
The private-median threshold is therefore at most
\begin{equation}\label{eq:step-002-median-threshold}
C_{PM}\frac{u_i}{\varepsilon_m}
\log^2\frac{eu_i}{\eta\delta_m}
\le6C_{PM}c_1^2\Xi_i<C_T\Xi_i\le t_i,
\end{equation}
including its strict sample requirement.  The first two conditions in
\eqref{eq:step-002-CT-conditions} and
\eqref{eq:step-002-Lambda-lower} give
\(t_i\ge128\log(1/\eta)\), proving the first inequality in
\eqref{eq:step-002-calibration-conclusions}.

Since \(t_i\le(C_T+1)\Xi_i\),
\(\log(1/\varepsilon)\le\Lambda_i\), and
\(2\log\Lambda_i\le\Lambda_i\) for \(\Lambda_i\ge\Lambda_0\),
\begin{equation}\label{eq:step-002-choosing-log}
\begin{split}
\log\frac{4t_i}{\eta\varepsilon_c\delta_c}
&\le\log\frac{64(C_T+1)s_i\Lambda_i^2}
                 {\eta\varepsilon^2\delta}\\
&\le\left(3+
\frac{\log(64(C_T+1)/(\eta e))}{\Lambda_0}\right)\Lambda_i.
\end{split}
\end{equation}
Because \(s_i\Lambda_i\ge2\Lambda_0\), the final condition in
\eqref{eq:step-002-CT-conditions} yields
\[
\frac{t_i}{12}
\ge\frac{C_Ts_i\Lambda_i^2}{12\varepsilon}
\ge\frac{64}{\varepsilon}
\log\frac{4t_i}{\eta\varepsilon_c\delta_c}
=\frac{16}{\varepsilon_c}
\log\frac{4t_i}{\eta\varepsilon_c\delta_c}.
\]
Finally, \(\Xi_i>1\) gives
\begin{equation}\label{eq:step-002-internal-size}
N_i=L_*t_i\le L_*(C_T\Xi_i+1)
\le L_*(C_T+1)\Xi_i=K_*\Xi_i.
\end{equation}
Together with \(K_Y\ge K_*\), this proves
\eqref{eq:step-002-quota-domination} including the exact ceiling.
\end{proof}

\begin{proposition}[Totalized quotient transition kernel]
\label{prop:step-002-kernel}
Under Assumptions~\ref{assump:vc-one-factors},
\ref{assump:countably-coded-evaluation}, and
\ref{assump:global-privacy-range}, Lemmas~\ref{lem:step-001-output-measurability}
and \ref{lem:step-001-logstar}, and
Lemmas~\ref{lem:step-002-order-geometry}--\ref{lem:step-002-calibration},
the rule defined below is a fixed permutation-invariant Markov kernel
\[
\bar A_i^{\mathrm{Yan}}:
(Q_i\times\{0,1\})^{q_i}\rightsquigarrow
(\mathcal H_i,\mathscr H_i)
\]
on every input, including inconsistent and nonrealizable inputs.
\end{proposition}

\begin{proof}
Freeze the reference \(\bar f_i^\circ\), the enumeration of \(Q_i\), the
induced first-occurrence enumeration of \(U_i\), implementations of the
private-median and choosing kernels, and every tie and fallback convention
before observing data.  On
\(\bar T=((q_j,b_j))_{j=1}^{q_i}\), draw a uniform permutation, retain the
first \(N_i\le q_i\) records, and split them into \(t_i\) consecutive blocks
of size \(L_*\).  Transform a record by
\begin{equation}\label{eq:step-002-record-transform}
(q,b)\longmapsto(\tau_i(q),b\oplus\bar f_i^\circ(q)).
\end{equation}
For each block \(T_b\), form the version space
\eqref{eq:step-002-version-space}, totalized core \(B_b\), and depth
\begin{equation}\label{eq:step-002-block-depth}
a_b=\begin{cases}
0,&B_b=\varnothing,\\
\ell_i(u_b),&B_b=I_i(u_b)\ne\varnothing.
\end{cases}
\end{equation}
Apply the private median on the fixed ordered set \(\{0,\ldots,H_i\}\) to
obtain \(z\).  At \(z=0\), output \(\bar f_i^\circ\).  For fixed
\(z\ge1\), define
\begin{equation}\label{eq:step-002-layer-summary}
p_z(T_b)=\begin{cases}
\text{the unique }u\in B_b\text{ with }\ell_i(u)=z,&a_b\ge z,\\
\bot,&a_b<z.
\end{cases}
\end{equation}
Use Lemma~\ref{lem:step-002-countable-choosing} on the \(t_i\) summaries
with solution set \(F_z=\{u\in U_i^+:\ell_i(u)=z\}\).  If it returns
\(\bot\), output the reference; if it returns \(\hat u\), output
\begin{equation}\label{eq:step-002-factor-output}
\bar h_{\hat u}(q)=\bar f_i^\circ(q)
\oplus\mathbf1\{\tau_i(q)\in I_i(\hat u)\}.
\end{equation}

Every branch is total: an inconsistent or empty version space has empty
core, an empty core has depth zero, a missing layer has summary \(\bot\),
and empty choosing support or a failed selection returns the fixed reference.
The input space is finite or countable discrete.  Hence for every
\(E\in\mathscr H_i\), the transition-probability function is measurable,
regardless of the set-theoretic form of a version space.  For fixed input,
the finite chain of kernels and deterministic maps gives a probability
measure on the standard-Borel output space from
Lemma~\ref{lem:step-001-output-measurability}.  Finally, precomposing the
input with a fixed permutation merely composes two permutations while the
sampled one remains uniform, proving permutation invariance.
\end{proof}

\begin{proposition}[All-input replacement privacy of the factor kernel]
\label{prop:step-002-factor-privacy}
Under the assumptions and construction of
Proposition~\ref{prop:step-002-kernel}, \(\bar A_i^{\mathrm{Yan}}\) is
replacement-\((\varepsilon/2,\delta/2)\)-DP on every adjacent pair of
length-\(q_i\) inputs, including every nonrealizable and totalization branch.
\end{proposition}

\begin{proof}
Couple the data-independent uniform permutation on two inputs differing in
one record.  If the changed position is not retained, all used blocks agree;
otherwise exactly one block changes by one replacement.  Consequently the
depth vector \eqref{eq:step-002-block-depth} changes in at most one
coordinate.  The calibrated private-median kernel is replacement-
\((\varepsilon_m,\delta_m)\)-DP on arbitrary depth vectors.

Condition on the released depth.  At \(z=0\), the second-stage output is
fixed.  At \(z\ge1\), only the changed block's summary can change, and
\eqref{eq:step-002-histogram-replacement} covers a leaving and entering
point, a transition to or from \(\bot\), and a change between two ordinary
points.  Lemma~\ref{lem:step-002-countable-choosing} supplies replacement-
\((\varepsilon_c,\delta_c)\)-privacy even with empty support.  Sequential
composition and postprocessing therefore give
\begin{equation}\label{eq:step-002-factor-privacy}
(\varepsilon_m+\varepsilon_c,\delta_m+\delta_c)
=\left(\frac\varepsilon2,\frac\delta2\right).
\end{equation}
Integrating the inequality over the common data-independent permutation
preserves it.  No consistency, realizability, successful private event, or
utility condition has been used, so the result holds on the full input
space and for every measurable output event.
\end{proof}

\begin{lemma}[Fixed-marginal core control and bad-block concentration]
\label{lem:step-002-good-cores}
Under Assumptions~\ref{assump:vc-one-factors} and
\ref{assump:countably-coded-evaluation}, fix
\(\bar c_i\in\bar C_i\) and a probability measure \(\bar D_i\) on \(Q_i\).
For iid realizable input and the symmetrized blocks of
Proposition~\ref{prop:step-002-kernel}, let
\(h_b^\circ=\mathbf1_{B_b}\), let
\(g_*=\bar c_i\oplus\bar f_i^\circ\), and call block \(b\) bad if its
transformed population risk exceeds \(1/64\).  Then
\begin{equation}\label{eq:step-002-marginal-bad-core}
\Pr[b\text{ is bad}]\le\frac1{48},
\end{equation}
the bad indicators are independent, and
\begin{equation}\label{eq:step-002-bad-core-tail}
\Pr\!\left[\#\{b:b\text{ is bad}\}\ge\frac{t_i}{12}\right]
\le e^{-t_i/128}\le\eta.
\end{equation}
\end{lemma}

\begin{proof}
A uniform permutation and restriction of an iid sequence is again iid; the
disjoint length-\(L_*\) blocks are mutually independent.  The transform
\eqref{eq:step-002-record-transform} maps them to a common distribution on
countable \(U_i\) labeled by \(g_*\).  By
Lemma~\ref{lem:step-002-core-vc}, the target and each data-dependent core
indicator lie in one fixed class of VC dimension at most one and have zero
empirical disagreement.  The VC-one generalization theorem of
\citet{yan2025vc1}, with accuracy \(1/64\), confidence \(1/48\), and block
size \eqref{eq:step-002-block-size}, gives
\eqref{eq:step-002-marginal-bad-core}.  Its event is measurable because the
quotient is countable.

Each indicator depends on only one independent block.  Applying Hoeffding's
inequality \citep{hoeffding1963} with mean sum at most \(t_i/48\) and
threshold \(t_i/12\) gives
\[
\exp\!\left(-2t_i\left(\frac1{12}-\frac1{48}\right)^2\right)
=e^{-t_i/128}.
\]
The first inequality in
\eqref{eq:step-002-calibration-conclusions} proves the final bound.  This is
a fixed marginal estimate followed by concentration, not a union bound
over blocks.
\end{proof}

\begin{proposition}[Unpadded fixed-confidence factor utility]
\label{prop:step-002-factor-utility}
Under Assumptions~\ref{assump:vc-one-factors},
\ref{assump:countably-coded-evaluation}, and
\ref{assump:global-privacy-range}, for every target
\(\bar c_i\in\bar C_i\) and probability measure \(\bar D_i\) on \(Q_i\),
if all \(q_i\) input records to \(\bar A_i^{\mathrm{Yan}}\) are iid and
labeled by \(\bar c_i\), then
\begin{equation}\label{eq:step-002-factor-utility}
\Pr\!\left[
\Pr_{q\sim\bar D_i}
[\bar A_i^{\mathrm{Yan}}(\bar T)(q)\ne\bar c_i(q)]
\le\frac1{64}\right]
\ge1-\frac1{4096}.
\end{equation}
The probability includes the sample, symmetrization, and internal randomness.
No utility claim is made for padded or nonrealizable input.
\end{proposition}

\begin{proof}
On a realizable transformed block, \(g_*\in V(T_b)\), so
\begin{equation}\label{eq:step-002-core-nesting}
B_b\subseteq P_{g_*},
\qquad h_b^\circ\le g_*\quad\text{pointwise}.
\end{equation}
Every nonempty core is a principal upset whose minimum lies on the one finite
chain \(P_{g_*}\), and depth is one-to-one and order-reversing along that
chain.

Intersect three events: the private median is a \(1/3\)-median; fewer than
\(t_i/12\) cores are bad; and, when \(z\ge1\), the choosing rule satisfies
\eqref{eq:step-002-choosing-utility}.  Each event fails with probability at
most \(\eta\), by
\eqref{eq:step-002-median-threshold},
Lemma~\ref{lem:step-002-good-cores}, and the choosing guarantee.  On the
median event, at least \(t_i/6\) block depths lie on each side of \(z\).

If \(z=0\), at least \(t_i/6\) blocks have empty core.  Since fewer than
\(t_i/12\) blocks are bad, one such core is good.  The output is transformed
zero, equal to that core indicator, and has risk at most \(1/64\).  This
also handles \(g_*=0\), for which every realizable core is empty.

Suppose \(z\ge1\).  Every block of depth at least \(z\) has a unique
layer-\(z\) point.  Since all cores lie in the target chain, these points
are one common point \(u_z^*\).  Its histogram score is at least
\(t_i/6\), and \eqref{eq:step-002-choosing-utility} together with
\eqref{eq:step-002-calibration-conclusions} leaves selected score at least
\begin{equation}\label{eq:step-002-positive-score}
\frac{t_i}{6}-\frac{t_i}{12}=\frac{t_i}{12}>0.
\end{equation}
The selection is not \(\bot\).  Positive score and uniqueness at depth
force it to equal \(u_z^*\), so the transformed output is
\begin{equation}\label{eq:step-002-output-nesting}
h_{\mathrm{out}}^\circ=\mathbf1_{I_i(u_z^*)}\le g_*.
\end{equation}

Among the at least \(t_i/6\) blocks of depth at most \(z\), one, say
\(b_0\), is good.  If its core is empty, its indicator is zero.  Otherwise
the core minimum has depth at most \(z\), so its principal upset is contained
in \(I_i(u_z^*)\).  Thus
\begin{equation}\label{eq:step-002-shallow-comparison}
h_{b_0}^\circ\le h_{\mathrm{out}}^\circ\le g_*.
\end{equation}
For binary functions under this pointwise order, disagreement with \(g_*\)
is reverse monotone, and hence
\begin{equation}\label{eq:step-002-transformed-risk}
\Pr[h_{\mathrm{out}}^\circ\ne g_*]
\le\Pr[h_{b_0}^\circ\ne g_*]\le\frac1{64}.
\end{equation}
Reference XOR cancels pointwise:
\begin{equation}\label{eq:step-002-xor-risk}
\mathbf1\{\bar f_i^\circ(q)\oplus
h_{\mathrm{out}}^\circ(\tau_i(q))\ne\bar c_i(q)\}
=\mathbf1\{h_{\mathrm{out}}^\circ(\tau_i(q))
\ne g_*(\tau_i(q))\}.
\end{equation}
The secondary quotient preserves this indicator, while
Lemma~\ref{lem:step-001-risk-pullback} further pulls it to raw factor risk
without residual.  Finally,
\begin{equation}\label{eq:step-002-factor-failure}
3\eta=\frac3{16384}<\frac1{4096},
\end{equation}
which proves \eqref{eq:step-002-factor-utility}.  The confidence \(1/48\)
was used only as a fixed block marginal; the concentration step avoids any
confidence fixed point or additional logarithm in the quota.
\end{proof}

\begin{proposition}[Measurable private factor certificate]
\label{prop:step-002-factor-certificate}
Under Assumptions~\ref{assump:vc-one-factors},
\ref{assump:countably-coded-evaluation}, and
\ref{assump:global-privacy-range}, for \(K_Y\ge K_*\), the fixed
permutation-symmetrized quotient-first rule
\(\bar A_i^{\mathrm{Yan}}\) is a Markov kernel on exactly \(q_i\) records,
is replacement-\((\varepsilon/2,\delta/2)\)-DP on every input, and satisfies
the unpadded iid utility bound \eqref{eq:step-002-factor-utility}.  Its
hidden constants depend only on the fixed numerical accuracy, confidence,
privacy split, and universal constants in the cited private-median,
VC-generalization, and choosing results, never on the class, quotient,
distribution, \(d_i\), \(\varepsilon\), or \(\delta\).
\end{proposition}

\begin{proof}
Lemma~\ref{lem:step-002-order-geometry} derives the finite order and all
principal-core properties, including the positive endpoint.  The improper
core class and its empirical consistency are
Lemma~\ref{lem:step-002-core-vc}.  The exact leaving-and-entering histogram,
countable extension, and every empty-support privacy direction are proved in
Lemma~\ref{lem:step-002-countable-choosing}.  The source thresholds,
Hoeffding confidence, choosing slack, and exact ceiling-safe bound
\(N_i\le q_i\), including \(d_i=1\), are
Lemma~\ref{lem:step-002-calibration}.  Proposition~\ref{prop:step-002-kernel}
then supplies the total measurable transition, and
Proposition~\ref{prop:step-002-factor-privacy} supplies all-input factor
privacy.  Lemma~\ref{lem:step-002-good-cores} and
Proposition~\ref{prop:step-002-factor-utility} provide the unpadded fixed-
confidence risk guarantee.  The construction remains the same one-factor
quotient-first rule when \(k=1\); no cross-factor restriction, finite-support
condition, or properness condition has been introduced.
\end{proof}

\subsection{Routed product kernel and support-two privacy}
\label{app:step-003}

\begin{lemma}[Measurable quotient routing and padded prefixes]
\label{lem:step-003-routing-measurability}
Under Assumptions~\ref{assump:canonical-product} and
\ref{assump:countably-coded-evaluation}, fix \(n\in\mathbb N\).  For every
factor \(i\), the count \(J_i:Z^n\to\{0,\ldots,n\}\), the map
\[
\bar T_i:Z^n\longrightarrow(Q_i\times\{0,1\})^{q_i}
\]
that takes the first \(q_i\) routed records and pads a shortage, and the
finite product \((\bar T_1,\ldots,\bar T_k)\) are measurable on every
labeled input, independently of realizability.
\end{lemma}

\begin{proof}
Let \(Y_i=Q_i\times\{0,1\}\) be countable discrete, adjoin
\(\dagger_i\notin Y_i\), and define
\begin{equation}\label{eq:step-003-route}
\varphi_i(x,y)=\begin{cases}
(\kappa_i(x),y),&x\in X_i,\\
\dagger_i,&x\notin X_i.
\end{cases}
\end{equation}
For any subset of the countable codomain, the preimage is a countable union
of measurable sets
\(\kappa_i^{-1}(\{q\})\times\{b\}\), together, when appropriate, with
\((X\setminus X_i)\times\{0,1\}\).  Thus \(\varphi_i\) is measurable and
\begin{equation}\label{eq:step-003-count}
J_i(S)=\sum_{j=1}^n\mathbf1\{\varphi_i(z_j)\ne\dagger_i\}
\end{equation}
is measurable.

For \(1\le r\le q_i\) and \(1\le j\le n\), let
\begin{equation}\label{eq:step-003-routed-rank-event}
F_{i,r,j}=\left\{S:\varphi_i(z_j)\ne\dagger_i,
\ \sum_{\ell<j}\mathbf1\{\varphi_i(z_\ell)\ne\dagger_i\}=r-1\right\}.
\end{equation}
This is the event that position \(j\) is the \(r\)-th routed record.  The
\(r\)-th padded-prefix coordinate equals that routed record on the unique
such event and equals \(\bar z_i^\circ\) when \(J_i<r\).  Hence, for
\(u\in Y_i\),
\begin{equation}\label{eq:step-003-prefix-fiber}
\{\bar T_{i,r}=u\}
=\bigl(\{J_i<r\}\cap\{\bar z_i^\circ=u\}\bigr)
\cup\bigcup_{j=1}^n
\bigl(F_{i,r,j}\cap\{\varphi_i(z_j)=u\}\bigr),
\end{equation}
where the first set is empty if \(u\ne\bar z_i^\circ\).  This fiber is
measurable.  Countable discreteness gives coordinate and tuple
measurability, and finite products preserve it.  No target, selector, or
consistency event is used.
\end{proof}

\begin{lemma}[One-edit stability of a padded prefix multiset]
\label{lem:step-003-padded-prefix}
Let \(Y\) be a set, \(q\ge1\), and \(p\in Y\).  For a finite list
\(w=(w_1,\ldots,w_J)\), define
\[
\operatorname{Pad}_{q,p}(w)=
(w_1,\ldots,w_{\min\{J,q\}},
\underbrace{p,\ldots,p}_{q-\min\{J,q\}\text{ copies}}).
\]
If \(w'\) is obtained by one replacement, deletion, or insertion, the two
padded-prefix multisets are equal or differ by one multiset replacement.
Consequently, the two tuples can be permuted to become equal or to differ
in exactly one coordinate.
\end{lemma}

\begin{proof}
The complete calculation is as follows:
\[
\begin{array}{c|c|c}
\text{edit}&\text{position/length}&\text{sole multiset change}\\ \hline
w_r\mapsto v&r>q&\text{none}\\
w_r\mapsto v&r\le q&w_r\mapsto v\\
\text{delete }w_r&r>q&\text{none}\\
\text{delete }w_r&r\le q,\ J\ge q+1&w_r\mapsto w_{q+1}\\
\text{delete }w_r&r\le q,\ J\le q&w_r\mapsto p\\
\text{insert }v\text{ at }r&r>q&\text{none}\\
\text{insert }v\text{ at }r\le q,\ J\ge q&w_q\mapsto v\\
\text{insert }v\text{ at }r\le q,\ J<q&p\mapsto v
\end{array}
\]
For instance, after deleting \(w_r\) with \(r\le q<J\), shifted entries
\(w_{r+1},\ldots,w_q\) occur in both multisets; only \(w_r\) leaves and
\(w_{q+1}\) enters.  If the list ends within the prefix, padding enters.
Insertion is the reverse calculation.  Coinciding entering and leaving
values can only turn the nominal replacement into equality.

The multiset intersection of two size-\(q\) multisets differing by at most
one replacement has multiplicity at least \(q-1\).  Align these common
occurrences in the first \(q-1\) positions and put the possible unmatched
occurrences last.  The resulting permutations give equality or ordered
one-coordinate replacement, including \(q=1\).
\end{proof}

\begin{proposition}[Support-two adjacency of routed inputs]
\label{prop:step-003-support-two}
Under Assumption~\ref{assump:canonical-product},
Lemma~\ref{lem:step-003-padded-prefix}, and
Proposition~\ref{prop:step-002-kernel}, if \(S,S'\in Z^n\) differ by one
replacement, there is \(I(S,S')\subseteq\{1,\ldots,k\}\),
\(|I(S,S')|\le2\), such that unaffected factor-input multisets agree; for
each affected factor the two inputs can be permuted into one-coordinate
replacement neighbors; and the corresponding factor output laws are
unchanged by those permutations.  This includes arbitrary labels, changed
rows outside the selected prefixes, same-block and cross-block moves, and
every full/short padding boundary.
\end{proposition}

\begin{proof}
Let the changed row be \(z=(x,y)\) in \(S\) and \(z'=(x',y')\) in \(S'\),
and let \(x\in X_a\), \(x'\in X_b\).  Write \(R_i(S)\) for the ordered
list of quotient records routed to factor \(i\); then
\(\bar T_i(S)=\operatorname{Pad}_{q_i,\bar z_i^\circ}(R_i(S))\).

If \(a=b\), only \(R_a\) changes, by one replacement at the same local
rank.  Lemma~\ref{lem:step-003-padded-prefix} makes the selected multiset
equal if the rank exceeds \(q_a\), and otherwise one-replacement different.
Thus \(I\subseteq\{a\}\).

If \(a\ne b\), every other routed list is unchanged.  The leaving list
deletes one record and the entering list inserts one record at the rank
determined by the unchanged global position.  The lemma gives at most one
multiset replacement in each: a deletion from a full prefix brings in the
next real record, or padding if none remains; deletion from a short prefix
brings in padding; insertion into a full prefix replaces the former last
selected record; insertion into a short prefix replaces padding; and an
edit after the selected real prefix changes nothing.  Hence
\(I\subseteq\{a,b\}\).

The same lemma supplies permutations \(\sigma_i,\sigma_i'\) aligning all
common multiplicities.  Permutation invariance from
Proposition~\ref{prop:step-002-kernel} yields, for every measurable
\(E_i\subseteq\mathcal H_i\),
\begin{equation}\label{eq:step-003-factor-permutation}
\begin{split}
\bar A_i^{\mathrm{Yan}}(\bar T_i(S),E_i)
&=\bar A_i^{\mathrm{Yan}}(\sigma_i\bar T_i(S),E_i),\\
\bar A_i^{\mathrm{Yan}}(\bar T_i(S'),E_i)
&=\bar A_i^{\mathrm{Yan}}(\sigma_i'\bar T_i(S'),E_i).
\end{split}
\end{equation}
The permutations are only a coupling for this fixed adjacent pair; the
learner need not choose them measurably.  Padding and arbitrary labels lie
in the total factor input space.  At \(k=1\), necessarily \(a=b=1\), so at
most the sole factor changes.
\end{proof}

\begin{proposition}[Measurable routed product kernel and decoder]
\label{prop:step-003-product-kernel}
Under Assumptions~\ref{assump:canonical-product} and
\ref{assump:countably-coded-evaluation}, Lemmas~\ref{lem:step-001-output-measurability}
and \ref{lem:step-001-risk-pullback},
Proposition~\ref{prop:step-002-kernel}, and
Lemma~\ref{lem:step-003-routing-measurability},
\begin{equation}\label{eq:step-003-product-kernel}
\bar A_n^{\oplus,Q}(S,d\bar h)
=\bigotimes_{i=1}^k
\bar A_i^{\mathrm{Yan}}(\bar T_i(S),d\bar h_i)
\end{equation}
defines a Markov kernel into
\((\mathcal H^{\oplus},\mathscr H^{\oplus})\), with independent factor
randomness.  Piecewise decoding is measurable under the finite-evaluation-
cylinder convention, so its pushforward is a decoded Markov kernel.  For
every fixed target and distribution, every displayed risk event is
measurable in the tuple output.
\end{proposition}

\begin{proof}
For \(E_i\in\mathscr H_i\),
\begin{equation}\label{eq:step-003-routed-coordinate-kernel}
K_i(S,E_i)=\bar A_i^{\mathrm{Yan}}(\bar T_i(S),E_i)
\end{equation}
is measurable in \(S\) by Lemma~\ref{lem:step-003-routing-measurability}
and the factor-kernel property.  For fixed \(S\), the finite product of
these probability measures exists, and for every measurable rectangle,
\begin{equation}\label{eq:step-003-product-rectangle}
\bar A_n^{\oplus,Q}\!\left(S,\prod_{i=1}^kE_i\right)
=\prod_{i=1}^kK_i(S,E_i),
\end{equation}
which is measurable in \(S\).

Let \(\mathcal D\) be the events in \(\mathscr H^{\oplus}\) whose
probability under \eqref{eq:step-003-product-kernel} is measurable in
\(S\).  It contains the whole space, is closed under complements because
the complementary probability is one minus the original probability, and
is closed under countable disjoint unions by pointwise convergence of
partial sums.  It is therefore a Dynkin system.  Measurable rectangles form
a pi-system contained in \(\mathcal D\) by
\eqref{eq:step-003-product-rectangle} and generate
\(\mathscr H^{\oplus}\).  The pi-lambda theorem gives
\(\mathcal D=\mathscr H^{\oplus}\), establishing kernel measurability for
arbitrary joint events.

Let \(\mathcal H_\Sigma\) denote the \(\Sigma\)-measurable binary
hypotheses and let \(\mathscr G_{\mathrm{ev}}\) be generated by their finite
evaluation cylinders.  Lemma~\ref{lem:step-001-output-measurability} makes
\[
\mathrm{Dec}:\mathcal H^{\oplus}\to\mathcal H_\Sigma,
\qquad \mathrm{Dec}(\bar h)=h_{\bar h},
\]
measurable.  Thus
\begin{equation}\label{eq:step-003-decoded-kernel}
\bar A_n^{\oplus,\mathrm{dec}}(S,G)
=\bar A_n^{\oplus,Q}(S,\mathrm{Dec}^{-1}(G))
\end{equation}
is a kernel for \(G\in\mathscr G_{\mathrm{ev}}\).  Finally,
Lemma~\ref{lem:step-001-risk-pullback} gives
\begin{equation}\label{eq:step-003-risk-event}
\{\bar h:R_D(h_{\bar h},c)\le t\}\in\mathscr H^{\oplus}
\end{equation}
for fixed \(c,D,t\).  This is the required output convention without a
standard-Borel assumption on the raw instance space.
\end{proof}

\begin{proposition}[Support-two composition for joint output events]
\label{prop:step-003-joint-composition}
Under Assumption~\ref{assump:global-privacy-range},
Proposition~\ref{prop:step-002-factor-privacy}, and
Propositions~\ref{prop:step-003-support-two} and
\ref{prop:step-003-product-kernel}, for every \(n\), adjacent
\(S,S'\in Z^n\), and \(E\in\mathscr H^{\oplus}\),
\begin{equation}\label{eq:step-003-global-dp}
\bar A_n^{\oplus,Q}(S,E)
\le e^\varepsilon\bar A_n^{\oplus,Q}(S',E)+\delta.
\end{equation}
The same inequality holds for every measurable deterministic
postprocessing, including decoding.  At most two factor privacy costs are
used, with no further dependence on \(k\).
\end{proposition}

\begin{proof}
Fix \(S,S'\), let \(I=I(S,S')\), and put \(m=|I|\le2\).  Define
\[
\mu_i=\bar A_i^{\mathrm{Yan}}(\bar T_i(S),\cdot),
\qquad
\nu_i=\bar A_i^{\mathrm{Yan}}(\bar T_i(S'),\cdot).
\]
For \(i\notin I\), permutation invariance gives \(\mu_i=\nu_i\).  For
\(i\in I\), \eqref{eq:step-003-factor-permutation} and factor privacy give,
for every \(B\in\mathscr H_i\),
\begin{equation}\label{eq:step-003-factor-dp-pair}
\mu_i(B)\le a\nu_i(B)+d,\qquad
\nu_i(B)\le a\mu_i(B)+d,\qquad
a=e^{\varepsilon/2},\quad d=\delta/2.
\end{equation}

We prove composition on the full product sigma-field.  For \(i\in I\), let
\(\lambda_i=\mu_i+\nu_i\), with densities
\(u_i=d\mu_i/d\lambda_i\) and \(v_i=d\nu_i/d\lambda_i\).  The elementary
hockey-stick identity
\[
\sup_B(\mu_i(B)-a\nu_i(B))
=\int(u_i-av_i)_+\,d\lambda_i
\]
follows by taking the positive-density set, and
\eqref{eq:step-003-factor-dp-pair} in both directions gives
\begin{equation}\label{eq:step-003-hockey-stick}
\int(u_i-av_i)_+\,d\lambda_i\le d,\qquad
\int(v_i-au_i)_+\,d\lambda_i\le d.
\end{equation}
Define subprobability densities
\begin{equation}\label{eq:step-003-truncated-densities}
u_i^0=\min\{u_i,av_i\},\qquad
v_i^0=\min\{v_i,au_i\}.
\end{equation}
Their measures \(\mu_i^0,\nu_i^0\) satisfy
\begin{equation}\label{eq:step-003-truncated-mass}
\mu_i^0(\mathcal H_i)\ge1-d,\qquad
\nu_i^0(\mathcal H_i)\ge1-d,
\end{equation}
and pointwise
\begin{equation}\label{eq:step-003-mutual-domination}
u_i^0\le av_i^0,\qquad v_i^0\le au_i^0.
\end{equation}
For example, if \(v_i^0=v_i\), the first inequality follows from
\(u_i^0\le av_i\); if \(v_i^0=au_i\), then
\(av_i^0=a^2u_i\ge u_i\ge u_i^0\).  The other direction is symmetric.

For unaffected indices set \(\mu_i^0=\nu_i^0=\mu_i\), and form
\begin{equation}\label{eq:step-003-product-submeasures}
\mu=\bigotimes_i\mu_i,\quad \nu=\bigotimes_i\nu_i,\quad
\mu^0=\bigotimes_i\mu_i^0,\quad \nu^0=\bigotimes_i\nu_i^0.
\end{equation}
Product densities under the product dominating measure show
\(\mu^0\le\mu\), \(\nu^0\le\nu\), and
\(\mu^0\le a^m\nu^0\).  Moreover,
\begin{equation}\label{eq:step-003-product-retained-mass}
\mu^0(\mathcal H^{\oplus})
=\prod_{i\in I}\mu_i^0(\mathcal H_i)
\ge(1-d)^m\ge1-md.
\end{equation}
Thus every joint event \(E\) satisfies
\begin{equation}\label{eq:step-003-composition-chain}
\begin{split}
\mu(E)
&\le\mu^0(E)+1-\mu^0(\mathcal H^{\oplus})\\
&\le a^m\nu^0(E)+md\\
&\le e^{m\varepsilon/2}\nu(E)+m\delta/2\\
&\le e^\varepsilon\nu(E)+\delta.
\end{split}
\end{equation}
These are exactly the output laws in
\eqref{eq:step-003-product-kernel}, proving
\eqref{eq:step-003-global-dp}.  For \(m=0\) there is equality; for \(m=1\)
the stronger factor parameters remain; and for \(m=2\) each factor budget
is used once.  No unchanged factor is composed and no rectangle restriction
appears.  Reversing \(S,S'\) proves the other ordered direction.  For a
measurable map \(g\), apply the inequality to \(g^{-1}(G)\); taking
\(g=\mathrm{Dec}\) proves decoded privacy.  Since factor privacy held on all
inputs, the proof includes inconsistent labels and padding.
\end{proof}

\begin{proposition}[Measurable routed privacy certificate]
\label{prop:step-003-global-kernel-certificate}
Under Assumptions~\ref{assump:canonical-product},
\ref{assump:countably-coded-evaluation}, and
\ref{assump:global-privacy-range}, the routed quotient-first rule
\(\bar A_n^{\oplus,Q}\) is a measurable Markov kernel for every fixed
sample size, its decoded risk events are measurable, and it is replacement-
\((\varepsilon,\delta)\)-DP on every labeled input pair.  A replacement
affects at most two factor laws.  When \(k=1\), the rule is the sole routed
factor kernel and in fact retains the stronger
\((\varepsilon/2,\delta/2)\) privacy parameters.
\end{proposition}

\begin{proof}
Lemma~\ref{lem:step-003-routing-measurability} proves measurable raw-to-
quotient preprocessing.  Lemma~\ref{lem:step-003-padded-prefix} accounts for
every replacement, deletion, insertion, and padding transition.
Proposition~\ref{prop:step-003-support-two} applies it to a global record
replacement and invokes permutation invariance only after aligning the
multisets.  Proposition~\ref{prop:step-003-product-kernel} gives the full
joint kernel, decoder, and risk-event convention.
Proposition~\ref{prop:step-003-joint-composition} then proves arbitrary-
event privacy with exact cost at most
\((2\varepsilon/2,2\delta/2)=(\varepsilon,\delta)\).  Its \(m=1\) branch
is the claimed one-factor specialization.
\end{proof}

\subsection{Weighted occupancy shortages}\label{app:step-004}

\begin{lemma}[Marginal occupancy and weighted expectation identity]
\label{lem:step-004-binomial-accounting}
Under Assumption~\ref{assump:canonical-product} and
Lemmas~\ref{lem:step-001-quotient-invariance} and
\ref{lem:step-001-output-measurability}, fix \(n\in\mathbb N\),
\(c\in C\), a probability measure \(D\) on \((X,\Sigma)\), and an iid
labeled sample
\(S=((x_j,c(x_j)))_{j=1}^n\sim D_c^n\).  With
\[
\rho_i=D(X_i),\qquad
J_i=\bigl|\{j:x_j\in X_i\}\bigr|,
\qquad
W_{\mathrm{short}}
=\sum_{i=1}^k\rho_i\mathbf 1\{J_i<q_i\},
\]
one has
\begin{equation}\label{eq:step-004-binomial-law}
J_i\sim\operatorname{Bin}(n,\rho_i)
\quad\text{for every }i.
\end{equation}
If \(\rho_i=0\), then
\(\rho_i\mathbf 1\{J_i<q_i\}=0\) for every sample realization and
\(J_i=0\) almost surely.  Moreover, \(W_{\mathrm{short}}\) is integrable
and
\begin{equation}\label{eq:step-004-expectation-identity}
\mathbb E W_{\mathrm{short}}
=\sum_{i=1}^k\rho_i\Pr(J_i<q_i).
\end{equation}
Neither conclusion requires independence among the counts for distinct
factors.
\end{lemma}

\begin{proof}
For each factor define \(\bar b^{(i)}\in\mathcal H^\oplus\) by
\[
\bar b^{(i)}_\ell(q)=\mathbf 1\{\ell=i\},
\qquad 1\le \ell\le k,\quad q\in Q_\ell.
\]
Lemma~\ref{lem:step-001-output-measurability} makes the decoded hypothesis
measurable, and the whole-domain partition in
Assumption~\ref{assump:canonical-product} gives
\(h_{\bar b^{(i)}}=\mathbf 1_{X_i}\).  Thus \(X_i\in\Sigma\).
Lemma~\ref{lem:step-001-quotient-invariance} gives quotient
representatives \(\bar c_i\) of the restrictions of \(c\); their tuple
decodes exactly to \(c\), so
Lemma~\ref{lem:step-001-output-measurability} also makes \(c\) measurable.
Consequently \(D_c\), \(D_c^n\), the masses \(\rho_i\), and the block
indicators below are well defined.

For fixed \(i\), put
\[
I_{j,i}=\mathbf 1\{x_j\in X_i\},\qquad 1\le j\le n.
\]
The iid sampling assumption makes \(I_{1,i},\ldots,I_{n,i}\) independent
Bernoulli variables with success probability \(D(X_i)=\rho_i\).  Since
\(J_i=\sum_{j=1}^n I_{j,i}\), this proves
\eqref{eq:step-004-binomial-law}.  Counts for distinct factors can be
dependent because each observation belongs to exactly one block, but no
subsequent argument uses their joint independence.

When \(\rho_i=0\), multiplication by the deterministic weight gives
\[
\rho_i\mathbf 1\{J_i<q_i\}=0
\]
for every realization.  Also \(I_{j,i}=0\) almost surely for each of the
finitely many sample indices, hence \(J_i=0\) almost surely.  This handles
zero mass without forming a conditional law on \(X_i\).

The finite measurable partition of all of \(X\) gives
\[
0\le W_{\mathrm{short}}
\le\sum_{i=1}^k\rho_i=D(X)=1.
\]
Thus all summands and their sum are integrable.  Finite linearity of
expectation now yields
\[
\mathbb E W_{\mathrm{short}}
=\sum_{i=1}^k\rho_i\mathbb E\mathbf 1\{J_i<q_i\}
=\sum_{i=1}^k\rho_i\Pr(J_i<q_i),
\]
which is \eqref{eq:step-004-expectation-identity}.  This is a finite-sum
identity, not a limit interchange or a union bound.
\end{proof}

\begin{lemma}[Explicit heavy-factor shortage tail]
\label{lem:step-004-heavy-shortage}
Let \(J\sim\operatorname{Bin}(n,p)\), put \(\mu=np\), and let
\(q\in\mathbb N\) satisfy \(q\ge1\).  If \(\mu\ge128q\), then
\begin{equation}\label{eq:step-004-heavy-tail}
\Pr(J<q)
\le \exp\!\left(-\frac{16129}{256}q\right)
\le e^{-16}.
\end{equation}
\end{lemma}

\begin{proof}
For any \(u\in(0,1)\), exponential Markov gives
\begin{equation}\label{eq:step-004-exp-markov}
\begin{split}
\Pr\bigl(J\le(1-u)\mu\bigr)
&=\Pr\bigl(e^{-uJ}\ge e^{-u(1-u)\mu}\bigr)\\
&\le e^{u(1-u)\mu}\,\mathbb E e^{-uJ}.
\end{split}
\end{equation}
The binomial transform and \(1+v\le e^v\) give
\begin{equation}\label{eq:step-004-binomial-transform}
\mathbb E e^{-uJ}
=(1-p+pe^{-u})^n
\le\exp\bigl(\mu(e^{-u}-1)\bigr).
\end{equation}
For \(u\ge0\),
\begin{equation}\label{eq:step-004-exp-quadratic}
e^{-u}\le1-u+\frac{u^2}{2}.
\end{equation}
Indeed, for
\(f(u)=1-u+u^2/2-e^{-u}\), one has
\(f(0)=f'(0)=0\) and \(f''(u)=1-e^{-u}\ge0\).
Substitution of \eqref{eq:step-004-binomial-transform} and
\eqref{eq:step-004-exp-quadratic} into
\eqref{eq:step-004-exp-markov} yields
\begin{equation}\label{eq:step-004-chernoff-derived}
\Pr\bigl(J\le(1-u)\mu\bigr)
\le
\exp\!\left(\mu\left[u(1-u)-u+\frac{u^2}{2}\right]\right)
=\exp\!\left(-\frac{\mu u^2}{2}\right).
\end{equation}

Under the stated hypothesis, \(q\le\mu/128\), and hence
\begin{equation}\label{eq:step-004-event-inclusion}
\{J<q\}\subseteq\left\{J\le\frac{\mu}{128}\right\}.
\end{equation}
Apply \eqref{eq:step-004-chernoff-derived} with \(u=127/128\).  Then
\begin{equation}\label{eq:step-004-explicit-exponent}
\frac{\mu}{2}\left(\frac{127}{128}\right)^2
\ge
\frac{128q}{2}\frac{127^2}{128^2}
=\frac{16129}{256}q.
\end{equation}
Since \(q\ge1\),
\begin{equation}\label{eq:step-004-exponent-sixteen}
\frac{16129}{256}q
\ge\frac{16129}{256}
>\frac{4096}{256}=16.
\end{equation}
Equations~\eqref{eq:step-004-chernoff-derived}--
\eqref{eq:step-004-exponent-sixteen} prove
\eqref{eq:step-004-heavy-tail}.
\end{proof}

\begin{proposition}[Weighted shortage certificate without balance]
\label{prop:step-004-weighted-shortage}
Under Assumption~\ref{assump:canonical-product},
Lemmas~\ref{lem:step-001-quotient-invariance},
\ref{lem:step-001-output-measurability},
\ref{lem:step-004-binomial-accounting}, and
\ref{lem:step-004-heavy-shortage}, fix \(n\in\mathbb N\), \(c\in C\),
a probability measure \(D\), and \(S\sim D_c^n\).  For the setting-defined
masses, counts, quotas, and shortage mass, suppose \(q_i\ge1\) and put
\(Q_\oplus=\sum_{i=1}^kq_i\).  Then, for every mass vector induced by
\(D\),
\begin{equation}\label{eq:step-004-weighted-shortage}
\mathbb E W_{\mathrm{short}}
\le\frac{128Q_\oplus}{n}+e^{-16}.
\end{equation}
No balance or positive-mass condition, factor-count independence, or
union bound over factors is required.  If \(k=1\) and \(n\ge q_1\), then
\begin{equation}\label{eq:step-004-one-factor-shortage}
W_{\mathrm{short}}=0
\quad\text{for every sample realization}.
\end{equation}
\end{proposition}

\begin{proof}
Put \(\mu_i=n\rho_i\) and split the finite factor set as
\begin{equation}\label{eq:step-004-light-heavy}
\mathcal L=\{i:\mu_i<128q_i\},
\qquad
\mathcal H=\{i:\mu_i\ge128q_i\}.
\end{equation}
For a light factor, \(\rho_i<128q_i/n\), and a shortage probability is at
most one.  Therefore the universally valid non-strict bound, including
the empty-light case, is
\begin{equation}\label{eq:step-004-light-bound}
\sum_{i\in\mathcal L}\rho_i\Pr(J_i<q_i)
\le\sum_{i\in\mathcal L}\rho_i
\le\frac{128}{n}\sum_{i\in\mathcal L}q_i
\le\frac{128Q_\oplus}{n}.
\end{equation}
A zero-mass factor is light because \(q_i\ge1\), and its summand is
identically zero by Lemma~\ref{lem:step-004-binomial-accounting}; no
conditional factor law is needed.

For \(i\in\mathcal H\), Lemma~\ref{lem:step-004-binomial-accounting}
gives \(J_i\sim\operatorname{Bin}(n,\rho_i)\), while
\(n\rho_i\ge128q_i\).  Lemma~\ref{lem:step-004-heavy-shortage} thus gives
\(\Pr(J_i<q_i)\le e^{-16}\).  Because the measurable blocks partition
all of \(X\), \(\sum_i\rho_i=1\), so
\begin{equation}\label{eq:step-004-heavy-bound}
\sum_{i\in\mathcal H}\rho_i\Pr(J_i<q_i)
\le e^{-16}\sum_{i\in\mathcal H}\rho_i
\le e^{-16}.
\end{equation}
This is a weighted average of marginal shortage probabilities, not a
bound on an all-factor union event.

Splitting the exact identity
\eqref{eq:step-004-expectation-identity} according to
\eqref{eq:step-004-light-heavy} and using
\eqref{eq:step-004-light-bound}--\eqref{eq:step-004-heavy-bound} proves
\eqref{eq:step-004-weighted-shortage}.

Finally, if \(k=1\), Assumption~\ref{assump:canonical-product} gives
\(X_1=X\), so \(\rho_1=1\) and every record is routed to factor one.
Hence \(J_1=n\) pointwise.  When \(n\ge q_1\), the shortage event is
empty and
\[
W_{\mathrm{short}}
=\rho_1\mathbf 1\{J_1<q_1\}=0,
\]
proving \eqref{eq:step-004-one-factor-shortage} exactly rather than through
the generic residual bound.
\end{proof}

\subsection{From factor utility to global PAC utility}
\label{app:step-005}

\begin{lemma}[Measurability of routed bad-factor mass]
\label{lem:step-005-measurable-bad-mass}
Under Assumptions~\ref{assump:canonical-product} and
\ref{assump:countably-coded-evaluation} and
Proposition~\ref{prop:step-002-kernel}, fix \(n\in\mathbb N\), \(c\in C\),
and a probability measure \(D\) on \((X,\Sigma)\).  In the joint experiment
consisting of \(S\sim D_c^n\) and the routed factor outputs, every
\(J_i\), routed input \(\bar T_i(S)\), factor risk, and event
\(\mathsf B_i\) is measurable.  Consequently \(W_{\rm bad}\) is
measurable and integrable, with
\begin{equation}\label{eq:step-005-bad-mass-range}
0\le W_{\rm bad}\le1.
\end{equation}
\end{lemma}

\begin{proof}
Assumption~\ref{assump:countably-coded-evaluation} gives \(X_i\in\Sigma\)
and makes \(\kappa_i:X_i\to Q_i\) measurable into a finite or countable
discrete quotient.  Hence
\begin{equation}\label{eq:step-005-count-measurable}
J_i(S)=\sum_{j=1}^n\mathbf1\{x_j\in X_i\}
\end{equation}
is measurable on the labeled sample space \(Z^n\).

For \(1\le r\le q_i\) and \(1\le j\le n\), let
\begin{equation}\label{eq:step-005-position-events}
E_{i,r,j}=
\left\{x_j\in X_i,\ 
\sum_{\ell<j}\mathbf1\{x_\ell\in X_i\}=r-1\right\}.
\end{equation}
For fixed \(i,r\), these measurable events are disjoint in \(j\).  On
\(E_{i,r,j}\), the \(r\)-th routed record is
\((\kappa_i(x_j),y_j)\), while on the measurable complement
\(\{J_i<r\}\) it is the fixed padding record \(\bar z_i^\circ\).
Finite measurable pasting makes every routed coordinate measurable into
the countable discrete space \(Q_i\times\{0,1\}\).  Thus
\begin{equation}\label{eq:step-005-prefix-map}
\bar T_i:Z^n\longrightarrow(Q_i\times\{0,1\})^{q_i}
\end{equation}
is measurable as an ordered prefix, not merely as a multiset.

Write
\begin{equation}\label{eq:step-005-factor-kernel}
K_i(t,E)=\bar A_i^{\rm Yan}(t,E),\qquad E\in\mathscr H_i.
\end{equation}
Proposition~\ref{prop:step-002-kernel} and
\eqref{eq:step-005-prefix-map} make
\(S\mapsto K_i(\bar T_i(S),E)\) measurable.  The prescribed independent
factor randomizations define the finite product kernel
\begin{equation}\label{eq:step-005-product-kernel}
K^\oplus(S,d\bar h)
=\bigotimes_{i=1}^kK_i(\bar T_i(S),d\bar h_i)
\end{equation}
on \((\mathcal H^\oplus,\mathscr H^\oplus)\).  For every measurable
rectangle \(E_1\times\cdots\times E_k\), its transition probability is
\(\prod_iK_i(\bar T_i(S),E_i)\), a measurable function of \(S\).
Measurable rectangles form a generating pi-system.  The events whose
transition probabilities are measurable form a Dynkin system: they
contain the whole space and are closed under complements and countable
disjoint unions.  The pi-lambda theorem therefore extends measurability to
all of \(\mathscr H^\oplus\).  For fixed \(S\),
\eqref{eq:step-005-product-kernel} is the ordinary finite product measure.
Consequently
\begin{equation}\label{eq:step-005-joint-law}
\mathbb P_{c,D,n}(dS,d\bar h)
=D_c^n(dS)K^\oplus(S,d\bar h)
\end{equation}
is a well-defined joint law.  If \(\bar H_i\) is the \(i\)-th output
coordinate, then
\begin{equation}\label{eq:step-005-conditional-marginal}
\mathbb P_{c,D,n}(\bar H_i\in E\mid S)
=K_i(\bar T_i(S),E).
\end{equation}
Only this marginal identity, not independence of failure events, is used
below.

Fix \(i\) with \(\rho_i>0\).  Enumerate the finite or countable \(Q_i\)
as \(q_{i,1},q_{i,2},\ldots\), stopping in the finite case, and define
\begin{equation}\label{eq:step-005-factor-risk-map}
\begin{split}
r_i(\bar h_i)
&=\Pr_{q\sim\bar D_i}
  [\bar h_i(q)\ne\bar c_i(q)]\\
&=\sum_m\bar D_i(\{q_{i,m}\})
  \mathbf1\{\bar h_i(q_{i,m})\ne\bar c_i(q_{i,m})\}.
\end{split}
\end{equation}
Its finite partial sums are \(\mathscr H_i\)-measurable because coordinate
evaluations generate \(\mathscr H_i\), and the series is their
nondecreasing pointwise limit.  Thus \(r_i\) and the strict event
\(\{r_i(\bar H_i)>1/64\}\) are measurable.  This uses countability only
in the quotient, not in the raw support of \(D_i\).

On the joint space \eqref{eq:step-005-joint-law}, for \(\rho_i>0\),
\begin{equation}\label{eq:step-005-bad-event}
\mathsf B_i
=\{J_i<q_i\}
 \cup
 \bigl(\{J_i\ge q_i\}
       \cap\{r_i(\bar H_i)>1/64\}\bigr),
\end{equation}
and \(\mathsf B_i=\varnothing\) when \(\rho_i=0\).  Every bad event is
therefore measurable.  The finite measurable partition gives
\(\sum_i\rho_i=1\), so
\begin{equation}\label{eq:step-005-bad-mass-definition}
W_{\rm bad}=\sum_{i=1}^k\rho_i\mathbf1_{\mathsf B_i}
\end{equation}
is measurable and satisfies \eqref{eq:step-005-bad-mass-range}, proving
integrability over both sample and learner randomness.
\end{proof}

\begin{lemma}[Conditional iid law of every unpadded first prefix]
\label{lem:step-005-prefix-iid}
Under Assumptions~\ref{assump:canonical-product} and
\ref{assump:countably-coded-evaluation},
Proposition~\ref{prop:step-002-factor-utility}, and
Lemma~\ref{lem:step-005-measurable-bad-mass}, fix \(n,c,D\) and the joint
experiment above.  For every \(i\) with \(\rho_i>0\), conditional on any
positive-probability block-position vector containing at least \(q_i\)
occurrences of \(i\), the unpadded input \(\bar T_i\) has law
\(((\bar D_i)_{\bar c_i})^{q_i}\), and
\begin{equation}\label{eq:step-005-conditional-failure}
\Pr\!\left[r_i(\bar H_i)>\frac1{64}
  \ \middle|\ \text{that block-position vector}\right]
\le\frac1{4096}.
\end{equation}
If the vector has fewer than \(q_i\) occurrences, the prefix is padded and
no factor-utility assertion is made.  If \(\rho_i=0\), neither \(D_i\)
nor \(\bar D_i\) is formed and \(\mathsf B_i=\varnothing\).
\end{lemma}

\begin{proof}
Let \(G_j\in\{1,\ldots,k\}\) be the unique block index with
\(x_j\in X_{G_j}\), and put \(G=(G_1,\ldots,G_n)\).  This is a measurable
finite-valued random vector.  For
\(a=(a_1,\ldots,a_n)\in\{1,\ldots,k\}^n\), iid sampling gives
\begin{equation}\label{eq:step-005-position-probability}
\Pr(G=a)=\prod_{j=1}^n\rho_{a_j}.
\end{equation}
Suppose this probability is positive.  Then all \(\rho_{a_j}>0\).  For
measurable \(E_j\subseteq X_{a_j}\),
\begin{equation}\label{eq:step-005-conditional-product}
\begin{split}
&\Pr(x_1\in E_1,\ldots,x_n\in E_n\mid G=a)\\
&\qquad=
\frac{\prod_{j=1}^nD(E_j)}{\prod_{j=1}^n\rho_{a_j}}
=\prod_{j=1}^nD_{a_j}(E_j).
\end{split}
\end{equation}
For general measurable coordinate sets, intersect them with their
prescribed blocks.  Equality on measurable rectangles identifies the full
conditional product law.  This conditions on one atom of a finite
partition and needs no regular conditional distribution on the raw space.

Fix \(i\) with \(\rho_i>0\), and define the labeled quotient law
\begin{equation}\label{eq:step-005-labeled-quotient-law}
\nu_i
=\bigl(q\mapsto(q,\bar c_i(q))\bigr)_\#\bar D_i
=:(\bar D_i)_{\bar c_i}
\quad\text{on }Q_i\times\{0,1\}.
\end{equation}
On a positive-probability atom \(\{G=a\}\),
\eqref{eq:step-005-conditional-product}, deterministic target labels, and
measurable quotient pushforward show that the records at positions with
\(a_j=i\) are independent with common law \(\nu_i\).  When there are at
least \(q_i\) such positions, their first \(q_i\) positions are
deterministic functions of \(a\), and hence
\begin{equation}\label{eq:step-005-prefix-law}
\mathcal L(\bar T_i\mid G=a)=\nu_i^{q_i}.
\end{equation}
Conditioning on positions has not conditioned on the within-block data or
on learner randomness.

By \eqref{eq:step-005-conditional-marginal}, after \(\bar T_i=t\) is
fixed the factor output has law \(K_i(t,\cdot)\).  Therefore
\eqref{eq:step-005-prefix-law} gives
\begin{equation}\label{eq:step-005-factor-utility-integral}
\begin{split}
&\Pr\!\left(r_i(\bar H_i)>\frac1{64}\mid G=a\right)\\
&\quad=
\int K_i\!\left(
t,\left\{\bar h_i:r_i(\bar h_i)>\frac1{64}\right\}
\right)\nu_i^{q_i}(dt)
\le\frac1{4096},
\end{split}
\end{equation}
where the last inequality is Proposition~\ref{prop:step-002-factor-utility}.
It includes the iid factor prefix and all internal factor randomness; no
other factor is conditioned to be good.

If \(a\) has fewer than \(q_i\) occurrences of \(i\), at least one
padding record is inserted, so the factor-utility premise is unavailable.
This entire branch is contained in
\(\{J_i<q_i\}\subseteq\mathsf B_i\).  If \(\rho_i=0\), every atom
containing \(i\) has probability zero by
\eqref{eq:step-005-position-probability}; the definition instead sets
\(\mathsf B_i=\varnothing\), with zero weight in every risk sum.
\end{proof}

\begin{proposition}[Weighted bad-mass expectation without factor unions]
\label{prop:step-005-weighted-bad-mass}
Under Assumptions~\ref{assump:canonical-product},
\ref{assump:countably-coded-evaluation}, and
\ref{assump:global-privacy-range},
Propositions~\ref{prop:step-002-factor-utility} and
\ref{prop:step-004-weighted-shortage}, and
Lemmas~\ref{lem:step-005-measurable-bad-mass} and
\ref{lem:step-005-prefix-iid}, for every \(n\in\mathbb N\), \(c\in C\),
and probability measure \(D\),
\begin{equation}\label{eq:step-005-bad-mass-expectation}
\mathbb E W_{\rm bad}
\le\frac{128Q_\oplus}{n}+e^{-16}+\frac1{4096}.
\end{equation}
The expectation is over the sample and all learner randomness.  No
independence among output or bad events and no factor union bound is used.
\end{proposition}

\begin{proof}
For \(\rho_i>0\), let
\begin{equation}\label{eq:step-005-no-shortage-failure}
F_i=\left\{J_i\ge q_i,\ r_i(\bar H_i)>\frac1{64}\right\}.
\end{equation}
The events \(\{J_i<q_i\}\) and \(F_i\) are disjoint, and
\eqref{eq:step-005-bad-event} gives
\begin{equation}\label{eq:step-005-bad-probability-split}
\Pr(\mathsf B_i)=\Pr(J_i<q_i)+\Pr(F_i).
\end{equation}
Summing \eqref{eq:step-005-factor-utility-integral} over the finitely many
positive-probability position atoms with at least \(q_i\) occurrences
gives
\begin{equation}\label{eq:step-005-marginal-failure}
\begin{split}
\Pr(F_i)
&=\sum_{\substack{a:\Pr(G=a)>0\\
                   |\{j:a_j=i\}|\ge q_i}}
  \Pr(G=a)
  \Pr\!\left(r_i(\bar H_i)>\frac1{64}\mid G=a\right)\\
&\le\frac1{4096}
  \sum_{\substack{a:\Pr(G=a)>0\\
                   |\{j:a_j=i\}|\ge q_i}}
  \Pr(G=a)\\
&=\frac1{4096}\Pr(J_i\ge q_i)
\le\frac1{4096}.
\end{split}
\end{equation}
This is a one-factor marginal calculation, not a bound on a union or a
factorization of joint probabilities.

For \(\rho_i=0\), \(\mathsf B_i=\varnothing\) and its weighted
probability is exactly zero.  Finite linearity, integrability from
Lemma~\ref{lem:step-005-measurable-bad-mass},
\eqref{eq:step-005-bad-probability-split}--
\eqref{eq:step-005-marginal-failure}, and
\(\sum_{i:\rho_i>0}\rho_i=1\) yield
\begin{equation}\label{eq:step-005-bad-versus-short}
\begin{split}
\mathbb EW_{\rm bad}
&=\sum_{i=1}^k\rho_i\Pr(\mathsf B_i)\\
&\le\sum_{i=1}^k\rho_i\Pr(J_i<q_i)
 +\frac1{4096}\sum_{i:\rho_i>0}\rho_i\\
&=\mathbb EW_{\rm short}+\frac1{4096}.
\end{split}
\end{equation}
Proposition~\ref{prop:step-004-weighted-shortage} supplies
\[
\mathbb EW_{\rm short}
\le\frac{128Q_\oplus}{n}+e^{-16},
\]
which proves \eqref{eq:step-005-bad-mass-expectation}.  Arbitrarily many
small-mass factors add no unweighted probability term.
\end{proof}

\begin{proposition}[Pointwise domination of exact global risk]
\label{prop:step-005-risk-domination}
Under Assumptions~\ref{assump:canonical-product} and
\ref{assump:countably-coded-evaluation} and
Lemma~\ref{lem:step-005-measurable-bad-mass}, for every fixed \(n,c,D\)
and every sample/output pair in the joint experiment,
\begin{equation}\label{eq:step-005-risk-domination}
R_D(h_{\bar H},c)\le\frac1{64}+W_{\rm bad}.
\end{equation}
The global risk is measurable.  The inequality holds on shortages and for
arbitrary factor outputs, without a simultaneous all-factor good event.
\end{proposition}

\begin{proof}
For \(\rho_i>0\), if \(\mathsf B_i\) does not occur, then
\(J_i\ge q_i\) and \(r_i(\bar H_i)\le1/64\); if it occurs, binary risk is
at most one.  Pointwise,
\begin{equation}\label{eq:step-005-local-risk-domination}
r_i(\bar H_i)
\le\frac1{64}\mathbf1_{\mathsf B_i^c}+\mathbf1_{\mathsf B_i}
=\frac1{64}+\frac{63}{64}\mathbf1_{\mathsf B_i}
\le\frac1{64}+\mathbf1_{\mathsf B_i}.
\end{equation}
This is a consequence of the bad-event definition, not a pointwise use of
the probabilistic utility theorem.

For every fixed quotient tuple \(\bar h\), block disjointness, the
positive-mass conditional laws, deterministic target labels, and quotient
pushforward give the exact identity
\begin{equation}\label{eq:step-005-exact-risk-decomposition}
\begin{split}
R_D(h_{\bar h},c)
&=\sum_{i:\rho_i>0}
D\bigl(\{x\in X_i:h_{\bar h}(x)\ne c(x)\}\bigr)\\
&=\sum_{i:\rho_i>0}\rho_i
\Pr_{q\sim\bar D_i}
[\bar h_i(q)\ne\bar c_i(q)]\\
&=\sum_{i:\rho_i>0}\rho_i r_i(\bar h_i).
\end{split}
\end{equation}
This is exact for arbitrary raw support.  Together with
\eqref{eq:step-005-factor-risk-map}, it also proves measurability of the
global risk as a function of \(\bar H\).

Apply \eqref{eq:step-005-local-risk-domination} in
\eqref{eq:step-005-exact-risk-decomposition}, use
\(\sum_{i:\rho_i>0}\rho_i=1\), and note that zero-mass factors have zero
weight and empty bad event.  Then
\[
R_D(h_{\bar H},c)
\le\sum_{i:\rho_i>0}\rho_i
\left(\frac1{64}+\mathbf1_{\mathsf B_i}\right)
=\frac1{64}+W_{\rm bad},
\]
which proves \eqref{eq:step-005-risk-domination} without requiring all
factors to be good simultaneously.
\end{proof}

\begin{proposition}[Universal numerical PAC closure and one-factor baseline]
\label{prop:step-005-pac-closure}
Under Assumptions~\ref{assump:canonical-product},
\ref{assump:countably-coded-evaluation}, and
\ref{assump:global-privacy-range},
Propositions~\ref{prop:step-002-factor-utility},
\ref{prop:step-004-weighted-shortage},
\ref{prop:step-005-weighted-bad-mass}, and
\ref{prop:step-005-risk-domination}, set
\begin{equation}\label{eq:step-005-cup}
C_{\rm up}=65536=2^{16}.
\end{equation}
For every \(c\in C\), probability measure \(D\) on \((X,\Sigma)\), and
integer \(n\ge\lceil C_{\rm up}Q_\oplus\rceil\),
\begin{equation}\label{eq:step-005-pac-conclusion}
\Pr_{S\sim D_c^n,\,\bar A_n^{\oplus,Q}}
\left[R_D(h_{\bar A_n^{\oplus,Q}(S)},c)\le\frac1{16}\right]
\ge\frac{15}{16}.
\end{equation}
If \(k=1\), the same threshold entails no shortage and the stronger
factor baseline
\begin{equation}\label{eq:step-005-one-factor-utility}
\Pr\left[R_D(h_{\bar H_1},c)\le\frac1{64}\right]
\ge1-\frac1{4096}
\end{equation}
holds directly.
\end{proposition}

\begin{proof}
Write \(Q=Q_\oplus\).  Since
\(n\ge\lceil65536Q\rceil\ge65536Q\),
Proposition~\ref{prop:step-005-weighted-bad-mass} gives
\begin{equation}\label{eq:step-005-numerical-expectation}
\mathbb EW_{\rm bad}
\le\frac{128}{65536}+e^{-16}+\frac1{4096}.
\end{equation}
The natural exponential satisfies
\(e=\sum_{m=0}^\infty1/m!>1+1=2\), so
\begin{equation}\label{eq:step-005-exp-comparison}
e^{-16}<2^{-16}=\frac1{65536}.
\end{equation}
Putting the terms in \eqref{eq:step-005-numerical-expectation} over the
common denominator \(65536\) gives
\begin{equation}\label{eq:step-005-expectation-threshold}
\mathbb EW_{\rm bad}
<\frac{128+1+16}{65536}
=\frac{145}{65536}
<\frac{192}{65536}
=\frac3{1024}.
\end{equation}

For completeness, nonnegativity gives the pointwise Markov inequality
\begin{equation}\label{eq:step-005-markov-indicator}
\mathbf1\left\{W_{\rm bad}>\frac3{64}\right\}
\le\frac{W_{\rm bad}}{3/64}.
\end{equation}
Taking expectations and using
\eqref{eq:step-005-expectation-threshold},
\begin{equation}\label{eq:step-005-markov-bound}
\Pr\left(W_{\rm bad}>\frac3{64}\right)
\le\frac{64}{3}\mathbb EW_{\rm bad}
<\frac{64}{3}\frac3{1024}=\frac1{16}.
\end{equation}
Proposition~\ref{prop:step-005-risk-domination} gives the exact event
inclusion
\begin{equation}\label{eq:step-005-risk-event-inclusion}
\left\{R_D(h_{\bar H},c)>\frac1{16}\right\}
\subseteq
\left\{W_{\rm bad}>\frac1{16}-\frac1{64}\right\}
=\left\{W_{\rm bad}>\frac3{64}\right\}.
\end{equation}
The strict boundary is correct: when \(W_{\rm bad}=3/64\), risk domination
gives risk at most \(1/16\), which is success under the closed event.
Equations~\eqref{eq:step-005-markov-bound}--
\eqref{eq:step-005-risk-event-inclusion} give failure probability strictly
below \(1/16\), hence the stated success probability, over the entire joint
law \eqref{eq:step-005-joint-law}.

For \(k=1\), one has \(X_1=X\), \(\rho_1=1\), and \(J_1=n\) for every
sample.  Since \(C_{\rm up}\ge1\) and \(Q_\oplus=q_1\), the threshold
implies \(n\ge q_1\).  The routed input is therefore the first \(q_1\)
iid records with no padding.  Their quotient images have the iid labeled
quotient law for arbitrary raw support, so
Proposition~\ref{prop:step-002-factor-utility} applies directly.  Identity
\eqref{eq:step-005-exact-risk-decomposition} reduces to
\(R_D(h_{\bar H_1},c)=r_1(\bar H_1)\), proving
\eqref{eq:step-005-one-factor-utility}.  The shortage bound and Markov
relaxation are not used in this specialization.
\end{proof}

\subsection{Public specialization of the quota sum}
\label{app:step-006}

\begin{lemma}[Endpoint-safe domination of heterogeneous logarithms]
\label{lem:step-006-log-domination}
Under Assumptions~\ref{assump:vc-one-factors} and
\ref{assump:global-privacy-range} and
Lemma~\ref{lem:step-001-logstar}, define
\begin{equation}\label{eq:step-006-log-definitions}
L_i=\log\!\left(\frac{e s_i}{\varepsilon\delta}\right),
\qquad
L=\log\!\left(\frac{eM}{\varepsilon\delta}\right),
\quad M=\sum_{i=1}^ks_i.
\end{equation}
Then all logarithm arguments are finite and strictly larger than \(20e\),
and
\begin{equation}\label{eq:step-006-log-domination}
2\le s_i\le M,\qquad
1<L_i\le L,\qquad
\sum_{i=1}^ks_iL_i^2\le ML^2.
\end{equation}
\end{lemma}

\begin{proof}
Lemma~\ref{lem:step-001-logstar} gives \(s_i\ge2\) and
\(M=\sum_{j=1}^ks_j\).  Since all summands are positive,
\begin{equation}\label{eq:step-006-si-upper}
M-s_i=\sum_{j\ne i}s_j\ge0,
\end{equation}
where the empty sum is zero when \(k=1\).  Hence \(2\le s_i\le M\).
The same lemma gives \(M\ge2k\), and thus \(M\ge2\).

Assumption~\ref{assump:global-privacy-range} gives
\(0<\varepsilon\delta<1/10\).  Therefore
\begin{equation}\label{eq:step-006-log-arguments}
\frac{es_i}{\varepsilon\delta}>20e>e,
\qquad
\frac{eM}{\varepsilon\delta}>20e>e.
\end{equation}
The natural logarithm is strictly increasing on \((0,\infty)\), so
\eqref{eq:step-006-si-upper} gives
\begin{equation}\label{eq:step-006-log-order}
1<\log(20e)<L_i\le L.
\end{equation}
Both logarithms are positive, so squaring preserves their order.  After
multiplying by the positive heterogeneous weights and summing,
\begin{equation}\label{eq:step-006-heterogeneous-sum}
\sum_{i=1}^ks_iL_i^2
\le\sum_{i=1}^ks_iL^2=ML^2.
\end{equation}
No balance among the \(s_i\)'s is used.

The endpoint checks are explicit in
\eqref{eq:step-006-log-arguments}: \(\varepsilon=1/10\) is allowed and
the arguments remain strictly above \(20e\) for every \(\delta<1\); as
either positive parameter approaches zero, the arguments increase.  The
excluded zero endpoints are never inserted into a logarithm.
\end{proof}

\begin{lemma}[Ceiling-resolved heterogeneous quota sum]
\label{lem:step-006-ceiling-sum}
Under Assumptions~\ref{assump:vc-one-factors} and
\ref{assump:global-privacy-range}, Lemma~\ref{lem:step-001-logstar}, and
Lemma~\ref{lem:step-006-log-domination}, each exact quota satisfies
\begin{equation}\label{eq:step-006-individual-ceiling}
q_i\le K_Y\frac{s_i}{\varepsilon}L_i^2+1,
\end{equation}
and
\begin{equation}\label{eq:step-006-quota-sum-with-k}
Q_\oplus
\le\frac{K_Y}{\varepsilon}\sum_{i=1}^ks_iL_i^2+k
\le K_Y\frac{M}{\varepsilon}L^2+k.
\end{equation}
\end{lemma}

\begin{proof}
If \(m=\lceil x\rceil\), then \(m-1<x\le m\), so
\(m<x+1\) and in particular \(\lceil x\rceil\le x+1\).  Apply this to
each finite positive quantity
\begin{equation}\label{eq:step-006-ceiling-input}
x_i=K_Y\frac{s_i}{\varepsilon}
\log^2\!\left(\frac{es_i}{\varepsilon\delta}\right)
=K_Y\frac{s_i}{\varepsilon}L_i^2.
\end{equation}
Because \(q_i=\lceil x_i\rceil\), this proves
\eqref{eq:step-006-individual-ceiling}, including when \(x_i\) is an
integer.  Summing one \(+1\) for each of the exactly \(k\) factors gives
the first inequality in \eqref{eq:step-006-quota-sum-with-k}; multiplying
\eqref{eq:step-006-heterogeneous-sum} by \(K_Y/\varepsilon>0\) gives the
second.  No ceiling has been discarded.
\end{proof}

\begin{proposition}[Public quota rate specialization]
\label{prop:step-006-public-quota-bridge}
Under Assumptions~\ref{assump:vc-one-factors} and
\ref{assump:global-privacy-range}, Lemma~\ref{lem:step-001-logstar}, and
Lemma~\ref{lem:step-006-ceiling-sum}, define the universal constant
\begin{equation}\label{eq:step-006-cquota}
C_{\rm quota}=\max\left\{1,K_Y+\frac1{20}\right\}.
\end{equation}
Then
\begin{equation}\label{eq:step-006-public-quota-rate}
Q_\oplus
\le C_{\rm quota}\frac{M}{\varepsilon}
\log^2\!\left(\frac{eM}{\varepsilon\delta}\right).
\end{equation}
The constant depends only on the fixed universal source constant \(K_Y\),
and the bound has neither a separate factor-count term nor a class,
quotient, or support cardinality term.
\end{proposition}

\begin{proof}
Let
\begin{equation}\label{eq:step-006-public-scale}
A=\frac{M}{\varepsilon}L^2
=\frac{M}{\varepsilon}
\log^2\!\left(\frac{eM}{\varepsilon\delta}\right)>0.
\end{equation}
Lemma~\ref{lem:step-001-logstar} gives \(M\ge2k\), hence
\begin{equation}\label{eq:step-006-k-half-m}
k\le\frac M2.
\end{equation}
Lemma~\ref{lem:step-006-log-domination} gives \(L>1\), while
Assumption~\ref{assump:global-privacy-range} gives
\(\varepsilon\le1/10\).  Therefore
\begin{equation}\label{eq:step-006-k-absorption}
k\le\frac M2
=\frac{\varepsilon}{2L^2}A
\le\frac1{20}A.
\end{equation}
Combining \eqref{eq:step-006-quota-sum-with-k} and
\eqref{eq:step-006-k-absorption},
\begin{equation}\label{eq:step-006-final-absorption}
Q_\oplus
\le K_YA+k
\le\left(K_Y+\frac1{20}\right)A
\le C_{\rm quota}A,
\end{equation}
which proves \eqref{eq:step-006-public-quota-rate}.

When every \(d_i=1\), the setting definition gives \(s_i=2\) and
\(M=2k\), so \eqref{eq:step-006-k-half-m} is tight but the same
absorption applies.  Heterogeneous factors are handled termwise in
\eqref{eq:step-006-heterogeneous-sum}.  When \(k=1\), \(M=s_1\) and
\(L_1=L\); the exact intermediate statement is
\(q_1\le K_Y(M/\varepsilon)L^2+1\), and
\eqref{eq:step-006-k-absorption} absorbs only that one ceiling.  Thus none
of these regimes introduces a cardinality term or changes the deterministic
fixed-sample mode.
\end{proof}

\subsection{An additive VC lower bound}
\label{app:step-007}

\begin{lemma}[Unseen labels against unrestricted randomized learners]
\label{lem:step-007-unseen-labels}
Under Assumptions~\ref{assump:canonical-product} and
\ref{assump:vc-one-factors} and
Lemma~\ref{lem:step-001-vc-cardinality}, fix \(n\in\mathbb N\) and any
unrestricted learner kernel \(A_n\) with measurable binary decoder
\(\omega\mapsto h_\omega\).  There are points \(x_i\in X_i\), targets
\(c^b\in C\) indexed by \(b\in\{0,1\}^k\), and the probability measure
\begin{equation}\label{eq:step-007-hard-distribution}
D=\frac1k\sum_{i=1}^k\delta_{x_i}
\end{equation}
such that, for uniform \(B\in\{0,1\}^k\), a sample
\(S\sim D_{c^B}^n\), and output from \(A_n(S,\cdot)\),
\begin{equation}\label{eq:step-007-unseen-risk}
\mathbb E_{B,S,A_n}R_D(h_{A_n(S)},c^B)
\ge\frac12\left(1-\frac1k\right)^n.
\end{equation}
No privacy, properness, or PAC premise is required.
\end{lemma}

\begin{proof}
Lemma~\ref{lem:step-001-vc-cardinality} supplies points \(x_i\in X_i\)
whose set \(W=\{x_1,\ldots,x_k\}\) is shattered.  The disjoint blocks make
them distinct.  For every \(b\in\{0,1\}^k\), fix one of the finitely many
full-product targets supplied by that lemma, so that
\begin{equation}\label{eq:step-007-target-labels}
c^b(x_i)=b_i,\qquad1\le i\le k.
\end{equation}
No measurable selector is involved because only finitely many targets are
fixed.

Represent the sample using independent indices
\(I_1,\ldots,I_n\), each uniform on \([k]=\{1,\ldots,k\}\):
\begin{equation}\label{eq:step-007-index-sample}
S_B(I)=\bigl((x_{I_t},B_{I_t})\bigr)_{t=1}^n.
\end{equation}
For \(j\in[k]\), let
\begin{equation}\label{eq:step-007-unseen-event}
U_j=\{I_t\ne j\text{ for every }1\le t\le n\}.
\end{equation}
Then
\begin{equation}\label{eq:step-007-unseen-probability}
\Pr(U_j)=\left(1-\frac1k\right)^n.
\end{equation}

For \(a\in\{0,1\}\), the finite-evaluation-cylinder requirement makes
\begin{equation}\label{eq:step-007-evaluation-events}
E_{j,a}=\{\omega:h_\omega(x_j)=a\}
\end{equation}
measurable.  The events \(E_{j,0}\) and \(E_{j,1}\) partition the entire
output space even for improper hypotheses.  Fix an index tuple on which
\(U_j\) occurs and all target bits except \(b_j\).  The dataset
\eqref{eq:step-007-index-sample} is identical for \(b_j=0\) and \(b_j=1\);
call it \(S\).  Averaging over the fair bit and integrating all learner
randomness through the same kernel gives
\begin{equation}\label{eq:step-007-fair-unseen-bit}
\begin{split}
&\frac12\Pr_{A_n(S)}[h_{A_n(S)}(x_j)\ne0]
+\frac12\Pr_{A_n(S)}[h_{A_n(S)}(x_j)\ne1]\\
&\qquad=\frac12A_n(S,E_{j,1})+\frac12A_n(S,E_{j,0})
=\frac12.
\end{split}
\end{equation}
Thus arbitrary output randomization and coupling cannot reveal an unseen
target bit.

Under \eqref{eq:step-007-hard-distribution}, the risk is the measurable
finite-evaluation sum
\begin{equation}\label{eq:step-007-finite-risk}
R_D(h_\omega,c^B)
=\frac1k\sum_{j=1}^k
\mathbf1\{h_\omega(x_j)\ne B_j\}.
\end{equation}
Dropping the nonnegative losses on seen points, applying
\eqref{eq:step-007-fair-unseen-bit} on \(U_j\), and using
\eqref{eq:step-007-unseen-probability},
\begin{equation}\label{eq:step-007-risk-sum}
\begin{split}
\mathbb ER_D(h_{A_n(S)},c^B)
&=\frac1k\sum_{j=1}^k
\Pr[h_{A_n(S)}(x_j)\ne B_j]\\
&\ge\frac1k\sum_{j=1}^k
\Pr[U_j,\ h_{A_n(S)}(x_j)\ne B_j]\\
&=\frac1k\sum_{j=1}^k\frac12\Pr(U_j)
=\frac12\left(1-\frac1k\right)^n.
\end{split}
\end{equation}
Repeated sample points are already included and can only leave more
coordinates unseen.
\end{proof}

\begin{lemma}[Exact PAC-event to expectation conversion]
\label{lem:step-007-pac-conversion}
Let \(L\) be measurable with \(0\le L\le1\).  If
\begin{equation}\label{eq:step-007-pac-event}
\Pr[L\le1/16]\ge15/16,
\end{equation}
then
\begin{equation}\label{eq:step-007-expectation-ceiling}
\mathbb EL\le\frac{31}{256}.
\end{equation}
Consequently,
\begin{equation}\label{eq:step-007-strict-conversion}
\mathbb EL>\frac{31}{256}
\quad\Longrightarrow\quad
\Pr[L>1/16]>1/16.
\end{equation}
\end{lemma}

\begin{proof}
Put \(p=\Pr[L>1/16]\).  On the complementary event \(L\le1/16\), and
everywhere \(L\le1\), so
\begin{equation}\label{eq:step-007-loss-split}
\mathbb EL\le\frac1{16}(1-p)+p
=\frac1{16}+\frac{15}{16}p.
\end{equation}
Under \eqref{eq:step-007-pac-event}, \(p\le1/16\), whence
\[
\mathbb EL\le\frac1{16}+\frac{15}{16}\frac1{16}
=\frac{31}{256}.
\]
The strict contrapositive proves \eqref{eq:step-007-strict-conversion}.
The good-event equality boundary is unchanged.
\end{proof}

\begin{proposition}[Additive VC PAC lower certificate]
\label{prop:step-007-additive-vc}
Under Assumptions~\ref{assump:canonical-product} and
\ref{assump:vc-one-factors}, Lemma~\ref{lem:step-001-vc-cardinality},
Lemma~\ref{lem:step-007-unseen-labels}, and
Lemma~\ref{lem:step-007-pac-conversion}, fix \(n\in\mathbb N\).  If an
unrestricted learner \(A_n\) satisfies, for every \(c\in C\) and every
allowed probability measure \(D'\),
\begin{equation}\label{eq:step-007-global-pac-premise}
\Pr_{S\sim(D'_c)^n,A_n}
\left[R_{D'}(h_{A_n(S)},c)\le\frac1{16}\right]
\ge\frac{15}{16},
\end{equation}
then
\begin{equation}\label{eq:step-007-additive-lower}
n\ge\frac{k}{2}.
\end{equation}
More precisely, if \(n<k/2\), the distribution
\eqref{eq:step-007-hard-distribution} and a deterministic target
\(c^{b_*}\in C\) satisfy
\begin{equation}\label{eq:step-007-deterministic-failure}
\Pr_{S\sim D_{c^{b_*}}^n,A_n}
\left[R_D(h_{A_n(S)},c^{b_*})>\frac1{16}\right]
>\frac1{16}.
\end{equation}
\end{proposition}

\begin{proof}
Since \(n\in\mathbb N=\{1,2,\ldots\}\), for \(k=1\) one has
\(n\ge1\ge k/2\), and for \(k=2\), \(n\ge1=k/2\).  Only \(k\ge3\)
can enter the contradiction branch \(n<k/2\).

For every integer \(m\ge1\) and \(x\in[0,1]\),
\begin{equation}\label{eq:step-007-bernoulli-inequality}
(1-x)^m\ge1-mx.
\end{equation}
It is equality at \(m=1\).  If it holds at \(m\), multiplication by
\(1-x\ge0\) gives
\[
(1-x)^{m+1}
\ge(1-mx)(1-x)
=1-(m+1)x+mx^2
\ge1-(m+1)x,
\]
proving \eqref{eq:step-007-bernoulli-inequality} by induction.

Suppose \(n<k/2\).  Lemma~\ref{lem:step-007-unseen-labels} and
\eqref{eq:step-007-bernoulli-inequality} with \(x=1/k\) imply
\begin{equation}\label{eq:step-007-strict-average-risk}
\mathbb E_{B,S,A_n}R_D(h_{A_n(S)},c^B)
\ge\frac12\left(1-\frac1k\right)^n
\ge\frac12\left(1-\frac nk\right)
>\frac14>\frac{31}{256}.
\end{equation}
For each \(b\), define
\begin{equation}\label{eq:step-007-target-loss}
\ell_b=\mathbb E_{S\sim D_{c^b}^n,A_n}
R_D(h_{A_n(S)},c^b).
\end{equation}
The left side of \eqref{eq:step-007-strict-average-risk} is the finite
average \(2^{-k}\sum_b\ell_b\).  Hence some deterministic \(b_*\) has
\begin{equation}\label{eq:step-007-target-extraction}
\ell_{b_*}>\frac14>\frac{31}{256}.
\end{equation}
Lemma~\ref{lem:step-007-pac-conversion}, applied contrapositively to the
\([0,1]\)-valued risk in \eqref{eq:step-007-target-loss}, proves
\eqref{eq:step-007-deterministic-failure}.  This contradicts
\eqref{eq:step-007-global-pac-premise}; therefore \(n<k/2\) is impossible.
Thus \eqref{eq:step-007-additive-lower} holds with the universal constant
\(a_{\rm VC}=1/2\), without privacy, properness, or output-space
restriction beyond measurable binary evaluation.
\end{proof}

\subsection{An expected-risk Littlestone floor for one factor}
\label{app:step-008}

For a labeled database \(U=((u_a,y_a))_{a=1}^N\) and a decoded binary
hypothesis \(h\), write
\[
L_U(h)=\frac1N\sum_{a=1}^N\mathbf1\{h(u_a)\ne y_a\}.
\]
When \(y_a=c(u_a)\), we also write \(L_U(h,c)\).  In this subsection
\([t]=\{1,\ldots,t\}\) and \(\ln\) denotes the natural logarithm.

\begin{proposition}[Expected empirical threshold obstruction]
\label{prop:step-008-almm-expected}
There is a universal \(a_0>0\) with the following property.  Let \(V\)
be a finite totally ordered set of size \(t\), let even \(N\ge72\), and
let
\[
B:(V\times\{0,1\})^N\rightsquigarrow\{0,1\}^{V}
\]
be an arbitrary-output randomized algorithm.  Suppose that
\begin{enumerate}
\item \(B\) is replacement-\((0.1,\eta)\)-DP on all labeled databases;
\item \(\eta\le1/(1000N^2\log_2N)\); and
\item for every threshold target \(c\) on \(V\) and every \(N\)-row
database \(U\) labeled by \(c\),
\begin{equation}\label{eq:step-008-empirical-loss-premise}
\mathbb E_B L_U(B(U),c)\le\frac18.
\end{equation}
\end{enumerate}
Then
\begin{equation}\label{eq:step-008-almm-obstruction}
N\ge a_0\log_2^*t.
\end{equation}
The output may be any bit vector, not necessarily a threshold.
\end{proposition}

\begin{proof}
We adapt the active internal reduction proving the finite-threshold lower
bound of \citet{alonLivniMalliarisMoran2019}; its printed PAC corollary is
not used as an expected-risk theorem.

For an \(N\)-row database \(U\) and \(v\in V\), put
\[
b_U(v)=\Pr[B(U)(v)=1].
\]
For an increasing database on distinct points, let
\(\operatorname{ord}_U(v)\) be the number of its instance coordinates at
or below \(v\).  Call \(K\subseteq V\) \(N\)-homogeneous if there are
\(p_0,\ldots,p_N\in[0,1]\) such that every increasing balanced
threshold-labeled database \(U\) on \(N\) distinct points of \(K\) and
every \(v\in K\setminus U_X\) satisfy
\begin{equation}\label{eq:step-008-homogeneity}
\left|b_U(v)-p_{\operatorname{ord}_U(v)}\right|
\le\frac1{100N}.
\end{equation}

The finite Ramsey lemma of \citet{alonLivniMalliarisMoran2019}, applied to
the \(N+1\) output-coordinate marginals rounded to multiples of
\(1/(100N)\), gives a universal \(C_R\) and an \(N\)-homogeneous set
\(K\) with
\begin{equation}\label{eq:step-008-ramsey-lower}
k:=|K|\ge
\frac{\log_2^{(N)}t}{2^{C_RN\log_2N}}
\end{equation}
whenever \(\log_2^*t>N\).  There are at most
\((100N+1)^{N+1}\) colors; the cited finite hypergraph Ramsey bound gives
the display after changing the fixed logarithm base and enlarging \(C_R\).
Only finitely many coordinate marginals are used, so arbitrary improper
bit-vector output is admissible.

We next expose the sole utility input to the cited reduction.  On its
balanced increasing realizable database \(U\),
\eqref{eq:step-008-empirical-loss-premise} gives
\begin{equation}\label{eq:step-008-balanced-accuracy}
\frac78
\le\mathbb E[1-L_U(B(U),c)]
=\frac1N\sum_{a\le N/2}(1-b_U(u_a))
 +\frac1N\sum_{a>N/2}b_U(u_a).
\end{equation}
Thus some positive coordinate has marginal at least \(3/4\), and some
negative coordinate has marginal at most \(1/4\).  Replace the selected
coordinate by the adjacent unused point in the reduction, retaining its
label.  Reverse replacement privacy gives, for the high marginal,
\begin{equation}\label{eq:step-008-high-marginal}
b_{U'}(u_a)
\ge e^{-0.1}(b_U(u_a)-\eta)
\ge e^{-0.1}(3/4-\eta)>2/3,
\end{equation}
whereas forward privacy gives, for the low marginal,
\begin{equation}\label{eq:step-008-low-marginal}
b_{U''}(u_a)
\le e^{0.1}b_U(u_a)+\eta
\le e^{0.1}/4+\eta<1/3.
\end{equation}
Both strict estimates follow from \(N\ge72\) and
\begin{equation}\label{eq:step-008-small-eta}
\eta\le\frac1{1000N^2\log_2N}<0.001.
\end{equation}
The removed points lie outside the modified databases, so homogeneity maps
\eqref{eq:step-008-high-marginal}--
\eqref{eq:step-008-low-marginal} to entries of the same list: the upper
entry is at least \(2/3-1/(100N)\), the lower at most
\(1/3+1/(100N)\), and their ordered index gap is at most \(N\).
Telescoping yields
\begin{equation}\label{eq:step-008-probability-jump}
\exists j\in\{1,\ldots,N\}:\qquad
p_j-p_{j-1}\ge
\frac1N\left(\frac13-\frac1{50N}\right)
\ge\frac1{4N}.
\end{equation}
This is exactly where the cited printed proof used PAC utility: its PAC
premise served only to imply an expected empirical loss at most \(1/8\),
which is already \eqref{eq:step-008-empirical-loss-premise}.

Suppose \(k>2N+1\) and put \(n=k-N\).  The probability-to-distribution
lemma of \citet{alonLivniMalliarisMoran2019}, instantiated with the jump
\eqref{eq:step-008-probability-jump}, chooses a balanced database with one
movable row and an interval \(J\subseteq K\) of exactly \(n\) unused
points.  For \(a\in[n]\), place the movable row at the \(a\)-th point of
\(J\), call the database \(U_a\), and let \(P_a\) be the law on
\(\{0,1\}^n\) obtained by restricting \(B(U_a)\) to \(J\).  Every pair
\(P_a,P_{a'}\) is \((0.1,\eta)\)-indistinguishable because the databases
differ in one row.  With \(r=(p_{j-1}+p_j)/2\), homogeneity and the jump
give
\begin{equation}\label{eq:step-008-cube-margin}
\begin{aligned}
\Pr_{z\sim P_a}[z(b)=1]&\le r-\frac1{10N},&&b<a,\\
\Pr_{z\sim P_a}[z(b)=1]&\ge r+\frac1{10N},&&b>a.
\end{aligned}
\end{equation}
Indeed, the jump \(1/(4N)\), minus two homogeneity errors, leaves more
than \(1/(5N)\) between the sides; taking the midpoint leaves the displayed
margin.  The output vectors remain arbitrary members of the full cube, and
all-input privacy covers every database used here.

The binary lemma and product-composition lemma of
\citet{alonLivniMalliarisMoran2019} imply
\begin{equation}\label{eq:step-008-binary-upper}
n\le2^{1000N^2(\log_2N)^2}.
\end{equation}
We include their quantitative contradiction to discharge the additive
privacy scale.  Set \(L=\log_2N\), suppose the displayed bound fails, and
define
\begin{equation}\label{eq:step-008-binary-parameters}
T=\left\lfloor1000N^2L^2\right\rfloor-1,
\qquad
D=\left\lceil50N^2\ln(6T)\right\rceil.
\end{equation}
For \(D\) independent cube vectors, threshold each coordinate's empirical
fraction of ones at \(r\), breaking ties toward zero, and perform \(T\)
binary-search decisions.  Since \(n>2^{T+1}\), each of the \(2^T\)
terminal intervals contains at least two indices.  Choose in every interval
a representative \(a\) unequal to every coordinate queried on its path,
and let \(E_a\) be the event that the search follows that path.  These
events are disjoint.  Under \(P_a^D\), each queried coordinate is at
distance at least \(1/(10N)\) from \(r\) on the side leading toward
\(a\).  Hoeffding's inequality \citep{hoeffding1963} and a union bound give
\begin{equation}\label{eq:step-008-path-probability}
P_a^D(E_a)
\ge1-T\exp\!\left(-\frac{D}{50N^2}\right)
\ge\frac56.
\end{equation}

For every pair of product laws, the cited product-composition result gives
indistinguishability with parameters \((0.1D,D\eta)\).  From
\eqref{eq:step-008-binary-parameters} and
\eqref{eq:step-008-small-eta},
\begin{equation}\label{eq:step-008-product-delta}
\begin{split}
D\eta
&\le\frac{\ln(6T)}{20L}+\frac1{1000N^2L}\\
&\le\frac{\ln(6000)+2\ln N+2\ln L}{20L}\\
&\qquad+\frac1{1000N^2L}<\frac15.
\end{split}
\end{equation}
The last inequality holds at \(N=72\) and its left side decreases
thereafter: \(\ln N/L=\ln2\) is constant, while \(1/L\) and
\(\ln L/L\) decrease for \(L\ge\log_2 72>e\).  Thus for any fixed
\(P_b^D\), reverse product privacy and
\eqref{eq:step-008-path-probability} imply
\begin{equation}\label{eq:step-008-reverse-product-privacy}
P_b^D(E_a)
\ge e^{-0.1D}\bigl(P_a^D(E_a)-D\eta\bigr)
>\frac12e^{-0.1D}.
\end{equation}
The same monotonic calculation gives \(\ln(6T)\le3.4L\), so
\(0.1D\le17N^2L+0.1\).  Since
\[
(T-1)\ln2
\ge(1000N^2L^2-3)\ln2
>17N^2L+0.1,
\]
the disjoint events contradict total mass:
\[
1\ge\sum_aP_b^D(E_a)
>2^T\frac12e^{-0.1D}>1.
\]
This proves \eqref{eq:step-008-binary-upper} without any further utility
input.

Take \(C_H=1001\).  If \(k\le2N+1\), then
\(k\le2^{C_HN^2L^2}\) immediately.  Otherwise
\eqref{eq:step-008-probability-jump},
\eqref{eq:step-008-cube-margin}, and
\eqref{eq:step-008-binary-upper} give
\begin{equation}\label{eq:step-008-homogeneous-upper}
k=N+n
\le N+2^{1000N^2L^2}
\le2^{C_HN^2L^2}.
\end{equation}
If \(\log_2^*t\le N\), the proposition follows after choosing
\(a_0\le1\).  Otherwise combine
\eqref{eq:step-008-ramsey-lower} and
\eqref{eq:step-008-homogeneous-upper}:
\begin{equation}\label{eq:step-008-iterate-upper}
\log_2^{(N)}t
\le k\,2^{C_RNL}
\le2^{C_1N^2L^2}
\end{equation}
for a universal \(C_1\).  At most
\(\log_2^*(2^{C_1N^2L^2})+O(1)=\log_2^*N+O(1)\) further logarithms reduce
the right side to at most one.  Thus
\[
\log_2^*t\le N+\log_2^*N+O(1)\le C_2N
\]
for a universal \(C_2\) and every \(N\ge72\).  Taking
\(a_0=1/C_2\), and decreasing it for the finitely many fixed-base
endpoints if needed, proves \eqref{eq:step-008-almm-obstruction}.

Evenness is used only for the balanced databases.  The fixed relabeling
between \(\{0,1\}\) and \(\{-1,+1\}\) preserves loss, replacement
adjacency, marginals, and privacy.  All cited reductions inspect only
finite coordinate marginals or restrictions, so the arbitrary-output
scope is unchanged.
\end{proof}

\begin{lemma}[Finite Shelah tasks and the log-star bridge]
\label{lem:step-008-shelah-tasks}
Under Assumption~\ref{assump:vc-one-factors} and
Lemma~\ref{lem:step-001-quotient-invariance}, fix \(i\) and put
\begin{equation}\label{eq:step-008-witness-scales}
r_i=\log_2^*d_i,
\qquad
t_i=\max\{1,\lfloor\ln d_i\rfloor\}.
\end{equation}
There are fixed quotient points \(q_{i,1},\ldots,q_{i,t_i}\), quotient
concepts \(\bar c_{i,1},\ldots,\bar c_{i,t_i}\), and raw representatives
\(x_{i,j}\in\kappa_i^{-1}(\{q_{i,j}\})\) such that, for
\(c_{i,\ell}=\bar c_{i,\ell}\circ\kappa_i\in C_i\),
\begin{equation}\label{eq:step-008-threshold-pattern}
\bar c_{i,\ell}(q_{i,j})=c_{i,\ell}(x_{i,j})
=\mathbf1\{\ell\le j\}.
\end{equation}
For \(m\ge8\), let \(N_m\) be the least even integer at least \(9m\).
For \(\mathbf j=(j_1,\ldots,j_{N_m})\in[t_i]^{N_m}\), define
\begin{equation}\label{eq:step-008-task-measure}
\mu_{i,\mathbf j}=\frac1{N_m}\sum_{a=1}^{N_m}
\delta_{x_{i,j_a}}.
\end{equation}
Then
\begin{equation}\label{eq:step-008-task-universe}
\mathfrak T_{i,m}
=\{(c_{i,\ell},\mu_{i,\mathbf j}):
\ell\in[t_i],\ \mathbf j\in[t_i]^{N_m}\}
\end{equation}
is a finite learner-independent universe of realizable factor tasks.  If
\(r_i\ge8\), then
\begin{equation}\label{eq:step-008-logstar-shift}
\log_2^*t_i\ge r_i-2.
\end{equation}
\end{lemma}

\begin{proof}
Lemma~\ref{lem:step-001-quotient-invariance} gives
\(\operatorname{LD}(\bar C_i)=d_i\).  Shelah's threshold theorem in
\citet{alonLivniMalliarisMoran2019} gives a threshold pattern of length
\(\lfloor\log d_i\rfloor\).  If the source logarithm is natural,
\eqref{eq:step-008-witness-scales} uses no greater length; if it is base
two, \(\ln d_i\le\log_2d_i\), so a prefix has the required length.  When
\(\lfloor\ln d_i\rfloor=0\), nonconstancy supplies a quotient point with
both labels available, and a concept labeled one gives the length-one
pattern.  Quotient pullback and finitely many raw representative choices
give \eqref{eq:step-008-threshold-pattern} without a measurable global
section.

Every \(\mu_{i,\mathbf j}\) is a probability measure on the raw measurable
space: a Dirac measure is defined on any measurable space, even if its
singleton is not itself measurable.  The targets are measurable setting
concepts.  Thus all tasks are legal and realizable, and the family has at
most \(t_i^{N_m+1}\) members fixed before the learner.

For the shift, let \(x=\log_2d_i\).  If \(r_i\ge8\), then \(d_i>1024\)
and \(x>10\).  Since \(x\ln2>2\),
\begin{equation}\label{eq:step-008-threshold-length}
t_i=\lfloor x\ln2\rfloor\ge\frac{\ln2}{2}x.
\end{equation}
At \(x=10\), \((\ln2)x/2>\log_2x\); the derivative of the left side is
\(\ln2/2\), while that of \(\log_2x\) is
\(1/(x\ln2)<\ln2/2\) for \(x\ge10\).  Hence
\(t_i\ge\log_2x\).  Monotonicity and two applications of the log-star
recursion give
\[
\log_2^*t_i
\ge\log_2^*(\log_2x)
=\log_2^*x-1=r_i-2,
\]
proving \eqref{eq:step-008-logstar-shift}.
\end{proof}

\begin{lemma}[Replacement-private resampling and exact risk identity]
\label{lem:step-008-resampling}
Fix the factor and witness of Lemma~\ref{lem:step-008-shelah-tasks}.  Let
\(m\ge8\), \(\delta>0\), and let
\[
A:(X_i\times\{0,1\})^m\rightsquigarrow(\Omega,\mathscr F)
\]
be total and replacement-\((0.1,\delta)\)-DP on every labeled input, with
measurable finite-evaluation cylinders.  Define \(N_m\) as above and set
\begin{equation}\label{eq:step-008-cdelta}
c_\delta=\frac{\ln2}{200000}.
\end{equation}
If
\begin{equation}\label{eq:step-008-candidate-delta}
0<\delta\le\frac{c_\delta}{m^2\ln(m+1)},
\end{equation}
then the following algorithm \(B_{i,m}\) is well defined on every labeled
\(N_m\)-row database over the raw witness: sample \(m\) row indices iid
uniformly with replacement, run \(A\) on those rows, and output its decoded
labels on \((x_{i,1},\ldots,x_{i,t_i})\).  It is replacement-
\((\varepsilon_m',\delta_m')\)-DP with
\begin{equation}\label{eq:step-008-resampled-privacy}
\varepsilon_m'=\frac{0.6m}{N_m}\le\frac1{15},
\qquad
\delta_m'=e^{0.6m/N_m}\frac{4m}{N_m}\delta
<\delta
\le\frac1{1000N_m^2\log_2N_m}.
\end{equation}
For every \((c_{i,\ell},\mu_{i,\mathbf j})\in\mathfrak T_{i,m}\), let
\[
U_{i,\ell,\mathbf j}
=\bigl((x_{i,j_a},c_{i,\ell}(x_{i,j_a}))\bigr)_{a=1}^{N_m}.
\]
Then
\begin{equation}\label{eq:step-008-risk-identity}
\mathbb E L_{U_{i,\ell,\mathbf j}}
\bigl(B_{i,m}(U_{i,\ell,\mathbf j}),c_{i,\ell}\bigr)
=\mathbb E_{S\sim(\mu_{i,\mathbf j})_{c_{i,\ell}}^m,A}
R_{\mu_{i,\mathbf j}}(h_{A(S)},c_{i,\ell}).
\end{equation}
At \(m=8\), \(N_m=72\).
\end{lemma}

\begin{proof}
Finite-evaluation measurability makes the output vector legal; it can be
arbitrary and need not be a threshold.  Totality on nonrealizable inputs is
inherited from \(A\).  By construction,
\begin{equation}\label{eq:step-008-even-rounding}
9m\le N_m\le9m+1\le10m.
\end{equation}
In particular \(N_m\ge2m\).  The with-replacement privacy-amplification
calculation preceding Lemma~5.9 of
\citet{bunNissimStemmerVadhan2015} applies to replacement-adjacent
databases, including every multiplicity of the changed row.  With base
privacy \(0.1\), it gives exactly the two parameters in
\eqref{eq:step-008-resampled-privacy}.  From \(N_m\ge9m\),
\(\varepsilon_m'\le0.6/9=1/15\).  For \(0\le y<1\), the exponential
series gives \(e^y\le(1-y)^{-1}\), so
\begin{equation}\label{eq:step-008-delta-contraction}
e^{0.6m/N_m}\frac{4m}{N_m}
\le e^{1/15}\frac49
\le\frac{15}{14}\frac49
=\frac{10}{21}<1.
\end{equation}
Thus \(\delta_m'<\delta\), and monotonicity also regards the resampled
algorithm as \((0.1,\delta_m')\)-DP.

For \(m\ge8\),
\begin{equation}\label{eq:step-008-square-domination}
10m\le(m+1)^2.
\end{equation}
Equations~\eqref{eq:step-008-even-rounding}--
\eqref{eq:step-008-square-domination} give
\begin{equation}\label{eq:step-008-delta-scale-translation}
N_m^2\log_2N_m
=\frac{N_m^2\ln N_m}{\ln2}
\le\frac{100m^2\,2\ln(m+1)}{\ln2}.
\end{equation}
Inverting this display and using
\eqref{eq:step-008-cdelta}--\eqref{eq:step-008-candidate-delta} yields
\[
\delta_m'<\delta
\le\frac{\ln2}{200000m^2\ln(m+1)}
\le\frac1{1000N_m^2\log_2N_m},
\]
completing the privacy claims.  At \(m=8\), \(N_m=72\) and
\[
\varepsilon_8'=\frac{0.6\cdot8}{72}=\frac1{15},
\qquad
\delta_8'=e^{1/15}\frac49\delta
\le\frac{10}{21}\delta<\delta,
\]
so no open endpoint is hidden.

Conditional on \(U_{i,\ell,\mathbf j}\), the resample is exactly iid from
\((\mu_{i,\mathbf j})_{c_{i,\ell}}\).  For every decoded \(h\), before
expectation,
\begin{equation}\label{eq:step-008-pointwise-risk-identity}
\begin{split}
L_{U_{i,\ell,\mathbf j}}(h,c_{i,\ell})
&=\frac1{N_m}\sum_{a=1}^{N_m}
\mathbf1\{h(x_{i,j_a})\ne c_{i,\ell}(x_{i,j_a})\}\\
&=R_{\mu_{i,\mathbf j}}(h,c_{i,\ell}).
\end{split}
\end{equation}
Taking expectation over sampled indices and learner randomness proves
\eqref{eq:step-008-risk-identity}; repeated rows retain their multiplicity
on both sides.
\end{proof}

\begin{proposition}[Unrestricted expected-risk factor floor]
\label{prop:step-008-factor-floor}
Under Assumptions~\ref{assump:vc-one-factors} and
\ref{assump:global-privacy-range}, Lemma~\ref{lem:step-001-quotient-invariance},
Proposition~\ref{prop:step-008-almm-expected},
Lemma~\ref{lem:step-008-shelah-tasks}, and
Lemma~\ref{lem:step-008-resampling}, fix \(i\), \(m\ge8\), \(\delta>0\),
and a total arbitrary-output learner
\[
A:(X_i\times\{0,1\})^m\rightsquigarrow(\Omega,\mathscr F)
\]
that is replacement-\((0.1,\delta)\)-DP on all labeled inputs and has
measurable finite evaluations.  For \(a_0\) from
Proposition~\ref{prop:step-008-almm-expected}, define
\begin{equation}\label{eq:step-008-floor-constants}
a_{\rm L}=\min\left\{1,\frac{a_0}{20}\right\},
\qquad
c_\delta=\frac{\ln2}{200000}.
\end{equation}
If \(\delta\le c_\delta/(m^2\ln(m+1))\) and \(A\) has expected risk at
most \(1/8\) on every realizable factor task, then
\begin{equation}\label{eq:step-008-factor-floor}
m\ge a_{\rm L}\log_2^*d_i.
\end{equation}
More strongly, for the fixed finite universe
\(\mathfrak T_{i,m}\) of Lemma~\ref{lem:step-008-shelah-tasks},
\begin{equation}\label{eq:step-008-strict-factor-floor}
m<a_{\rm L}\log_2^*d_i
\quad\Longrightarrow\quad
\exists(c,D)\in\mathfrak T_{i,m}:
\mathbb E_{S\sim D_c^m,A}R_D(h_{A(S)},c)>\frac18.
\end{equation}
\end{proposition}

\begin{proof}
Write \(r_i=\log_2^*d_i\).  If \(r_i<8\), then
\(m\ge8>r_i\), and \eqref{eq:step-008-factor-floor} follows from
\(a_{\rm L}\le1\).

Suppose \(r_i\ge8\).  If the learner has expected risk at most \(1/8\)
on every realizable task, it does so on every member of
\(\mathfrak T_{i,m}\).  Lemma~\ref{lem:step-008-resampling} constructs an
even-\(N_m\) arbitrary-output empirical threshold learner satisfying all
hypotheses of Proposition~\ref{prop:step-008-almm-expected}.  That
proposition, Lemma~\ref{lem:step-008-shelah-tasks}, and \(N_m\le10m\)
give
\begin{equation}\label{eq:step-008-final-floor-calculation}
10m\ge N_m
\ge a_0\log_2^*t_i
\ge a_0(r_i-2)
\ge\frac{a_0}{2}r_i.
\end{equation}
Thus \(m\ge(a_0/20)r_i\ge a_{\rm L}r_i\).

If instead \(m<a_{\rm L}r_i\), the low-\(r_i\) case is impossible.  In
the remaining case, were every task in \(\mathfrak T_{i,m}\) to have
expected risk at most \(1/8\), the same calculation
\eqref{eq:step-008-final-floor-calculation} would contradict the strict
hypothesis.  Therefore logical negation gives one task with risk strictly
greater than \(1/8\), proving
\eqref{eq:step-008-strict-factor-floor}.  The finite task universe was
fixed before \(A\).
\end{proof}

\subsection{Candidate-wise hard-factor calibration}
\label{app:step-009}

\begin{lemma}[Low-factor mass calibration]
\label{lem:step-009-low-mass}
Under Assumptions~\ref{assump:canonical-product} and
\ref{assump:vc-one-factors}, Lemma~\ref{lem:step-001-logstar},
Proposition~\ref{prop:step-007-additive-vc}, and
Proposition~\ref{prop:step-008-factor-floor}, fix one candidate
\(n\in\mathbb N\) and a learner satisfying the global
\((1/16,1/16)\)-PAC guarantee.  Put
\(a_{\rm VC}=1/2\), take \(a_{\rm L}\in(0,1]\) from
Proposition~\ref{prop:step-008-factor-floor}, and define the universal
constants
\begin{equation}\label{eq:step-009-rzero}
R_0=\left\lceil\frac{32}{a_{\rm L}}\right\rceil
\end{equation}
and
\begin{equation}\label{eq:step-009-clow}
c_{\rm low}
=\min\left\{
\frac{a_{\rm VC}}{128(R_0+2)},\frac{a_{\rm L}}{64}
\right\}
=\min\left\{
\frac1{256(R_0+2)},\frac{a_{\rm L}}{64}
\right\}>0.
\end{equation}
With
\begin{equation}\label{eq:step-009-hard-set}
\pi_i=\omega_i=\frac{s_i}{M},
\qquad
H=\{i:r_i>R_0\},
\qquad
w_L=\sum_{i\notin H}\pi_i,
\end{equation}
if \(n<c_{\rm low}M\), then
\begin{equation}\label{eq:step-009-low-mass}
w_L<\frac1{128}.
\end{equation}
In particular, \(H\ne\varnothing\).
\end{lemma}

\begin{proof}
Lemma~\ref{lem:step-001-logstar} gives \(s_i\le r_i+2\).  For
\(i\notin H\), including \(r_i=R_0\),
\begin{equation}\label{eq:step-009-low-factor-size}
s_i\le r_i+2\le R_0+2.
\end{equation}
The PAC premise and Proposition~\ref{prop:step-007-additive-vc} give
\begin{equation}\label{eq:step-009-vc-factor-count}
n\ge a_{\rm VC}k,
\qquad k\le\frac n{a_{\rm VC}}.
\end{equation}
Since every \(s_i\ge2\), \(M>0\).  Summing
\eqref{eq:step-009-low-factor-size}, applying
\eqref{eq:step-009-vc-factor-count}, and only then using the strict
contradiction hypothesis gives
\begin{equation}\label{eq:step-009-low-mass-chain}
\begin{split}
w_L
&=\frac1M\sum_{i\notin H}s_i
\le\frac{(R_0+2)k}{M}\\
&\le\frac{R_0+2}{a_{\rm VC}}\frac nM
<\frac{R_0+2}{a_{\rm VC}}c_{\rm low}
\le\frac1{128}.
\end{split}
\end{equation}
Strictness comes from \(n<c_{\rm low}M\), even when the first branch in
\eqref{eq:step-009-clow} is attained.  Because
\(\sum_i\pi_i=1\), an empty \(H\) would give \(w_L=1\), a contradiction.
If every factor is active, the empty low-factor sum is zero.  Factors with
\(r_i=0\) or \(r_i=R_0\) are included in the low set and already paid for
by \eqref{eq:step-009-low-mass-chain}.
\end{proof}

\begin{lemma}[Strict active-factor budget calibration]
\label{lem:step-009-active-budget}
Under Lemma~\ref{lem:step-001-logstar},
Proposition~\ref{prop:step-008-factor-floor}, and
Lemma~\ref{lem:step-009-low-mass}, use the exact constants
\(a_{\rm VC}=1/2\), \(0<a_{\rm L}\le1\), \(R_0\), \(c_{\rm low}\), and
the set \(H\) in \eqref{eq:step-009-rzero}--
\eqref{eq:step-009-hard-set}.  At the same candidate, if
\(n<c_{\rm low}M\), then every \(i\in H\) satisfies
\begin{equation}\label{eq:step-009-active-intermediate}
4n\pi_i<\frac{a_{\rm L}r_i}{8},
\qquad
\left\lceil4n\pi_i\right\rceil<a_{\rm L}r_i,
\qquad
8<a_{\rm L}r_i.
\end{equation}
Consequently
\begin{equation}\label{eq:step-009-active-budget}
m_{n,i}=\max\left\{8,\left\lceil4n\pi_i\right\rceil\right\}
\quad\text{satisfies}\quad
8\le m_{n,i}<a_{\rm L}r_i.
\end{equation}
\end{lemma}

\begin{proof}
Fix \(i\in H\).  From \(r_i>R_0\) and
\eqref{eq:step-009-rzero},
\begin{equation}\label{eq:step-009-active-slack}
a_{\rm L}r_i>a_{\rm L}R_0\ge32.
\end{equation}
In particular,
\begin{equation}\label{eq:step-009-floor-ceiling-slack}
1<\frac{a_{\rm L}r_i}{32},
\qquad
8<\frac{a_{\rm L}r_i}{4}<a_{\rm L}r_i.
\end{equation}
Since \(a_{\rm L}\le1\), \(R_0\ge32\), so
\(r_i>R_0\ge32>2\) and
\begin{equation}\label{eq:step-009-plus-two}
r_i+2<2r_i.
\end{equation}
Using \(\pi_i=s_i/M\), the strict hypothesis, the accepted comparison
\(s_i\le r_i+2\), and \(c_{\rm low}\le a_{\rm L}/64\),
\begin{equation}\label{eq:step-009-active-mass-chain}
\begin{split}
4n\pi_i
&=4\frac nM s_i
<4c_{\rm low}s_i\\
&\le4c_{\rm low}(r_i+2)
\le\frac{a_{\rm L}}{16}(r_i+2)
<\frac{a_{\rm L}r_i}{8}.
\end{split}
\end{equation}
This explicitly absorbs the additive \(+2\).

For every real \(x\), including an integer, \(\lceil x\rceil<x+1\).
Thus \eqref{eq:step-009-active-mass-chain} and the first inequality in
\eqref{eq:step-009-floor-ceiling-slack} give
\begin{equation}\label{eq:step-009-ceiling-absorption}
\begin{split}
\left\lceil4n\pi_i\right\rceil
&<4n\pi_i+1\\
&<\frac{a_{\rm L}r_i}{8}
 +\frac{a_{\rm L}r_i}{32}
=\frac5{32}a_{\rm L}r_i
<a_{\rm L}r_i.
\end{split}
\end{equation}
Equation~\eqref{eq:step-009-floor-ceiling-slack} separately makes the
floor eight strictly subcritical.  Both entries of the maximum in
\eqref{eq:step-009-active-budget} are therefore strictly below
\(a_{\rm L}r_i\), including all ceiling and maximum equality cases.
\end{proof}

\begin{proposition}[Candidate-wise eligibility of every active factor]
\label{prop:step-009-almm-eligibility}
Under Assumption~\ref{assump:candidate-delta-budget},
Proposition~\ref{prop:step-008-factor-floor}, and
Lemma~\ref{lem:step-009-active-budget}, use the constants and set fixed by
Lemma~\ref{lem:step-009-low-mass}.  At the same fixed candidate, assume
\(n<c_{\rm low}M\).  Then every \(i\in H\) satisfies
\begin{equation}\label{eq:step-009-eligibility-budget}
m_{n,i}\in\mathbb N,
\qquad m_{n,i}\ge8,
\qquad m_{n,i}<a_{\rm L}r_i,
\end{equation}
and
\begin{equation}\label{eq:step-009-eligibility-delta}
0<\delta\le
\frac{c_\delta}{m_{n,i}^2\log(m_{n,i}+1)}.
\end{equation}
The other candidate-budget conjunct remains
\begin{equation}\label{eq:step-009-global-candidate-delta}
0<\delta\le\frac1{n\log(n+1)}.
\end{equation}
These are exactly the numerical hypotheses required to apply
Proposition~\ref{prop:step-008-factor-floor} at \(m=m_{n,i}\) once a
total unrestricted replacement-\((0.1,\delta)\)-DP factor learner is
constructed.
\end{proposition}

\begin{proof}
The budget definition gives \(m_{n,i}\in\mathbb N\) and
\(m_{n,i}\ge8\); Lemma~\ref{lem:step-009-active-budget}, whose hypotheses
use this same candidate and the same strict inequality
\(n<c_{\rm low}M\), gives strict subcriticality.

Assumption~\ref{assump:candidate-delta-budget} states at this one candidate
that
\[
0<\delta\le
\min\left\{
\frac1{n\log(n+1)},
\min_{1\le j\le k}
\frac{c_\delta}{m_{n,j}^2\log(m_{n,j}+1)}
\right\}.
\]
The first coordinate of the minimum gives
\eqref{eq:step-009-global-candidate-delta}.  The inner minimum is no
larger than its \(i\)-th entry, giving
\eqref{eq:step-009-eligibility-delta} for every \(i\in H\).  All
denominators are positive because \(n\ge1\) and \(m_{n,i}\ge8\).
Equality at either delta cap is allowed; only the subcritical budget
inequality is strict.  This is a deterministic implication from one finite
minimum, not a probabilistic union bound or a uniform-in-candidate schedule.
\end{proof}

\subsection{Learner-independent hard priors}
\label{app:step-010}

Fix \(i\in H\) and abbreviate \(m=m_{n,i}\).  The constructions below
depend on the fixed class, candidate, factor, and \(\delta\), but are made
before a factor learner is quantified.

\begin{lemma}[Finite quotient-cell coding of the hard tasks]
\label{lem:step-010-finite-coding}
Under Assumption~\ref{assump:countably-coded-evaluation},
Lemma~\ref{lem:step-001-quotient-invariance}, and
Lemma~\ref{lem:step-008-shelah-tasks}, let
\begin{equation}\label{eq:step-010-finite-sets}
F_i=\{x_{i,1},\ldots,x_{i,t_i}\},
\qquad
\mathcal V_i=\{0,1\}^{[t_i]},
\qquad
\Theta_i=[t_i]\times[t_i]^{N_m}.
\end{equation}
There exist:
\begin{enumerate}
\item a measurable retraction
\(\tau_i:(X_i,\Sigma_i)\to(F_i,2^{F_i})\) fixing every \(x_{i,j}\);
\item for each \(v\in\mathcal V_i\), a measurable, possibly improper
hypothesis \(e_i(v):X_i\to\{0,1\}\) with
\(e_i(v)(x_{i,j})=v_j\), such that \(v\mapsto e_i(v)\) has measurable
finite-evaluation cylinders; and
\item for every \(\theta=(\ell,\mathbf j)\in\Theta_i\), an exact finite
labeled law \(P_\theta\) on \(F_i\times\{0,1\}\) and a risk vector
\(r_\theta:\mathcal V_i\to[0,1]\) such that, for
\begin{equation}\label{eq:step-010-indexed-task}
T_\theta=(c_\theta,D_\theta)
=\left(c_{i,\ell},\frac1{N_m}\sum_{a=1}^{N_m}
\delta_{x_{i,j_a}}\right),
\end{equation}
one has
\begin{equation}\label{eq:step-010-exact-coded-risk}
R_{D_\theta}(e_i(v),c_\theta)=r_\theta(v)
\end{equation}
for every \(v\), and the iid labeled sample law is the finite mixture
indexed by \(P_\theta^m\).
\end{enumerate}
All these sets are finite and nonempty, including when \(t_i=1\), and no
measurable section of \(\kappa_i\) is used.
\end{lemma}

\begin{proof}
Put
\begin{equation}\label{eq:step-010-witness-cells}
G_{i,j}=\kappa_i^{-1}(\{q_{i,j}\})\in\Sigma_i.
\end{equation}
The quotient points are distinct: if \(j<j'\) but
\(q_{i,j}=q_{i,j'}\), the threshold pattern with \(\ell=j+1\) would give
labels zero and one at the same quotient point.  Hence the cells and their
chosen representatives are pairwise distinct.

Define
\begin{equation}\label{eq:step-010-retraction}
\tau_i(x)=
\begin{cases}
x_{i,j},&x\in G_{i,j}\text{ for some }j\in[t_i],\\
x_{i,1},&x\notin\bigcup_{j=1}^{t_i}G_{i,j}.
\end{cases}
\end{equation}
For \(j\ge2\), the fiber of \(x_{i,j}\) is \(G_{i,j}\); the fiber of
\(x_{i,1}\) is \(X_i\setminus\bigcup_{j=2}^{t_i}G_{i,j}\).  All fibers
are measurable, so \(\tau_i\) is measurable into the finite discrete
space and fixes each witness point.  When \(t_i=1\), it is simply the
constant map.

For \(v=(v_1,\ldots,v_{t_i})\), define
\begin{equation}\label{eq:step-010-output-extension}
e_i(v)(x)=
\begin{cases}
v_j,&x\in G_{i,j}\text{ for some }j\in[t_i],\\
0,&x\notin\bigcup_{j=1}^{t_i}G_{i,j}.
\end{cases}
\end{equation}
The inverse image of one is a finite union of measurable cells.  Since the
output space is finite discrete, for arbitrary raw points \(z_1,\ldots,z_s\)
and bits \(b_1,\ldots,b_s\), the set of \(v\)'s satisfying
\(e_i(v)(z_a)=b_a\) for all \(a\) is measurable.  The extension is allowed
to be improper.

For \(\theta=(\ell,(j_1,\ldots,j_{N_m}))\), put
\begin{equation}\label{eq:step-010-row-weights}
w_{\theta,j}=\frac1{N_m}
\bigl|\{a\in[N_m]:j_a=j\}\bigr|
\end{equation}
and define
\begin{equation}\label{eq:step-010-finite-training-law}
P_\theta(x_{i,j},b)
=w_{\theta,j}\mathbf1\{b=c_{i,\ell}(x_{i,j})\}.
\end{equation}
These nonnegative weights sum to one.  Define
\begin{equation}\label{eq:step-010-risk-vector}
\begin{split}
r_\theta(v)
&=\sum_{j=1}^{t_i}w_{\theta,j}
\mathbf1\{v_j\ne c_{i,\ell}(x_{i,j})\}\\
&=\frac1{N_m}\sum_{a=1}^{N_m}
\mathbf1\{v_{j_a}\ne c_{i,\ell}(x_{i,j_a})\}.
\end{split}
\end{equation}
Equations~\eqref{eq:step-010-indexed-task},
\eqref{eq:step-010-output-extension}, and
\eqref{eq:step-010-risk-vector} prove
\eqref{eq:step-010-exact-coded-risk}; the iid labeled law is the mixture
over \((F_i\times\{0,1\})^m\) with weights \(P_\theta^m\).

This finite representation does not assert measurability of raw
singletons.  Each \(D_\theta\) is a finite convex combination of Dirac
probability measures, valid on any measurable space, while the measurable
quotient cells are what make \(\tau_i\) and \(e_i\) legal.
\end{proof}

\begin{proposition}[Exact finite representation of unrestricted private learners]
\label{prop:step-010-kernel-equivalence}
Under Assumption~\ref{assump:countably-coded-evaluation} and
Lemma~\ref{lem:step-010-finite-coding}, put
\begin{equation}\label{eq:step-010-input-alphabet}
\mathcal Z_i=F_i\times\{0,1\},
\qquad
\mathcal U_i=\mathcal Z_i^m,
\end{equation}
and let \(\mathsf{Adj}_i^{\to}\) contain every ordered pair of elements of
\(\mathcal U_i\) differing in exactly one row.  Let \(\mathcal P_i\) be
the arrays \(p=(p_{u,v})\) satisfying
\begin{equation}\label{eq:step-010-stochasticity}
p_{u,v}\ge0,
\qquad
\sum_{v\in\mathcal V_i}p_{u,v}=1
\quad(u\in\mathcal U_i),
\end{equation}
and, for every \((u,u')\in\mathsf{Adj}_i^{\to}\) and every
\(E\subseteq\mathcal V_i\),
\begin{equation}\label{eq:step-010-finite-dp}
\sum_{v\in E}p_{u,v}
\le e^{0.1}\sum_{v\in E}p_{u',v}+\delta.
\end{equation}
For \(\theta\in\Theta_i\), define
\begin{equation}\label{eq:step-010-payoff}
L_\theta(p)
=\sum_{u\in\mathcal U_i}P_\theta^m(u)
\sum_{v\in\mathcal V_i}p_{u,v}r_\theta(v).
\end{equation}
Then every unrestricted total raw-space factor learner on \(m\) rows with
measurable finite evaluations and all-input replacement-
\((0.1,\delta)\)-DP induces \(p^A\in\mathcal P_i\) whose exact expected
risk on \(T_\theta\) is \(L_\theta(p^A)\).  Conversely, every
\(p\in\mathcal P_i\) extends to such a total raw-space learner \(A^p\),
with finite output \(\mathcal V_i\), decoder \(e_i\), and exact expected
risk \(L_\theta(p)\) on every \(T_\theta\).
\end{proposition}

\begin{proof}
Let \(A:(X_i\times\{0,1\})^m\rightsquigarrow(\Omega,\mathscr F)\) be
an arbitrary learner of the stated type, with decoder \(h_\omega\).  The
finite evaluation map
\begin{equation}\label{eq:step-010-evaluation-map}
\operatorname{ev}_i(\omega)
=(h_\omega(x_{i,1}),\ldots,h_\omega(x_{i,t_i}))
\in\mathcal V_i
\end{equation}
is measurable: its singleton fibers are finite-evaluation cylinders and
the inverse image of any \(E\subseteq\mathcal V_i\) is a finite union of
them.  For \(u\in\mathcal U_i\), viewed as a raw labeled input, set
\begin{equation}\label{eq:step-010-restricted-kernel}
p^A_{u,v}=A(u,\operatorname{ev}_i^{-1}(\{v\})).
\end{equation}
The fibers partition \(\Omega\), so
\eqref{eq:step-010-stochasticity} holds.  For any ordered adjacent
\((u,u')\), whether realizable or not, apply raw all-input privacy directly
to \(\operatorname{ev}_i^{-1}(E)\) to obtain
\eqref{eq:step-010-finite-dp}.  Because ordered adjacency includes the
swapped pair, both directions are retained.  Requiring all events is
essential: summing atomwise approximate-DP inequalities would multiply
\(\delta\).

Expanding the finite sample law and grouping outputs by
\eqref{eq:step-010-evaluation-map},
\begin{equation}\label{eq:step-010-restriction-payoff}
\begin{split}
&\mathbb E_{S\sim(D_\theta)_{c_\theta}^m,A}
R_{D_\theta}(h_{A(S)},c_\theta)\\
&\qquad=\sum_{u\in\mathcal U_i}P_\theta^m(u)
\sum_{v\in\mathcal V_i}p^A_{u,v}r_\theta(v)
=L_\theta(p^A).
\end{split}
\end{equation}
This is exact for arbitrary improper output because the task risk depends
only on the witness evaluations.

Conversely, fix \(p\in\mathcal P_i\).  Extend the retraction to labeled
records by
\begin{equation}\label{eq:step-010-labeled-retraction}
\widehat\tau_i(x,y)=(\tau_i(x),y)
\end{equation}
and coordinatewise to databases.  Define
\begin{equation}\label{eq:step-010-extended-kernel}
A^p(S,E)=\sum_{v\in E}p_{\widehat\tau_i^m(S),v},
\qquad E\subseteq\mathcal V_i,
\end{equation}
and decode \(v\) as \(e_i(v)\).  Measurability of
\(\widehat\tau_i^m\) into a finite discrete space makes this a Markov
kernel; row stochasticity makes it total on every raw input, including
inconsistent and nonrealizable inputs.

For adjacent raw datasets, their preprocessed images are either equal or
differ in the same one coordinate.  Equal images give identical output
laws; otherwise \eqref{eq:step-010-finite-dp} applies to every output
event, and the swapped ordered pair supplies the reverse direction.  Thus
there is no composition or added privacy cost.  On a task draw, every raw
record is a chosen representative fixed by \(\tau_i\), so preprocessing is
the identity and
\begin{equation}\label{eq:step-010-extension-payoff}
\mathbb E_{S\sim(D_\theta)_{c_\theta}^m,A^p}
R_{D_\theta}(e_i(A^p(S)),c_\theta)
=L_\theta(p).
\end{equation}
This proves both directions with exact privacy and risk.
\end{proof}

\begin{lemma}[Compact finite private-learning game]
\label{lem:step-010-compact-game}
Under Proposition~\ref{prop:step-010-kernel-equivalence},
\(\mathcal P_i\) is a nonempty compact convex polytope in
\(\mathbb R^{\mathcal U_i\times\mathcal V_i}\).  It has a nonempty finite
vertex set \(\mathcal W_i\) and
\begin{equation}\label{eq:step-010-vertex-hull}
\mathcal P_i=\operatorname{conv}(\mathcal W_i).
\end{equation}
Every \(L_\theta\) is affine and continuous, and
\begin{equation}\label{eq:step-010-worst-task-payoff}
g_i(p)=\max_{\theta\in\Theta_i}L_\theta(p)
\end{equation}
is continuous and attains its minimum on \(\mathcal P_i\).
\end{lemma}

\begin{proof}
The accepted construction and eligibility give \(t_i\ge1\) and \(m\ge8\),
and therefore
\begin{equation}\label{eq:step-010-cardinalities}
|\mathcal V_i|=2^{t_i}<\infty,
\qquad
|\mathcal U_i|=(2t_i)^m<\infty,
\qquad
|\Theta_i|=t_i^{N_m+1}<\infty.
\end{equation}
There are finitely many ordered adjacent pairs and finitely many subsets of
\(\mathcal V_i\).  Thus
\eqref{eq:step-010-stochasticity}--\eqref{eq:step-010-finite-dp} form a
finite system of affine equalities and closed affine inequalities.  The
set is closed and convex, and stochasticity gives
\(0\le p_{u,v}\le1\), hence boundedness.

It is nonempty: fix \(v^0\in\mathcal V_i\) and put
\begin{equation}\label{eq:step-010-constant-kernel}
p^0_{u,v}=\mathbf1\{v=v^0\}
\quad\text{for every }u.
\end{equation}
For an event containing \(v^0\), privacy reads
\(1\le e^{0.1}+\delta\); otherwise both sides are zero.  Thus the
constant kernel is feasible.  A nonempty closed bounded set in this
finite-dimensional space is compact.  The finite polytope vertex theorem
then gives a finite nonempty extreme-point set and
\eqref{eq:step-010-vertex-hull}.

Equation~\eqref{eq:step-010-payoff} is a finite linear combination of the
coordinates \(p_{u,v}\), so every payoff is affine, continuous, and lies
in \([0,1]\).  The maximum in
\eqref{eq:step-010-worst-task-payoff} is over a finite nonempty set, hence
continuous, and compactness gives a minimizer.  No property of an infinite
learner or output space is used for this attainment.
\end{proof}

\begin{proposition}[Strict compact-game value]
\label{prop:step-010-strict-value}
Under Propositions~\ref{prop:step-008-factor-floor} and
\ref{prop:step-009-almm-eligibility},
Proposition~\ref{prop:step-010-kernel-equivalence}, and
Lemma~\ref{lem:step-010-compact-game}, for every fixed \(i\in H\),
\begin{equation}\label{eq:step-010-strict-value}
\gamma_i
=\min_{p\in\mathcal P_i}\max_{\theta\in\Theta_i}L_\theta(p)
>\frac18.
\end{equation}
Thus \(\eta_i=\gamma_i-1/8\) is a well-defined positive finite-game
margin, not an additional assumption or universal numerical constant.
\end{proposition}

\begin{proof}
Fix \(p\in\mathcal P_i\).  Proposition~\ref{prop:step-010-kernel-equivalence}
extends it to a total learner \(A^p\) defined on all labeled inputs, with
all-event replacement-\((0.1,\delta)\)-privacy, measurable possibly
improper output, and exact task loss \(L_\theta(p)\).  At this same factor
and candidate, Proposition~\ref{prop:step-009-almm-eligibility} supplies
\begin{equation}\label{eq:step-010-floor-eligibility}
m\ge8,
\qquad m<a_{\rm L}\log_2^*d_i,
\qquad
0<\delta\le\frac{c_\delta}{m^2\log(m+1)}.
\end{equation}
Proposition~\ref{prop:step-008-factor-floor} therefore gives an indexed
task with expected risk strictly greater than \(1/8\).  By exact payoff
preservation,
\begin{equation}\label{eq:step-010-pointwise-strict}
g_i(p)=\max_{\theta\in\Theta_i}L_\theta(p)>\frac18
\quad\text{for every }p\in\mathcal P_i.
\end{equation}
Lemma~\ref{lem:step-010-compact-game} supplies an actual minimizer
\(p_i^*\).  Evaluating \eqref{eq:step-010-pointwise-strict} at that point
gives
\[
\gamma_i=g_i(p_i^*)>\frac18.
\]
Attainment is essential: taking only an infimum of pointwise strict
inequalities would yield merely a non-strict bound.
\end{proof}

\begin{proposition}[Finite learner-independent hard prior]
\label{prop:step-010-hard-prior}
Under Proposition~\ref{prop:step-010-strict-value} and
Lemma~\ref{lem:step-010-compact-game}, for every fixed \(i\in H\) there is
a probability distribution \(\mu_i\) on \(\Theta_i\) such that
\begin{equation}\label{eq:step-010-indexed-prior}
\sum_{\theta\in\Theta_i}\mu_i(\theta)L_\theta(p)
\ge\gamma_i>\frac18
\quad\text{for every }p\in\mathcal P_i.
\end{equation}
Its pushforward \(\nu_i\) under \(\theta\mapsto T_\theta\) is a finite
prior over realizable finite-support factor tasks, fixed independently of
every learner.  Every unrestricted total replacement-\((0.1,\delta)\)-DP
factor learner \(A\) on \(m_{n,i}\) rows with measurable finite evaluations
satisfies
\begin{equation}\label{eq:step-010-hard-prior}
\mathbb E_{(c_i,D_i)\sim\nu_i}
\mathbb E_{S_i\sim(D_i)_{c_i}^{m_{n,i}},A}
R_{D_i}(h_{A(S_i)},c_i)>\frac18.
\end{equation}
\end{proposition}

\begin{proof}
Write the finite vertex set as
\(\mathcal W_i=\{p^1,\ldots,p^J\}\).  Every
\(p\in\mathcal P_i\) has a representation
\(p=\sum_{a=1}^J\lambda_ap^a\), \(\lambda\in\Delta([J])\), and affinity
gives
\begin{equation}\label{eq:step-010-affine-vertex-payoff}
L_\theta(p)=\sum_{a=1}^J\lambda_aL_\theta(p^a).
\end{equation}
Apply the finite von Neumann minimax theorem to the real matrix
\begin{equation}\label{eq:step-010-game-matrix}
M_{a,\theta}=L_\theta(p^a),
\qquad(a,\theta)\in[J]\times\Theta_i.
\end{equation}
Using \eqref{eq:step-010-vertex-hull},
\eqref{eq:step-010-affine-vertex-payoff}, and the fact that a linear
functional on the task simplex is maximized by a pure task,
\begin{equation}\label{eq:step-010-minimax}
\begin{split}
\gamma_i
&=\min_{\lambda\in\Delta([J])}
  \max_{\mu\in\Delta(\Theta_i)}
  \sum_{a,\theta}\lambda_a\mu(\theta)M_{a,\theta}\\
&=\max_{\mu\in\Delta(\Theta_i)}
  \min_{\lambda\in\Delta([J])}
  \sum_{a,\theta}\lambda_a\mu(\theta)M_{a,\theta}.
\end{split}
\end{equation}
The finite theorem also gives an attaining maximizing mixture \(\mu_i\).
For any \(p\), choose a representing \(\lambda\); the last line of
\eqref{eq:step-010-minimax} gives
\begin{equation}\label{eq:step-010-minimax-guarantee}
\sum_\theta\mu_i(\theta)L_\theta(p)
=\sum_{a,\theta}\lambda_a\mu_i(\theta)M_{a,\theta}
\ge\gamma_i.
\end{equation}

Every indexed task is realizable and finitely supported.  If two indices
encode the same task, pushforward simply adds their masses and leaves all
expectations unchanged.  Finally, any raw learner in the proposition maps
by Proposition~\ref{prop:step-010-kernel-equivalence} to
\(p^A\in\mathcal P_i\).  Combining
\eqref{eq:step-010-minimax-guarantee},
\eqref{eq:step-010-restriction-payoff}, and
Proposition~\ref{prop:step-010-strict-value} gives
\[
\mathbb E_{(c_i,D_i)\sim\nu_i}
\mathbb E_{S_i,A}R_{D_i}(h_{A(S_i)},c_i)
\ge\gamma_i>\frac18.
\]
The finite game, its optimizer, and \(\nu_i\) are chosen before \(A\),
which is the required prior-before-learner order.
\end{proof}

\subsection{Uniform control of factor-sample overflow}
\label{app:step-011}

\begin{lemma}[Optimized binomial exponential envelope]
\label{lem:step-011-binomial-envelope}
Let \(n\in\mathbb N\), \(0<p\le1\),
\(L\sim\operatorname{Bin}(n,p)\), and \(\mu=np\).  For every real
\(t>\mu\),
\begin{equation}\label{eq:step-011-binomial-envelope}
\Pr(L\ge t)
\le B(\mu,t)
:=\exp\!\left\{t-\mu-t\log\frac{t}{\mu}\right\}
=e^{-\mu}\left(\frac{e\mu}{t}\right)^t.
\end{equation}
For fixed \(t>0\), \(B(\mu,t)\) is strictly increasing on
\(0<\mu<t\); for fixed \(\mu>0\), it is strictly decreasing on
\(t>\mu\).  If \(t>n\), the probability on the left is zero.
\end{lemma}

\begin{proof}
For \(\lambda>0\), exponential Markov and the exact binomial MGF give
\begin{equation}\label{eq:step-011-exponential-markov}
\Pr(L\ge t)
\le e^{-\lambda t}\mathbb Ee^{\lambda L},
\end{equation}
and
\begin{equation}\label{eq:step-011-binomial-mgf}
\mathbb Ee^{\lambda L}
=(1-p+pe^\lambda)^n
=\bigl(1+p(e^\lambda-1)\bigr)^n
\le\exp\{\mu(e^\lambda-1)\}.
\end{equation}
Here \(1+u\le e^u\) for \(u\ge0\), since
\(e^u-1-u\) vanishes at zero and has nonnegative derivative
\(e^u-1\).  Because \(t>\mu>0\), choose
\begin{equation}\label{eq:step-011-optimal-lambda}
\lambda=\log(t/\mu)>0.
\end{equation}
Substitution gives exponent
\(-t\log(t/\mu)+\mu(t/\mu-1)=t-\mu-t\log(t/\mu)\), proving
\eqref{eq:step-011-binomial-envelope}.  Finite support handles \(t>n\).

Finally,
\begin{equation}\label{eq:step-011-mu-derivative}
\frac{\partial}{\partial\mu}\log B(\mu,t)
=-1+\frac t\mu>0
\quad(0<\mu<t),
\end{equation}
and
\begin{equation}\label{eq:step-011-t-derivative}
\frac{\partial}{\partial t}\log B(\mu,t)
=\log\frac\mu t<0
\quad(t>\mu).
\end{equation}
Since \(B>0\), these signs prove the monotonicity claims.
\end{proof}

\begin{lemma}[Exact numerical slack for overflow]
\label{lem:step-011-numerical-slack}
With \(e\) the base of the natural logarithm,
\begin{equation}\label{eq:step-011-eta-zero}
\eta_0=e^7\left(\frac29\right)^9
<\frac3{2048}.
\end{equation}
Moreover, \(\log4>1\).
\end{lemma}

\begin{proof}
For \(\ell\in\mathbb N_0\),
\begin{equation}\label{eq:step-011-factorial-tail}
(6+\ell)!\ge6!\,7^\ell.
\end{equation}
The exponential and geometric series therefore give
\begin{equation}\label{eq:step-011-e-rational-bound}
\begin{split}
e
&=\sum_{j=0}^5\frac1{j!}
 +\sum_{\ell=0}^\infty\frac1{(6+\ell)!}\\
&\le\frac{163}{60}+\frac1{720}\sum_{\ell=0}^\infty7^{-\ell}
=\frac{163}{60}+\frac7{4320}
=\frac{11743}{4320}<\frac{87}{32}.
\end{split}
\end{equation}
The final comparison is exact because
\begin{equation}\label{eq:step-011-e-cross-product}
11743\cdot32=375776<375840=87\cdot4320.
\end{equation}
Thus \(e<87/32<4\), and strict monotonicity of the natural logarithm gives
\(1=\log e<\log4\).

Using \(87=3\cdot29\) and \(9=3^2\),
\begin{equation}\label{eq:step-011-eta-rational}
\eta_0
<\left(\frac{87}{32}\right)^7\left(\frac29\right)^9
=\frac{29^7}{2^{26}3^{11}}.
\end{equation}
Exact integer arithmetic gives
\begin{equation}\label{eq:step-011-integer-slack}
29^7=17{,}249{,}876{,}309
<17{,}414{,}258{,}688=3^{12}2^{15}.
\end{equation}
Substitution into \eqref{eq:step-011-eta-rational} yields
\(\eta_0<3^{12}2^{15}/(2^{26}3^{11})=3/2^{11}\), as claimed.
\end{proof}

\begin{lemma}[Floor-eight overflow at small mean]
\label{lem:step-011-small-mean}
Under Lemma~\ref{lem:step-011-binomial-envelope}, let
\(L\sim\operatorname{Bin}(n,p)\), \(0<p\le1\),
\(0<\mu=np\le2\), and
\(m=\max\{8,\lceil4\mu\rceil\}\).  Then
\begin{equation}\label{eq:step-011-small-mean-bound}
\Pr(L>m)\le e^7(2/9)^9=\eta_0.
\end{equation}
This includes \(\mu=2\).
\end{lemma}

\begin{proof}
If \(n\le8\), then \(L\le n\le8\le m\) pointwise.  Otherwise
\(n\ge9\), and integer-valued \(L,m\) with \(m\ge8\) give
\begin{equation}\label{eq:step-011-small-event}
\{L>m\}=\{L\ge m+1\}\subseteq\{L\ge9\}.
\end{equation}
Since \(0<\mu\le2<9\), Lemma~\ref{lem:step-011-binomial-envelope} at
\(t=9\), followed by its fixed-threshold monotonicity, yields
\begin{equation}\label{eq:step-011-small-envelope}
\Pr(L>m)\le B(\mu,9)\le B(2,9)
=e^{-2}\left(\frac{2e}{9}\right)^9
=e^7\left(\frac29\right)^9.
\end{equation}
At \(\mu=2\), this is exactly the transition endpoint.
\end{proof}

\begin{lemma}[Factor-four overflow at large mean]
\label{lem:step-011-large-mean}
Under Lemmas~\ref{lem:step-011-binomial-envelope} and
\ref{lem:step-011-numerical-slack}, let
\(L\sim\operatorname{Bin}(n,p)\), \(0<p\le1\), \(\mu=np\ge2\), and
\(m=\max\{8,\lceil4\mu\rceil\}\).  Then
\begin{equation}\label{eq:step-011-large-mean-bound}
\Pr(L>m)\le e^7(2/9)^9=\eta_0.
\end{equation}
\end{lemma}

\begin{proof}
If \(m\ge n\), finite support gives zero probability.  Otherwise, because
\(L,m\) are integers and \(m\ge\lceil4\mu\rceil\),
\begin{equation}\label{eq:step-011-large-event}
\{L>m\}=\{L\ge m+1\}
\subseteq\{L\ge4\mu+1\},
\end{equation}
where
\begin{equation}\label{eq:step-011-ceiling-successor}
m+1\ge\lceil4\mu\rceil+1\ge4\mu+1.
\end{equation}
The real threshold \(4\mu+1\) exceeds \(\mu\), so
\begin{equation}\label{eq:step-011-large-envelope}
\Pr(L>m)\le B(\mu,4\mu+1).
\end{equation}

Let
\begin{equation}\label{eq:step-011-log-envelope}
f(\mu)=\log B(\mu,4\mu+1)
=3\mu+1-(4\mu+1)\log\left(4+\frac1\mu\right).
\end{equation}
Direct differentiation gives
\begin{equation}\label{eq:step-011-large-derivative}
f'(\mu)=3+\frac1\mu
-4\log\left(4+\frac1\mu\right).
\end{equation}
For \(\mu\ge2\), \(1/\mu\le1/2\),
\(4+1/\mu>4\), and Lemma~\ref{lem:step-011-numerical-slack} gives
\(\log4>1\).  Therefore
\begin{equation}\label{eq:step-011-negative-derivative}
f'(\mu)<3+\frac12-4\log4
<\frac72-4=-\frac12<0.
\end{equation}
Thus \(f\) decreases on \([2,\infty)\), and
\begin{equation}\label{eq:step-011-large-endpoint}
B(\mu,4\mu+1)\le B(2,9)=e^7(2/9)^9.
\end{equation}
Equations~\eqref{eq:step-011-large-envelope} and
\eqref{eq:step-011-large-endpoint} prove the claim, with exact agreement
at \(\mu=2\).
\end{proof}

\begin{lemma}[Finite-support endpoints and one-factor zero overflow]
\label{lem:step-011-zero-overflow}
Let \(L\sim\operatorname{Bin}(n,p)\), \(p\in[0,1]\), and
\(m=\max\{8,\lceil4np\rceil\}\).  Then
\begin{equation}\label{eq:step-011-zero-overflow}
\Pr(L>m)=0
\end{equation}
whenever \(n\le8\), \(p=0\), or \(p=1\).  Under
Assumptions~\ref{assump:canonical-product} and
\ref{assump:vc-one-factors} and the exact weights of
Lemma~\ref{lem:step-009-low-mass}, this applies with \(p=\pi_i\), and in
particular gives zero overflow when \(k=1\).
\end{lemma}

\begin{proof}
Binomial support gives \(L\le n\).  Thus \(n\le8\) implies
\(L\le8\le m\).  If \(p=0\), then \(L=0\) almost surely.  If \(p=1\),
then \(L=n\) almost surely and
\begin{equation}\label{eq:step-011-p-one-budget}
m=\max\{8,\lceil4n\rceil\}=\max\{8,4n\}\ge n.
\end{equation}
When \(k=1\), \(M=s_1\), so \(\pi_1=s_1/M=1\) and the last case gives
\(L_1=n\le m_{n,1}\) almost surely.
\end{proof}

\begin{proposition}[Uniform marginal overflow certificate]
\label{prop:step-011-overflow-certificate}
Under Assumptions~\ref{assump:canonical-product} and
\ref{assump:vc-one-factors} and Lemma~\ref{lem:step-009-low-mass}, for the
same fixed candidate and every factor, if
\(L_i\sim\operatorname{Bin}(n,\pi_i)\), then
\begin{equation}\label{eq:step-011-final-overflow}
\Pr[L_i>m_{n,i}]
\le\eta_0=e^7\left(\frac29\right)^9
<\frac3{2048}.
\end{equation}
The probability is zero when \(\pi_i=0\), \(\pi_i=1\), or \(n\le8\),
and hence is zero in the one-factor case.
\end{proposition}

\begin{proof}
Lemma~\ref{lem:step-011-zero-overflow} handles the stated exact endpoints.
Otherwise put \(\mu=n\pi_i\).  Lemma~\ref{lem:step-011-small-mean}
handles \(0<\mu\le2\), while Lemma~\ref{lem:step-011-large-mean} handles
\(\mu>2\); they agree at \(\mu=2\).  Lemma~\ref{lem:step-011-numerical-slack}
supplies the strict comparison with \(3/2048\).  This is a uniform
factorwise marginal statement, not a union bound or a claim of independence
among counts.
\end{proof}

\subsection{One-use hidden-factor simulation}
\label{app:step-012}

All active priors \(\nu_j\) are fixed before the arbitrary global learner
is quantified.  A hidden learner may depend on that already fixed global
learner, but none of its hard priors does.

\begin{lemma}[Prior-fixed common ideal mixture experiment]
\label{lem:step-012-ideal-experiment}
Under Assumption~\ref{assump:canonical-product},
Proposition~\ref{prop:step-009-almm-eligibility}, and
Proposition~\ref{prop:step-010-hard-prior}, one can fix, before every global
learner, a law \(\lambda_j\) on realizable factor tasks for each \(j\),
with \(\lambda_j=\nu_j\) on \(H\) and \(\lambda_j\) a deterministic point
mass off \(H\).  Independent tasks from these laws determine a measurable
full-product target \(c^\star\in C\), a probability measure \(D^\star\)
with \(D^\star(X_j)=\pi_j\), and an ideal sample
\(S^\star\sim(D^\star_{c^\star})^n\).  Its route counts satisfy
\(L_j\sim\operatorname{Bin}(n,\pi_j)\), and every task is chosen before
the single call to a fixed global kernel.
\end{lemma}

\begin{proof}
For each \(j\notin H\), fix \(c_j^\circ\in C_j\) and
\(x_j^\circ\in X_j\), and set
\begin{equation}\label{eq:step-012-low-reference-task}
T_j^\circ=(c_j^\circ,\delta_{x_j^\circ}),
\qquad
\lambda_j=\delta_{T_j^\circ}.
\end{equation}
Only finitely many choices are made.  For \(j\in H\), set
\(\lambda_j=\nu_j\).  Proposition~\ref{prop:step-010-hard-prior} makes
each active law finitely supported on realizable finite-support tasks and
fixed before the learner.

Draw independently
\begin{equation}\label{eq:step-012-task-vector}
T_j=(c_j,D_j)\sim\lambda_j,
\qquad j\in[k].
\end{equation}
Assumption~\ref{assump:canonical-product} supplies the unique
\(c^\star\in C\) restricting to \(c_j\) on \(X_j\).  Define
\begin{equation}\label{eq:step-012-mixture-distribution}
D^\star(B)=\sum_{j=1}^k\pi_jD_j(B\cap X_j),
\qquad B\in\Sigma.
\end{equation}
The measurable disjoint blocks and weights summing to one make this a
probability measure with \(D^\star(X_j)=\pi_j\).

Draw iid route indices
\begin{equation}\label{eq:step-012-route}
I_1,\ldots,I_n\in[k],
\qquad\Pr(I_r=j)=\pi_j,
\end{equation}
and set
\begin{equation}\label{eq:step-012-route-counts}
\ell_{j,r}=|\{s\le r:I_s=j\}|,
\qquad L_j=\ell_{j,n}.
\end{equation}
Conditional on the task vector and route, draw independently over all
\(j\) and \(a\le n\vee m_{n,j}\),
\begin{equation}\label{eq:step-012-row-arrays}
U_{j,a}=(X_{j,a},c_j(X_{j,a})),
\qquad X_{j,a}\sim D_j,
\end{equation}
and put
\begin{equation}\label{eq:step-012-ideal-sample}
S_r^\star=U_{I_r,\ell_{I_r,r}},
\qquad S^\star=(S_1^\star,\ldots,S_n^\star).
\end{equation}
Distinct global slots use distinct independent row-array coordinates.  For
measurable labeled-record events \(B_1,\ldots,B_n\), conditional on the
task vector,
\begin{equation}\label{eq:step-012-iid-mixture-calculation}
\begin{split}
&\Pr(S_1^\star\in B_1,\ldots,S_n^\star\in B_n\mid(T_j)_j)\\
&\quad=\sum_{j_1,\ldots,j_n}
\prod_{r=1}^n\pi_{j_r}(D_{j_r})_{c_{j_r}}(B_r)\\
&\quad=\prod_{r=1}^n\sum_{j=1}^k
\pi_j(D_j)_{c_j}(B_r)
=\prod_{r=1}^nD^\star_{c^\star}(B_r).
\end{split}
\end{equation}
Thus the sample is iid from the displayed mixture, and
\eqref{eq:step-012-route} gives the binomial marginal counts.

All task, support-point, and route indices range over finite sets, so the
experiment is a finite mixture of Dirac laws and is measurable without a
sigma-field on an infinite task family.  After the tasks, route, and arrays
are drawn and the sample assembled, make one call
\begin{equation}\label{eq:step-012-ideal-output}
\Omega^\star\sim A_n(S^\star,\cdot).
\end{equation}
\end{proof}

\begin{proposition}[Total measurable one-use hidden-factor kernel]
\label{prop:step-012-hidden-kernel}
Under Assumption~\ref{assump:countably-coded-evaluation},
Lemma~\ref{lem:step-010-finite-coding},
Proposition~\ref{prop:step-010-hard-prior}, and
Lemma~\ref{lem:step-012-ideal-experiment}, fix the task laws, then an
arbitrary total global Markov kernel
\(A_n:Z^n\rightsquigarrow(\Omega,\mathscr F)\) with measurable finite
evaluations, and then \(i\in H\).  There is a total Markov kernel
\begin{equation}\label{eq:step-012-hidden-kernel}
K_i^{A_n}:
(X_i\times\{0,1\})^{m_{n,i}}
\rightsquigarrow(\mathcal V_i,2^{\mathcal V_i})
\end{equation}
with decoder \(e_i\), defined on every labeled factor database.  Its
overflow branch is input-independent and occurs before any input row is
read; off overflow, input row \(a\) is inserted only into the \(a\)-th
requested \(i\)-slot.  All other tasks and rows are sampled independently
of the input before the one call to \(A_n\).
\end{proposition}

\begin{proof}
Write \(m_i=m_{n,i}\) and define the finite-evaluation restriction
\begin{equation}\label{eq:step-012-output-restriction}
\rho_i:\Omega\to\mathcal V_i,
\qquad
\rho_i(\omega)
=(h_\omega(x_{i,1}),\ldots,h_\omega(x_{i,t_i})).
\end{equation}
For every \(E\subseteq\mathcal V_i\), the preimage
\(\rho_i^{-1}(E)\) is a finite union of finite-evaluation cylinders and is
measurable.  Fix a fallback \(v_i^\circ\in\mathcal V_i\).

On arbitrary input
\begin{equation}\label{eq:step-012-factor-input}
s=(z_1,\ldots,z_{m_i})
\in(X_i\times\{0,1\})^{m_i},
\end{equation}
the kernel samples without inspecting \(s\): all tasks \(T_j\sim\lambda_j\)
for \(j\ne i\), the entire route and hence
\(\mathcal O_i=\{L_i>m_i\}\), and, only on \(\mathcal O_i^c\), the rows
needed at non-\(i\) slots.  On overflow it immediately returns
\(v_i^\circ\), without reading \(s\), requesting an unavailable row,
assembling a database, or calling \(A_n\).

On \(\mathcal O_i^c\), for a route slot with \(I_r=i\), the occurrence
number \(\ell_{i,r}\) lies in \([m_i]\).  If \(\xi\) denotes all realized
input-independent auxiliary variables, define
\begin{equation}\label{eq:step-012-assembly-map}
\bigl(\Phi_{i,\xi}(s)\bigr)_r=
\begin{cases}
z_{\ell_{i,r}},&I_r=i,\\
W_r,&I_r\ne i,
\end{cases}
\end{equation}
where \(W_r\sim(D_{I_r})_{c_{I_r}}\) was already drawn.  Each requested
input row is used exactly once, and every row \(z_a\) with \(a>L_i\) is
unused.  The kernel calls \(A_n\) once on this database and returns
\(\rho_i(\omega)\).

Let \(\mathsf Q_i\) be the finite law of \(\xi\).  For
\(E\subseteq\mathcal V_i\), define, without evaluating
\(\Phi_{i,\xi}\) on overflow,
\begin{equation}\label{eq:step-012-conditional-kernel}
K_{i,\xi}(s,E)=
\begin{cases}
\mathbf1_E(v_i^\circ),&\mathcal O_i(\xi),\\
A_n(\Phi_{i,\xi}(s),\rho_i^{-1}(E)),&\mathcal O_i(\xi)^c,
\end{cases}
\end{equation}
and
\begin{equation}\label{eq:step-012-kernel-mixture}
K_i^{A_n}(s,E)=\int K_{i,\xi}(s,E)\,\mathsf Q_i(d\xi).
\end{equation}
For fixed nonoverflow \(\xi\), every coordinate of \(\Phi_{i,\xi}\) is
an input projection or a constant labeled record.  The trace-space
inclusion is measurable because \(X_i\in\Sigma\); hence the assembly map
and the nonoverflow branch are measurable.  The overflow branch is
constant, and the integral is a finite sum.  For each input it is a
probability law.  It is total on inconsistent and nonrealizable inputs,
and Lemma~\ref{lem:step-010-finite-coding} supplies the measurable
possibly improper decoder \(e_i\).
\end{proof}

\begin{proposition}[All-input privacy from one-use insertion]
\label{prop:step-012-hidden-privacy}
Under Assumption~\ref{assump:global-privacy-range} and
Proposition~\ref{prop:step-012-hidden-kernel}, if \(A_n\) is all-input
replacement-\((\varepsilon,\delta)\)-DP, then every \(K_i^{A_n}\),
\(i\in H\), is all-input replacement-\((\varepsilon,\delta)\)-DP and
hence replacement-\((0.1,\delta)\)-DP.  No realizability premise or
privacy composition is used.
\end{proposition}

\begin{proof}
Fix adjacent factor databases \(s,s'\), differing only in row
\(a\in[m_i]\), and one auxiliary realization \(\xi\).  On overflow, both
inputs return the same point mass.  Off overflow, if \(a>L_i\), the row is
unused and
\begin{equation}\label{eq:step-012-unused-row}
\Phi_{i,\xi}(s)=\Phi_{i,\xi}(s').
\end{equation}
If \(a\le L_i\), there is a unique requested slot
\begin{equation}\label{eq:step-012-unique-slot}
r_i(a)=\min\{r:\ell_{i,r}=a\}.
\end{equation}
The assembly map inserts row \(a\) there and nowhere else.  Thus the two
global datasets are equal or differ in exactly one labeled record, even
for inconsistent labels.

For each \(E\subseteq\mathcal V_i\), global privacy applied once to the
measurable event \(\rho_i^{-1}(E)\), together with equality in the other
branches, gives
\begin{equation}\label{eq:step-012-conditional-dp}
K_{i,\xi}(s,E)
\le e^\varepsilon K_{i,\xi}(s',E)+\delta.
\end{equation}
The law \(\mathsf Q_i\) is input-independent.  Integrating the same
conditional inequality yields
\begin{equation}\label{eq:step-012-hidden-dp}
\begin{split}
K_i^{A_n}(s,E)
&\le\int\bigl(e^\varepsilon K_{i,\xi}(s',E)+\delta\bigr)
\,\mathsf Q_i(d\xi)\\
&=e^\varepsilon K_i^{A_n}(s',E)+\delta.
\end{split}
\end{equation}
Exchanging \(s,s'\) supplies the reverse direction.  The additive term is
one \(\delta\), since this is input-independent mixing of one-call kernels,
not composition.  Finally \(\varepsilon\le0.1\) gives
\(e^\varepsilon\le e^{0.1}\).
\end{proof}

\begin{proposition}[Exact marginal identity coupling]
\label{prop:step-012-identity-coupling}
Under Assumption~\ref{assump:canonical-product},
Lemmas~\ref{lem:step-001-risk-pullback},
\ref{lem:step-010-finite-coding}, and
\ref{lem:step-012-ideal-experiment},
Propositions~\ref{prop:step-010-hard-prior},
\ref{prop:step-011-overflow-certificate},
\ref{prop:step-012-hidden-kernel}, and
\ref{prop:step-012-hidden-privacy}, fix \(i\in H\).  When
\(T_i=(c_i,D_i)\sim\nu_i\) and the hidden-kernel input is iid from
\((D_i)_{c_i}^{m_i}\), its experiment can be coupled to the common ideal
experiment so that, on \(\mathcal O_i^c\), the task vector, global target,
mixture distribution, global dataset, learner-output coordinate, finite
restriction, and factor risk are identical.  If \(R_i^{\rm tr}\) and
\(R_i^{\rm id}\) are the truncated and ideal risks, then
\begin{equation}\label{eq:step-012-pointwise-coupling}
|R_i^{\rm tr}-R_i^{\rm id}|
\le\mathbf1_{\mathcal O_i}
\end{equation}
and
\begin{equation}\label{eq:step-012-expected-coupling}
|\mathbb ER_i^{\rm tr}-\mathbb ER_i^{\rm id}|
\le\Pr(\mathcal O_i)\le\eta_0.
\end{equation}
The same holds conditional on any task in \(\operatorname{supp}\nu_i\).
This is a factor-marginal coupling, not a joint-factor independence claim.
\end{proposition}

\begin{proof}
Use the variables of Lemma~\ref{lem:step-012-ideal-experiment}.  Define
the factor input
\begin{equation}\label{eq:step-012-coupled-factor-input}
S_i^{\rm in}=(U_{i,1},\ldots,U_{i,m_i}).
\end{equation}
Conditional on \(T_i\), it has law \((D_i)_{c_i}^{m_i}\) and is independent
of \(T_{-i}\), the route, and the other-factor arrays, precisely the
auxiliary marginal required by the hidden kernel.  At a non-\(i\) slot use
\begin{equation}\label{eq:step-012-coupled-other-row}
W_r=U_{I_r,\ell_{I_r,r}}.
\end{equation}
On \(\mathcal O_i^c\), if \(I_r=i\),
\begin{equation}\label{eq:step-012-i-slot-identity}
\bigl(\Phi_{i,\xi}(S_i^{\rm in})\bigr)_r
=U_{i,\ell_{i,r}}=S_r^\star;
\end{equation}
if \(I_r\ne i\), the same follows from
\eqref{eq:step-012-coupled-other-row}.  Hence
\begin{equation}\label{eq:step-012-database-identity}
\Phi_{i,\xi}(S_i^{\rm in})=S^\star
\quad\text{on }\mathcal O_i^c.
\end{equation}
The task vector, target \(c^\star\), and mixture \(D^\star\) are literally
shared.  Full-product realizability is used here, but was not used in the
privacy proof.

Always sample \(\Omega^\star\) from
\eqref{eq:step-012-ideal-output}.  Define
\begin{equation}\label{eq:step-012-coupled-output-vectors}
V_i^{\rm id}=\rho_i(\Omega^\star),
\qquad
V_i^{\rm tr}=
\begin{cases}
\rho_i(\Omega^\star),&\mathcal O_i^c,\\
v_i^\circ,&\mathcal O_i.
\end{cases}
\end{equation}
The database identity makes this the correct hidden-kernel marginal off
overflow; the fallback branch gives the correct marginal on overflow.  No
random-seed representation of \(A_n\) is needed.

Set
\begin{equation}\label{eq:step-012-coupled-risks}
R_i^{\rm id}=R_{D_i}(h_{\Omega^\star}|_{X_i},c_i),
\qquad
R_i^{\rm tr}=R_{D_i}(e_i(V_i^{\rm tr}),c_i).
\end{equation}
Every task in the prior is supported on \(F_i\), and finite coding gives
\begin{equation}\label{eq:step-012-evaluation-identity}
e_i(\rho_i(\omega))(x_{i,a})=h_\omega(x_{i,a}),
\qquad a\in[t_i].
\end{equation}
Therefore \(R_i^{\rm tr}=R_i^{\rm id}\) off overflow.  Both risks are
measurable finite sums in \([0,1]\), so their difference is at most one on
overflow, proving \eqref{eq:step-012-pointwise-coupling}.  Expectation and
Proposition~\ref{prop:step-011-overflow-certificate} give
\eqref{eq:step-012-expected-coupling}.  The route is task-independent, so
conditioning on a fixed supported task does not change the bound.
\end{proof}

\begin{proposition}[Boundary traces and the one-factor specialization]
\label{prop:step-012-boundaries}
Under Propositions~\ref{prop:step-012-hidden-kernel},
\ref{prop:step-012-hidden-privacy}, and
\ref{prop:step-012-identity-coupling}, the following hold:
\begin{enumerate}
\item If \(L_i=0\), the factor input is never read and ideal/truncated
identity is exact.
\item Replacing an unused row changes no global row; replacing a used row
changes only its unique requested slot.
\item On overflow, output is input-independent and selected before an
unavailable row is requested.
\item Privacy covers every nonrealizable factor input.
\item If \(H=\{i\}\), the construction remains valid with no other active
prior and fixed tasks on all low factors.
\item At the algebraic endpoint \(\pi_i=0\), the construction reads no
input and has zero overflow and zero discrepancy.
\item If \(k=1\) on the active branch, then \(L_1=n\le m_{n,1}\), the
first \(n\) input rows are inserted once, and the coupling loss is zero.
\end{enumerate}
\end{proposition}

\begin{proof}
If \(L_i=0\), the assembly map contains no factor-input coordinate and
still equals the ideal database.  The used, unused, overflow, and
nonrealizable cases are exactly the exhaustive branches in
\eqref{eq:step-012-unused-row}--\eqref{eq:step-012-hidden-dp}.  If only
one factor is active, the product over other active priors is the unique
empty-product law, while \eqref{eq:step-012-low-reference-task} supplies
all low-factor tasks.  At \(\pi_i=0\), the route gives \(L_i=0\) surely,
and the overflow certificate gives probability zero.

For \(k=1\), \(\pi_1=1\), the route is deterministic, and
\begin{equation}\label{eq:step-012-one-factor-budget}
L_1=n,
\qquad
m_{n,1}=\max\{8,\lceil4n\rceil\}
=\max\{8,4n\}\ge n.
\end{equation}
Hence \(\mathcal O_1=\varnothing\).  There are no other tasks or rows;
input row \(a\) is inserted at global slot \(a\) for \(1\le a\le n\),
and all remaining input rows are unused.  The coupling identities hold
everywhere, so
\begin{equation}\label{eq:step-012-one-factor-zero-loss}
|\mathbb ER_1^{\rm tr}-\mathbb ER_1^{\rm id}|=0.
\end{equation}
\end{proof}

\subsection{Product-prior tensorization of exact risk}
\label{app:step-013}

Fix one lower-bound candidate and let \(H\), \(w_L\), and the active
hard priors \((\nu_i)_{i\in H}\) be as in
Lemma~\ref{lem:step-009-low-mass} and
Proposition~\ref{prop:step-010-hard-prior}.  For every \(j\notin H\),
let \(T_j^\circ\) be the deterministic realizable task used in
Lemma~\ref{lem:step-012-ideal-experiment}, and set
\begin{equation}\label{eq:step-013-product-task-law}
\lambda_j=
\begin{cases}
\nu_j,&j\in H,\\
\delta_{T_j^\circ},&j\notin H,
\end{cases}
\qquad
\Lambda=\bigotimes_{j=1}^k\lambda_j.
\end{equation}
The finite family of priors, and hence \(\Lambda\), is fixed before the
global learner.  After an arbitrary global learner \(A_n\) is fixed,
write \(\mathbb P_\star\) for the common experiment of
Lemma~\ref{lem:step-012-ideal-experiment}: draw all factor tasks from
\(\Lambda\), form the full target \(c^\star\) and mixture \(D^\star\),
draw the iid global sample \(S^\star\), and finally draw
\(\Omega^\star\sim A_n(S^\star,\cdot)\).  Let
\begin{equation}\label{eq:step-013-common-risks}
R_i^\star
=R_{D_i}(h_{\Omega^\star}|_{X_i},c_i),
\qquad
R^\star
=R_{D^\star}(h_{\Omega^\star},c^\star).
\end{equation}
These are measurable \([0,1]\)-valued random variables under the one
common law \(\mathbb P_\star\); write \(\mathbb E_\star\) for
expectation under that law.

\begin{lemma}[Active marginal floors in one common experiment]
\label{lem:step-013-common-marginal-floors}
Under Assumption~\ref{assump:canonical-product},
Proposition~\ref{prop:step-010-hard-prior},
Lemma~\ref{lem:step-012-ideal-experiment}, and
Propositions~\ref{prop:step-012-hidden-kernel},
\ref{prop:step-012-hidden-privacy}, and
\ref{prop:step-012-identity-coupling}, fix all active priors before an
arbitrary global learner and define the common law and risks by
\eqref{eq:step-013-product-task-law}--\eqref{eq:step-013-common-risks}.
Then, for every \(i\in H\),
\begin{equation}\label{eq:step-013-marginal-floor}
\mathbb E_\star R_i^\star>\frac18-\eta_0.
\end{equation}
Every expectation in \eqref{eq:step-013-marginal-floor} is a marginal
of the same common experiment; no independence among factor outputs is
required.
\end{lemma}

\begin{proof}
Fix \(i\in H\).  Propositions~\ref{prop:step-012-hidden-kernel} and
\ref{prop:step-012-hidden-privacy} construct, from the fixed common task
laws and the arbitrary global learner, a total unrestricted
replacement-\((0.1,\delta)\)-DP factor learner \(K_i^{A_n}\) on exactly
\(m_{n,i}\) rows.  This learner lies in the universal learner
quantifier of Proposition~\ref{prop:step-010-hard-prior}.  Define
\begin{equation}\label{eq:step-013-truncated-risk}
R_i^{\rm tr}
=R_{D_i}\bigl(e_i(K_i^{A_n}(S_i)),c_i\bigr),
\qquad
(c_i,D_i)\sim\nu_i,
\quad
S_i\sim(D_i)_{c_i}^{m_{n,i}}.
\end{equation}
The hard-prior conclusion gives the strict inequality
\begin{equation}\label{eq:step-013-truncated-floor}
\mathbb E R_i^{\rm tr}>\frac18.
\end{equation}

Proposition~\ref{prop:step-012-identity-coupling} identifies the ideal
counterpart of \eqref{eq:step-013-truncated-risk} with the \(i\)-th
marginal of \(\mathbb P_\star\).  Its same-target coupling, together
with Proposition~\ref{prop:step-011-overflow-certificate}, gives
\begin{equation}\label{eq:step-013-marginal-coupling}
\left|\mathbb E R_i^{\rm tr}
      -\mathbb E_\star R_i^\star\right|
\le \Pr(\mathcal O_i)
\le\eta_0.
\end{equation}
Consequently,
\begin{equation}\label{eq:step-013-floor-transfer}
\mathbb E_\star R_i^\star
\ge \mathbb E R_i^{\rm tr}-\eta_0
>\frac18-\eta_0.
\end{equation}
Thus the ideal/truncated discrepancy is charged exactly once.  The
finite evaluation map, task restriction, construction of the full target, and raw
decoding create no additional residual by
Proposition~\ref{prop:step-012-identity-coupling}.

The same marginal argument applies to every active \(i\), but each
right-hand variable \(R_i^\star\) remains the factor risk on the
already fixed common law \(\mathbb P_\star\).  The hidden kernels need
not be released or run jointly, so this argument never combines
inequalities from different ideal experiments.
\end{proof}

\begin{lemma}[Exact disjoint-risk sum in the common experiment]
\label{lem:step-013-exact-risk-sum}
Under Assumption~\ref{assump:canonical-product},
Lemma~\ref{lem:step-001-risk-pullback}, and
Lemma~\ref{lem:step-012-ideal-experiment}, the risks in
\eqref{eq:step-013-common-risks} satisfy, on every outcome of the common
experiment,
\begin{equation}\label{eq:step-013-pointwise-risk-sum}
R^\star=\sum_{i=1}^k\pi_iR_i^\star.
\end{equation}
Consequently,
\begin{equation}\label{eq:step-013-expected-risk-sum}
\mathbb E_\star R^\star
=\sum_{i=1}^k\pi_i\mathbb E_\star R_i^\star.
\end{equation}
Neither equality requires independence among the restrictions of the
global learner output.
\end{lemma}

\begin{proof}
Fix one outcome of the finite task draw and one decoded global output
\(h_{\Omega^\star}\).  Lemma~\ref{lem:step-012-ideal-experiment} gives
\begin{equation}\label{eq:step-013-common-target-mixture}
c^\star|_{X_i}=c_i,
\qquad
D^\star(B)=\sum_{i=1}^k\pi_iD_i(B\cap X_i),
\end{equation}
where \(D_i\) is supported on \(X_i\).  The mistake set is the disjoint
union
\begin{equation}\label{eq:step-013-mistake-union}
\{x:h_{\Omega^\star}(x)\ne c^\star(x)\}
=\bigsqcup_{i=1}^k
\{x\in X_i:h_{\Omega^\star}(x)\ne c_i(x)\}.
\end{equation}
Applying \eqref{eq:step-013-common-target-mixture} to
\eqref{eq:step-013-mistake-union}, equivalently applying
Lemma~\ref{lem:step-001-risk-pullback}, yields
\begin{align}
R^\star
&=D^\star\{h_{\Omega^\star}\ne c^\star\}\notag\\
&=\sum_{i=1}^k\pi_i
  D_i\{x\in X_i:h_{\Omega^\star}(x)\ne c_i(x)\}\notag\\
&=\sum_{i=1}^k\pi_iR_i^\star.
\label{eq:step-013-expanded-risk-sum}
\end{align}
This is an equality for the realized joint output, not merely an
equality in distribution or an approximation.  Arbitrary correlation
among the restrictions of \(h_{\Omega^\star}\) is present on both
sides.  Because every summand is measurable and bounded and \(k<\infty\),
finite linearity of expectation applied to
\eqref{eq:step-013-expanded-risk-sum} proves
\eqref{eq:step-013-expected-risk-sum}.  There is no infinite
interchange or conditioning conversion.
\end{proof}

\begin{proposition}[Common-prior weighted tensorization]
\label{prop:step-013-tensorization}
Under Lemma~\ref{lem:step-009-low-mass} and
Lemmas~\ref{lem:step-013-common-marginal-floors} and
\ref{lem:step-013-exact-risk-sum}, the one common product-prior ideal
experiment satisfies
\begin{equation}\label{eq:step-013-tensorized-floor}
\mathbb E_\star
R_{D^\star}(h_{\Omega^\star},c^\star)
>(1-w_L)\left(\frac18-\eta_0\right).
\end{equation}
The only losses are the low-factor mass \(w_L\) and one marginal
residual \(\eta_0\) per unit of active weight.
\end{proposition}

\begin{proof}
All factor risks are nonnegative.  Lemma~\ref{lem:step-013-exact-risk-sum}
therefore gives
\begin{equation}\label{eq:step-013-drop-low-factors}
\mathbb E_\star R^\star
=\sum_{i=1}^k\pi_i\mathbb E_\star R_i^\star
\ge\sum_{i\in H}\pi_i\mathbb E_\star R_i^\star.
\end{equation}
Apply \eqref{eq:step-013-truncated-floor} and
\eqref{eq:step-013-marginal-coupling} before summing.  Since every
active weight is positive,
\begin{align}
\sum_{i\in H}\pi_i\mathbb E_\star R_i^\star
&\ge
  \sum_{i\in H}\pi_i
  \bigl(\mathbb E R_i^{\rm tr}-\eta_0\bigr)
\notag\\
&>
  \sum_{i\in H}\pi_i
  \left(\frac18-\eta_0\right).
\label{eq:step-013-weighted-floors}
\end{align}
The first line charges factor \(i\)'s residual as
\(\pi_i\Pr(\mathcal O_i)\le\pi_i\eta_0\).  Hence the complete residual
is
\begin{equation}\label{eq:step-013-weighted-residual}
\sum_{i\in H}\pi_i\eta_0
=\eta_0\sum_{i\in H}\pi_i,
\end{equation}
not the probability of a union of overflow events and not an unweighted
multiple of \(\eta_0\).

Lemma~\ref{lem:step-009-low-mass} gives
\begin{equation}\label{eq:step-013-active-mass}
\sum_{i\in H}\pi_i
=1-\sum_{i\notin H}\pi_i
=1-w_L.
\end{equation}
It also gives \(H\ne\varnothing\), while every
\(\pi_i=s_i/M\) is positive.  Thus the second line of
\eqref{eq:step-013-weighted-floors} remains strict.  Combining
\eqref{eq:step-013-drop-low-factors}--\eqref{eq:step-013-active-mass}
proves \eqref{eq:step-013-tensorized-floor}.

The low-factor tasks were fixed before the learner and remain part of
the same target, distribution, sample, and output.  Their risks may be
arbitrarily correlated with all active risks and may be zero; only
their nonnegativity is used.  The proof also uses only marginal
expectations and never assumes independence of factor outputs or
overflow events.
\end{proof}

\begin{proposition}[Exact one-factor active-branch floor]
\label{prop:step-013-one-factor-baseline}
Under Lemma~\ref{lem:step-009-low-mass},
Proposition~\ref{prop:step-010-hard-prior},
Lemma~\ref{lem:step-011-zero-overflow},
Propositions~\ref{prop:step-012-identity-coupling} and
\ref{prop:step-012-boundaries}, and
Lemma~\ref{lem:step-013-exact-risk-sum}, suppose \(k=1\) on the active
fixed-candidate contradiction branch.  Then
\begin{equation}\label{eq:step-013-one-factor-floor}
\mathcal O_1=\varnothing\quad\text{almost surely},
\qquad
\mathbb E_\star R^\star>\frac18.
\end{equation}
Thus the one-factor baseline has no overflow or coupling loss.
\end{proposition}

\begin{proof}
Lemma~\ref{lem:step-009-low-mass} gives \(H\ne\varnothing\) on the
contradiction branch.  When \(k=1\), this forces
\begin{equation}\label{eq:step-013-one-factor-active-set}
H=\{1\},
\qquad
\pi_1=1,
\qquad
w_L=0.
\end{equation}
Lemma~\ref{lem:step-011-zero-overflow} gives \(L_1=n\) and
\begin{equation}\label{eq:step-013-one-factor-budget}
m_{n,1}=\max\{8,4n\}\ge n,
\end{equation}
so \(\mathcal O_1=\varnothing\) exactly.
Proposition~\ref{prop:step-012-boundaries} then shows that the one-use
learner inserts the first \(n\) factor rows once and that the truncated
and ideal risks coincide pointwise:
\begin{equation}\label{eq:step-013-one-factor-risk-identity}
R_1^{\rm tr}=R_1^\star.
\end{equation}
Proposition~\ref{prop:step-010-hard-prior} gives
\(\mathbb E R_1^{\rm tr}>1/8\).  Taking expectations in
\eqref{eq:step-013-one-factor-risk-identity} and applying
Lemma~\ref{lem:step-013-exact-risk-sum} with \(\pi_1=1\) gives
\begin{equation}\label{eq:step-013-one-factor-chain}
\mathbb E_\star R^\star
=\mathbb E_\star R_1^\star
=\mathbb E R_1^{\rm tr}
>\frac18.
\end{equation}
If the sole factor were outside \(H\), then \(H=\varnothing\), which
Lemma~\ref{lem:step-009-low-mass} already excludes on this branch.
\end{proof}

\begin{proof}[Proof of the product-prior tensorization conclusion]
Proposition~\ref{prop:step-010-hard-prior} fixes every active prior
independently of the learner, and
Lemma~\ref{lem:step-012-ideal-experiment} places their finite product
and the deterministic low-factor tasks in one experiment before data
or learner randomness.  For each active factor,
Propositions~\ref{prop:step-012-hidden-kernel} and
\ref{prop:step-012-hidden-privacy} supply an eligible unrestricted
factor learner.  Lemma~\ref{lem:step-013-common-marginal-floors} then
transfers that factor's hard-prior floor to the corresponding marginal
of the same common experiment, with sole residual \(\eta_0\).

Lemma~\ref{lem:step-013-exact-risk-sum} gives the pointwise finite
weighted risk identity, and
Proposition~\ref{prop:step-013-tensorization} combines it with the
active-mass identity from Lemma~\ref{lem:step-009-low-mass} to obtain
\eqref{eq:step-013-tensorized-floor}.  Equation
\eqref{eq:step-013-weighted-residual} records that each marginal loss
is charged once by its weight, with no union bound or independence
assumption.  Finally,
Proposition~\ref{prop:step-013-one-factor-baseline} gives the exact
zero-overflow floor \(>1/8\) when \(k=1\) on the active branch.
\end{proof}

\subsection{Candidate-wise unrestricted PAC lower bound}
\label{app:step-014}

Retain the common finite product-task law \(\Lambda\) from
\eqref{eq:step-013-product-task-law}, and let
\begin{equation}\label{eq:step-014-task-support}
\mathcal T=\operatorname{supp}(\Lambda)
\end{equation}
be its finite set of positive-mass atoms.  For
\(\tau=((c_i^\tau,D_i^\tau))_{i=1}^k\in\mathcal T\),
Lemma~\ref{lem:step-012-ideal-experiment} gives the full-product target
\(c^\tau\in C\), characterized by
\(c^\tau|_{X_i}=c_i^\tau\), and the probability measure
\begin{equation}\label{eq:step-014-atom-mixture}
D^\tau(B)=\sum_{i=1}^k
\pi_iD_i^\tau(B\cap X_i),
\qquad B\in\Sigma.
\end{equation}
Because every \(D_i^\tau\) is supported on \(X_i\),
\begin{equation}\label{eq:step-014-atom-masses}
D^\tau(X_i)=\pi_i=\omega_i,
\qquad
\sum_{i=1}^k\pi_i=1.
\end{equation}
For a fixed learner and candidate, define
\begin{equation}\label{eq:step-014-atom-loss}
L_\tau
=R_{D^\tau}(h_{\Omega_\tau},c^\tau),
\qquad
S_\tau\sim(D^\tau_{c^\tau})^n,
\quad
\Omega_\tau\sim A_n(S_\tau,\cdot).
\end{equation}
The random variable \(L_\tau\) is measurable and belongs to \([0,1]\).
The target, mixture, iid sample, and learner output in
\eqref{eq:step-014-atom-loss} are precisely the conditional objects of
the common experiment.  Since \(\mathcal T\) is finite and each of its
atoms has positive mass, Lemma~\ref{lem:step-013-exact-risk-sum} and
finite conditioning give
\begin{equation}\label{eq:step-014-finite-conditioning}
\mathbb E_\star R^\star
=\sum_{\tau\in\mathcal T}
\Lambda(\tau)\,\mathbb E L_\tau.
\end{equation}

\begin{lemma}[Exact PAC ceiling under finite task averaging]
\label{lem:step-014-finite-prior-ceiling}
Under the measurable binary-risk convention,
Lemma~\ref{lem:step-007-pac-conversion},
Lemma~\ref{lem:step-012-ideal-experiment}, and
Lemma~\ref{lem:step-013-exact-risk-sum}, fix one candidate \(n\), one
learner \(A_n\), and define the finite common experiment and atom
losses by \eqref{eq:step-014-task-support}--%
\eqref{eq:step-014-finite-conditioning}.  If every deterministic
\(\tau\in\mathcal T\) satisfies
\begin{equation}\label{eq:step-014-atom-pac-premise}
\Pr[L_\tau\le1/16]\ge\frac{15}{16},
\end{equation}
then
\begin{equation}\label{eq:step-014-taskwise-ceiling}
\mathbb E L_\tau\le\frac{31}{256}
\quad(\tau\in\mathcal T),
\qquad
\mathbb E_\star R^\star\le\frac{31}{256}.
\end{equation}
For each fixed atom, one also has the strict implication
\begin{equation}\label{eq:step-014-strict-event-conversion}
\mathbb E L_\tau>\frac{31}{256}
\quad\Longrightarrow\quad
\Pr[L_\tau>1/16]>\frac1{16}.
\end{equation}
\end{lemma}

\begin{proof}
Fix \(\tau\in\mathcal T\) and put
\begin{equation}\label{eq:step-014-bad-probability}
p_\tau=\Pr[L_\tau>1/16].
\end{equation}
This event is exactly the complement of the event in
\eqref{eq:step-014-atom-pac-premise}; loss equal to \(1/16\) remains on
the good side.  Since \(0\le L_\tau\le1\),
\begin{align}
\mathbb E L_\tau
&=\mathbb E[L_\tau\mathbf1\{L_\tau\le1/16\}]
  +\mathbb E[L_\tau\mathbf1\{L_\tau>1/16\}]
\notag\\
&\le\frac1{16}(1-p_\tau)+p_\tau
=\frac1{16}+\frac{15}{16}p_\tau.
\label{eq:step-014-bounded-loss-envelope}
\end{align}
Under \eqref{eq:step-014-atom-pac-premise},
\(p_\tau\le1/16\), and hence
\begin{equation}\label{eq:step-014-exact-ceiling-arithmetic}
\mathbb E L_\tau
\le\frac1{16}+\frac{15}{16}\frac1{16}
=\frac{16}{256}+\frac{15}{256}
=\frac{31}{256}.
\end{equation}
This is also the direct current-notation instance of
Lemma~\ref{lem:step-007-pac-conversion}.  Taking the strict
contrapositive of \eqref{eq:step-014-exact-ceiling-arithmetic} proves
\eqref{eq:step-014-strict-event-conversion}.  In particular, equality
\(p_\tau=1/16\) produces only the required non-strict ceiling.

Average \eqref{eq:step-014-exact-ceiling-arithmetic} with
\eqref{eq:step-014-finite-conditioning}.  Because \(\mathcal T\) is
finite and \(\Lambda\) is a probability law,
\begin{equation}\label{eq:step-014-prior-ceiling}
\mathbb E_\star R^\star
\le\sum_{\tau\in\mathcal T}
\Lambda(\tau)\frac{31}{256}
=\frac{31}{256}.
\end{equation}
No infinite interchange, measurable selection, or change from
taskwise probability to the common prior average occurs.
\end{proof}

\begin{lemma}[Strict rational separation of the tensorized floor]
\label{lem:step-014-rational-gap}
Under Lemma~\ref{lem:step-009-low-mass} and
Proposition~\ref{prop:step-013-tensorization}, fix the same candidate,
learner, and common experiment on the branch
\(n<c_{\rm low}M_{\oplus}(C)\).  If
\begin{equation}\label{eq:step-014-defect-bounds}
w_L<\frac1{128},
\qquad
\eta_0<\frac3{2048},
\end{equation}
then
\begin{equation}\label{eq:step-014-rational-gap}
\mathbb E_\star R^\star
>\frac{32131}{262144}
>\frac{31}{256}.
\end{equation}
\end{lemma}

\begin{proof}
The strict bounds in \eqref{eq:step-014-defect-bounds} give
\begin{equation}\label{eq:step-014-low-mass-complement}
1-w_L>\frac{127}{128}>0
\end{equation}
and
\begin{equation}\label{eq:step-014-overflow-complement}
\frac18-\eta_0
>\frac18-\frac3{2048}
=\frac{256-3}{2048}
=\frac{253}{2048}>0.
\end{equation}
Proposition~\ref{prop:step-013-tensorization}, followed by
\eqref{eq:step-014-low-mass-complement} and
\eqref{eq:step-014-overflow-complement}, yields
\begin{align}
\mathbb E_\star R^\star
&>(1-w_L)\left(\frac18-\eta_0\right)\notag\\
&>\frac{127}{128}\frac{253}{2048}
=\frac{127\cdot253}{128\cdot2048}
=\frac{32131}{262144}.
\label{eq:step-014-product-gap}
\end{align}
The final comparison is exact integer arithmetic:
\begin{equation}\label{eq:step-014-integer-gap}
\frac{31}{256}
=\frac{31\cdot1024}{262144}
=\frac{31744}{262144},
\qquad
32131-31744=387>0.
\end{equation}
Thus
\begin{equation}\label{eq:step-014-positive-margin}
\frac{32131}{262144}-\frac{31}{256}
=\frac{387}{262144}>0.
\end{equation}
Every comparison is strict; no decimal estimate or asymptotic
absorption is used.
\end{proof}

\begin{proposition}[Deterministic PAC-failure atom]
\label{prop:step-014-deterministic-atom}
Under Lemma~\ref{lem:step-012-ideal-experiment},
Lemma~\ref{lem:step-013-exact-risk-sum}, and
Lemma~\ref{lem:step-014-finite-prior-ceiling}, define
\(\mathcal T,\Lambda,L_\tau\) by
\eqref{eq:step-014-task-support}--%
\eqref{eq:step-014-finite-conditioning}.  If
\begin{equation}\label{eq:step-014-average-above-ceiling}
\mathbb E_\star R^\star>\frac{31}{256},
\end{equation}
then some positive-mass atom \(\tau_\star\in\mathcal T\) has a
deterministic full-product target \(c^{\tau_\star}\in C\) and the legal
probability measure
\begin{equation}\label{eq:step-014-selected-mixture}
D^{\tau_\star}(B)
=\sum_{i=1}^k\pi_i
D_i^{\tau_\star}(B\cap X_i)
\end{equation}
such that \(D^{\tau_\star}(X_i)=\pi_i=\omega_i\) for every \(i\), and
\begin{equation}\label{eq:step-014-selected-pac-failure}
\Pr_{\substack{S\sim(D^{\tau_\star}_{c^{\tau_\star}})^n\\
                \Omega\sim A_n(S,\cdot)}}
\left[
R_{D^{\tau_\star}}(h_\Omega,c^{\tau_\star})
>\frac1{16}
\right]
>\frac1{16}.
\end{equation}
\end{proposition}

\begin{proof}
If every atom obeyed
\(\mathbb E L_\tau\le31/256\), then
\eqref{eq:step-014-finite-conditioning} would imply
\(\mathbb E_\star R^\star\le31/256\), contrary to
\eqref{eq:step-014-average-above-ceiling}.  Hence a positive-mass
\(\tau_\star\in\mathcal T\) satisfies
\begin{equation}\label{eq:step-014-selected-expectation}
\mathbb E L_{\tau_\star}>\frac{31}{256}.
\end{equation}
The implication \eqref{eq:step-014-strict-event-conversion} in
Lemma~\ref{lem:step-014-finite-prior-ceiling} gives
\eqref{eq:step-014-selected-pac-failure}, with exactly the same
candidate, sample law, learner kernel, and binary risk.

Lemma~\ref{lem:step-012-ideal-experiment} constructs
\(c^{\tau_\star}\) by full Cartesian construction, so it belongs to
\(C\), rather than merely to a factorwise relaxation.  The same lemma
and \eqref{eq:step-014-atom-mixture}--%
\eqref{eq:step-014-atom-masses} show that
\eqref{eq:step-014-selected-mixture} is a probability measure with the
required block masses.  Lemma~\ref{lem:step-013-exact-risk-sum}
confirms that the selected loss is the exact global risk of this target
and mixture.  The mixture is legal on the selected factor measures'
actual supports and requires no product-distribution, public-knowledge,
quotient-output, finite-cardinality, or properness condition.  Finite
support describes this witness's provenance, not a restriction on the
theorem's distribution family or learner.
\end{proof}

\begin{proposition}[Fixed-candidate lower closure]
\label{prop:step-014-candidate-closure}
Under Assumption~\ref{assump:candidate-delta-budget},
Lemma~\ref{lem:step-007-pac-conversion},
Lemma~\ref{lem:step-009-low-mass},
Proposition~\ref{prop:step-009-almm-eligibility},
Proposition~\ref{prop:step-013-tensorization},
Lemmas~\ref{lem:step-014-finite-prior-ceiling} and
\ref{lem:step-014-rational-gap}, and
Proposition~\ref{prop:step-014-deterministic-atom}, fix a candidate
\(n\in\mathbb N\) satisfying the two delta-budget conjuncts in
Assumption~\ref{assump:candidate-delta-budget}.  If an unrestricted
measurable replacement-\((\varepsilon,\delta)\)-DP learner \(A_n\)
satisfies the universal \((1/16,1/16)\)-PAC guarantee at this candidate,
then
\begin{equation}\label{eq:step-014-candidate-lower}
n\ge c_{\rm low}M_{\oplus}(C).
\end{equation}
On the temporary branch
\(n<c_{\rm low}M_{\oplus}(C)\), the contradiction is witnessed by the
deterministic target and mixture in
\eqref{eq:step-014-selected-mixture}--%
\eqref{eq:step-014-selected-pac-failure}.
\end{proposition}

\begin{proof}
Write \(M=M_{\oplus}(C)\), and suppose for contradiction that
\begin{equation}\label{eq:step-014-contradiction-branch}
n<c_{\rm low}M.
\end{equation}
The universal PAC guarantee and
\eqref{eq:step-014-contradiction-branch} discharge the local premises
of Lemma~\ref{lem:step-009-low-mass}.
Assumption~\ref{assump:candidate-delta-budget} discharges the exact
numerical premises of
Proposition~\ref{prop:step-009-almm-eligibility} at this same \(n\) and
its exact budgets \((m_{n,i})_i\).  In particular, both
\(\delta\le1/[n\log(n+1)]\) and every
\(\delta\le c_\delta/[m_{n,i}^2\log(m_{n,i}+1)]\) remain in force, and
equality in either cap is allowed.

The lower chain culminating in
Proposition~\ref{prop:step-013-tensorization} therefore applies to the
same \(A_n\), without changing its output space or restricting it to
proper or factorwise hypotheses.  Lemma~\ref{lem:step-014-rational-gap}
gives
\begin{equation}\label{eq:step-014-lower-side}
\mathbb E_\star R^\star
>\frac{32131}{262144}
>\frac{31}{256}.
\end{equation}
The universal PAC guarantee applies to every deterministic
\((c^\tau,D^\tau)\) in the same experiment, so
Lemma~\ref{lem:step-014-finite-prior-ceiling} gives
\begin{equation}\label{eq:step-014-upper-side}
\mathbb E_\star R^\star\le\frac{31}{256},
\end{equation}
contradicting \eqref{eq:step-014-lower-side}.  Equivalently,
Proposition~\ref{prop:step-014-deterministic-atom} extracts from
\eqref{eq:step-014-lower-side} a specific atom satisfying
\eqref{eq:step-014-selected-pac-failure}.  Thus
\eqref{eq:step-014-contradiction-branch} is impossible, and its
negation is exactly \eqref{eq:step-014-candidate-lower}.

There is no small-sample exception.  At \(n=1\), every denominator in
the two candidate delta caps is positive, and the same comparison
applies whenever \eqref{eq:step-014-contradiction-branch} is entered.
One active factor suffices because its positive weight preserves the
strict tensorized lower bound.  If \(k=1\),
Proposition~\ref{prop:step-013-one-factor-baseline} strengthens
\eqref{eq:step-014-lower-side} to
\begin{equation}\label{eq:step-014-one-factor-gap}
\mathbb E_\star R^\star>\frac18>\frac{31}{256}
\end{equation}
with zero overflow.  Lemma~\ref{lem:step-014-finite-prior-ceiling} and
Proposition~\ref{prop:step-014-deterministic-atom} then recover the
exact one-factor strict PAC failure event.  The inactive \(k=1\) branch
is excluded by Lemma~\ref{lem:step-009-low-mass} under the PAC and
contradiction premises.  Thus the unrestricted one-factor VC and
Littlestone baseline is not weakened to an expectation-only statement.
\end{proof}

\begin{proof}[Proof of the candidate-wise lower conclusion]
At the fixed candidate, Lemma~\ref{lem:step-009-low-mass} supplies the
universal \(c_{\rm low}\), the strict bound \(w_L<1/128\), and a
nonempty active set on the branch
\(n<c_{\rm low}M_{\oplus}(C)\).
Proposition~\ref{prop:step-009-almm-eligibility} verifies every exact
factor budget and both candidate delta caps at that same candidate.
Lemma~\ref{lem:step-012-ideal-experiment} supplies the finite common
experiment, full-product targets, legal block mixtures, iid samples,
and prior-before-learner order, while
Lemma~\ref{lem:step-013-exact-risk-sum} supplies exact global risk.

Proposition~\ref{prop:step-013-tensorization} and
Lemma~\ref{lem:step-014-rational-gap} give
\begin{equation}\label{eq:step-014-assembled-gap}
\mathbb E_\star R^\star
>\frac{127}{128}
  \left(\frac18-\frac3{2048}\right)
=\frac{32131}{262144}
>\frac{31}{256}.
\end{equation}
Lemma~\ref{lem:step-014-finite-prior-ceiling} gives the opposite
ceiling under the alleged universal PAC guarantee, and
Proposition~\ref{prop:step-014-deterministic-atom} removes the averaging
device by producing one deterministic full-product task whose strict
failure probability exceeds \(1/16\).
Proposition~\ref{prop:step-014-candidate-closure} composes these results
without switching candidate, learner, risk, or probability space and
proves \(n\ge c_{\rm low}M_{\oplus}(C)\).
\end{proof}

\subsection{Upper, lower, and conditional characterization}
\label{app:step-015}

Throughout this subsection, write
\begin{equation}\label{eq:step-015-public-constants}
M=M_{\oplus}(C),
\qquad
C_{\rm up}=65536,
\qquad
C_{\rm quota}=\max\left\{1,K_Y+\frac1{20}\right\}.
\end{equation}
The first quantity is defined by the setting; the two constants are
the universal constants in
Propositions~\ref{prop:step-005-pac-closure} and
\ref{prop:step-006-public-quota-bridge}.

\begin{proposition}[Measurable global upper theorem and public quota bound]
\label{prop:step-015-global-upper}
Under Assumptions~\ref{assump:canonical-product},
\ref{assump:vc-one-factors},
\ref{assump:countably-coded-evaluation}, and
\ref{assump:global-privacy-range}, together with
Lemmas~\ref{lem:step-001-output-measurability} and
\ref{lem:step-001-risk-pullback} and
Propositions~\ref{prop:step-003-product-kernel},
\ref{prop:step-003-joint-composition},
\ref{prop:step-005-pac-closure}, and
\ref{prop:step-006-public-quota-bridge}, the following holds without
Assumption~\ref{assump:candidate-delta-budget}.  For every integer
\begin{equation}\label{eq:step-015-upper-threshold}
n\ge\left\lceil C_{\rm up}Q_\oplus\right\rceil,
\end{equation}
the quotient-first routed rule \(\bar A_n^{\oplus,Q}\) is a measurable
Markov kernel into \((\mathcal H^\oplus,\mathscr H^\oplus)\), is
replacement-\((\varepsilon,\delta)\)-DP on all labeled inputs, and, for
every \(c\in C\) and every probability measure \(D\) on
\((X,\Sigma)\), satisfies
\begin{equation}\label{eq:step-015-upper-pac}
\Pr_{\substack{S\sim D_c^n\\
                \bar H\sim\bar A_n^{\oplus,Q}(S,\cdot)}}
\left[
R_D(h_{\bar H},c)\le\frac1{16}
\right]
\ge\frac{15}{16}.
\end{equation}
Consequently,
\begin{equation}\label{eq:step-015-upper-consequences}
\operatorname{SC}_{\varepsilon,\delta}(C)
\le\left\lceil C_{\rm up}Q_\oplus\right\rceil,
\qquad
Q_\oplus
\le C_{\rm quota}\frac{M}{\varepsilon}
\log^2\!\left(\frac{eM}{\varepsilon\delta}\right).
\end{equation}
\end{proposition}

\begin{proof}
Proposition~\ref{prop:step-003-product-kernel} supplies, at every fixed
\(n\), the tuple-valued Markov kernel in the statement, the measurable
decoder, and measurable fixed-\(c,D\) risk events.
Proposition~\ref{prop:step-003-joint-composition} supplies the pointwise
all-input privacy inequality for every measurable tuple event and
every measurable decoded postprocessing.  Neither result uses the
candidate delta condition.

Proposition~\ref{prop:step-005-pac-closure} applies under
\eqref{eq:step-015-upper-threshold} and gives
\eqref{eq:step-015-upper-pac}, jointly over the iid sample and all
learner randomness and uniformly over every target and
arbitrary-support \(D\).  Lemma~\ref{lem:step-001-risk-pullback}
identifies this as exact raw distributional risk rather than a quotient
surrogate, while
Lemma~\ref{lem:step-001-output-measurability} shows that the decoded
tuple is a legal unrestricted learner output.

Apply these conclusions at the integer
\(\lceil C_{\rm up}Q_\oplus\rceil\).  The set of successful sample
sizes in the definition of
\(\operatorname{SC}_{\varepsilon,\delta}(C)\) is nonempty and contains
that integer.  Its least element exists by well-ordering of
\(\mathbb N\), belongs to the same set, and is no larger, which proves
the first inequality in \eqref{eq:step-015-upper-consequences}.
Proposition~\ref{prop:step-006-public-quota-bridge} is exactly the
second inequality.  This establishes the upper clause for every
\(0<\delta<1\) in
Assumption~\ref{assump:global-privacy-range}, without using either
lower delta cap.
\end{proof}

\begin{proposition}[Candidate-wise unrestricted lower theorem]
\label{prop:step-015-candidate-lower}
Under Assumptions~\ref{assump:canonical-product},
\ref{assump:vc-one-factors},
\ref{assump:countably-coded-evaluation}, and
\ref{assump:global-privacy-range}, fix one integer \(n\ge1\) satisfying
Assumption~\ref{assump:candidate-delta-budget}.  Under
Propositions~\ref{prop:step-014-deterministic-atom} and
\ref{prop:step-014-candidate-closure}, every unrestricted measurable
replacement-\((\varepsilon,\delta)\)-DP learner \(A_n\) that satisfies
the universal \((1/16,1/16)\)-PAC guarantee at this candidate obeys
\begin{equation}\label{eq:step-015-candidate-lower}
n\ge c_{\rm low}M.
\end{equation}
On the contradicted branch \(n<c_{\rm low}M\), there is a deterministic
full-product target \(c\in C\) and a legal block-mixture probability
measure \(D\) such that
\begin{equation}\label{eq:step-015-lower-witness-masses}
D(X_i)=\omega_i\qquad(1\le i\le k)
\end{equation}
and
\begin{equation}\label{eq:step-015-lower-witness-event}
\Pr_{\substack{S\sim D_c^n\\
                \Omega\sim A_n(S,\cdot)}}
\left[R_D(h_\Omega,c)>\frac1{16}\right]
>\frac1{16}.
\end{equation}
\end{proposition}

\begin{proof}
The two conjuncts of
Assumption~\ref{assump:candidate-delta-budget} are
\begin{equation}\label{eq:step-015-global-delta-cap}
\delta\le\frac1{n\log(n+1)}
\end{equation}
and
\begin{equation}\label{eq:step-015-factor-delta-caps}
\delta\le
\frac{c_\delta}{m_{n,i}^2\log(m_{n,i}+1)}
\qquad(1\le i\le k).
\end{equation}
They are checked at the same candidate and its exact budgets; neither
is discarded or inferred from the other.
Proposition~\ref{prop:step-014-candidate-closure} applies verbatim to
the same \(n,A_n\), output convention, and exact risk and yields
\eqref{eq:step-015-candidate-lower}.  Its learner is not restricted to
be proper, factorwise, quotient-output, or computationally bounded.

On the contradiction branch,
Proposition~\ref{prop:step-014-deterministic-atom}, as used by
Proposition~\ref{prop:step-014-candidate-closure}, removes the finite
proof prior and supplies
\eqref{eq:step-015-lower-witness-masses}--%
\eqref{eq:step-015-lower-witness-event}.  The event is strict because
the PAC success event is the closed event \(R_D\le1/16\).
All candidate quantities remain fixed, so there is no asymptotic or
uniform-in-candidate inference.
\end{proof}

\begin{proposition}[Conditional sample-complexity sandwich]
\label{prop:step-015-conditional-sandwich}
Under Proposition~\ref{prop:step-015-global-upper}, define
\begin{equation}\label{eq:step-015-attained-sample-complexity}
n_*=\operatorname{SC}_{\varepsilon,\delta}(C).
\end{equation}
If
\begin{equation}\label{eq:step-015-attained-global-cap}
0<\delta\le\frac1{n_*\log(n_*+1)}
\end{equation}
and
\begin{equation}\label{eq:step-015-attained-factor-caps}
\delta\le
\frac{c_\delta}
{m_{n_*,i}^2\log(m_{n_*,i}+1)}
\qquad(1\le i\le k)
\end{equation}
both hold, then
Proposition~\ref{prop:step-015-candidate-lower} applies at this integer
and
\begin{equation}\label{eq:step-015-conditional-sandwich}
c_{\rm low}M
\le n_*
\le\left\lceil C_{\rm up}Q_\oplus\right\rceil
\le C_{\rm up}C_{\rm quota}
\frac{M}{\varepsilon}
\log^2\!\left(\frac{eM}{\varepsilon\delta}\right).
\end{equation}
If either \eqref{eq:step-015-attained-global-cap} or any inequality in
\eqref{eq:step-015-attained-factor-caps} fails, the upper conclusion of
Proposition~\ref{prop:step-015-global-upper} remains valid, but no lower
bound at \(n_*\) is asserted.
\end{proposition}

\begin{proof}
Proposition~\ref{prop:step-015-global-upper} shows that the defining
set for sample complexity is nonempty.  Therefore \(n_*\in\mathbb N\)
is finite and attained by an unrestricted private learner satisfying
the universal PAC guarantee.
Equations~\eqref{eq:step-015-attained-global-cap} and
\eqref{eq:step-015-attained-factor-caps} are exactly
Assumption~\ref{assump:candidate-delta-budget} at this attained
candidate.  Proposition~\ref{prop:step-015-candidate-lower} therefore
gives the first inequality in
\eqref{eq:step-015-conditional-sandwich}; the second comes from
Proposition~\ref{prop:step-015-global-upper}.

Every \(q_i\) is an integer, so \(Q_\oplus\in\mathbb N\), and
\(C_{\rm up}=65536\) is an integer.  Hence
\begin{equation}\label{eq:step-015-exact-ceiling-removal}
\left\lceil C_{\rm up}Q_\oplus\right\rceil
=C_{\rm up}Q_\oplus.
\end{equation}
Multiplying the quota bound in
Proposition~\ref{prop:step-006-public-quota-bridge} by the positive
constant \(C_{\rm up}\) gives
\begin{equation}\label{eq:step-015-final-rate-bridge}
C_{\rm up}Q_\oplus
\le C_{\rm up}C_{\rm quota}
\frac{M}{\varepsilon}
\log^2\!\left(\frac{eM}{\varepsilon\delta}\right),
\end{equation}
which proves the last inequality in
\eqref{eq:step-015-conditional-sandwich} without an untracked ceiling
term.  The order is essential: identify \(n_*\), check both numerical
conditions there, and only then invoke the lower result.  Failure of
that check cannot affect the independently proved arbitrary-\(\delta\)
upper clause.
\end{proof}

\begin{proposition}[Exact one-factor quotient-first upper reduction]
\label{prop:step-015-one-factor-upper}
Under the assumptions and results of
Proposition~\ref{prop:step-015-global-upper}, suppose \(k=1\).  Then the
routed rule is exactly the measurable quotient-first factor kernel on
an unpadded quota-length iid prefix, followed by the same measurable
quotient decoder.  More precisely,
\begin{equation}\label{eq:step-015-one-factor-upper-identities}
X_1=X,
\qquad
M=s_1,
\qquad
Q_\oplus=q_1,
\qquad
J_1(S)=n\quad(S\in Z^n),
\end{equation}
and
\begin{equation}\label{eq:step-015-one-factor-quota}
q_1=\left\lceil
K_Y\frac{M}{\varepsilon}
\log^2\!\left(\frac{eM}{\varepsilon\delta}\right)
\right\rceil
\le C_{\rm quota}\frac{M}{\varepsilon}
\log^2\!\left(\frac{eM}{\varepsilon\delta}\right).
\end{equation}
At every theorem sample size there is zero shortage and zero
quotient-to-raw risk residual, and the same learner directly satisfies
\begin{equation}\label{eq:step-015-one-factor-strong-utility}
\Pr\left[R_D(h_{\bar H_1},c)\le\frac1{64}\right]
\ge\frac{4095}{4096}.
\end{equation}
\end{proposition}

\begin{proof}
When \(k=1\), the sole block is the whole domain, giving
\(X_1=X\) and \(J_1(S)=n\) pointwise.  The definitions give
\(M=s_1\), \(Q_\oplus=q_1\), and the equality in
\eqref{eq:step-015-one-factor-quota}.  Since \(C_{\rm up}\ge1\), the
threshold \(n\ge\lceil C_{\rm up}q_1\rceil\) implies \(n\ge q_1\).
Thus the routed tuple has one coordinate and its input is the first
\(q_1\) real quotient records, with no padding.  By construction that
coordinate law is precisely the factor kernel
\(\bar A_1^{\rm Yan}\), rather than a new or relaxed learner.

Proposition~\ref{prop:step-003-product-kernel} gives its measurable
tuple form, and Lemma~\ref{lem:step-001-output-measurability} gives the
same quotient decoder.
Proposition~\ref{prop:step-003-joint-composition} uses only the sole
factor privacy cost, so there is no cross-factor composition.
Proposition~\ref{prop:step-005-pac-closure} bypasses weighted shortage
and the Markov relaxation in this specialization and gives
\eqref{eq:step-015-one-factor-strong-utility} directly.
Lemma~\ref{lem:step-001-risk-pullback} makes this exact raw risk for
every arbitrary-support \(D\), with zero quotient or reference
residual.  Finally,
Proposition~\ref{prop:step-006-public-quota-bridge} gives the inequality
in \eqref{eq:step-015-one-factor-quota}.  Thus this is the same
measurable quotient-first factor learner at the same quota order, with
only the setting-required measurable totalization and decoding made
explicit.
\end{proof}

\begin{proposition}[Exact one-factor unrestricted lower reduction]
\label{prop:step-015-one-factor-lower}
Under Assumption~\ref{assump:candidate-delta-budget} at one fixed
candidate and the assumptions and results of
Proposition~\ref{prop:step-015-candidate-lower}, suppose \(k=1\).  Then
\begin{equation}\label{eq:step-015-one-factor-lower-identities}
\omega_1=1,
\qquad
M=s_1=1+\log^*(d_1+1),
\qquad
m_{n,1}=\max\{8,4n\}.
\end{equation}
The lower mixture requests exactly \(n\) rows from factor one, so
\begin{equation}\label{eq:step-015-one-factor-zero-overflow}
n\le m_{n,1},
\qquad
\Pr[\text{factor-one overflow}]=0.
\end{equation}
Proposition~\ref{prop:step-014-candidate-closure} therefore yields
\begin{equation}\label{eq:step-015-one-factor-lower-rate}
n\ge c_{\rm low}\bigl(1+\log^*(d_1+1)\bigr),
\end{equation}
and, on the contradicted branch, the exact strict PAC-failure event for
the same arbitrary improper learner.  Both conditions
\eqref{eq:step-015-global-delta-cap}--%
\eqref{eq:step-015-factor-delta-caps} remain in force.
\end{proposition}

\begin{proof}
When \(k=1\), the definitions give \(M=s_1\),
\(\omega_1=s_1/M=1\), and
\(\lceil4n\omega_1\rceil=4n\), proving
\eqref{eq:step-015-one-factor-lower-identities}.  Every one of the
\(n\) global mixture rows belongs to factor one, so the one-use
simulation requests exactly \(n\) factor rows.  Since
\begin{equation}\label{eq:step-015-one-factor-budget-chain}
m_{n,1}=\max\{8,4n\}\ge4n\ge n,
\end{equation}
the overflow event is empty, which proves
\eqref{eq:step-015-one-factor-zero-overflow} exactly rather than by a
tail estimate.

The one-factor part of
Proposition~\ref{prop:step-014-candidate-closure} preserves both source
branches: bounded iterated-log complexity is covered by the ordinary
unrestricted VC floor, while the high-complexity branch uses the
unrestricted expected-risk Littlestone interface.  That proposition
also records that no low-factor mass is lost at weight one and converts
the strict expected-risk floor back to the event \(R_D>1/16\).
Substituting \(M=s_1\) proves
\eqref{eq:step-015-one-factor-lower-rate}.  The learner remains
arbitrary, improper, joint, and computationally unrestricted; no
factorwise, quotient-coded, finite-output, or efficiency condition is
introduced.
\end{proof}

\begin{proof}[Proof of the conditional characterization]
Proposition~\ref{prop:step-015-global-upper} composes the measurable
kernel, all-input privacy, arbitrary-distribution PAC, exact-risk, and
public quota interfaces.  It proves a finite attained
\(n_*=\operatorname{SC}_{\varepsilon,\delta}(C)\) for every
\(0<\delta<1\) permitted by
Assumption~\ref{assump:global-privacy-range}, without invoking
Assumption~\ref{assump:candidate-delta-budget}.

Proposition~\ref{prop:step-015-candidate-lower} separately retains both
\[
\delta\le\frac1{n\log(n+1)}
\quad\text{and}\quad
\delta\le
\frac{c_\delta}{m_{n,i}^2\log(m_{n,i}+1)}
\quad(1\le i\le k)
\]
at one fixed candidate, and exports the exact unrestricted lower
theorem and deterministic strict PAC-failure witness.  No claim is
made at a candidate where either condition fails.
Proposition~\ref{prop:step-015-conditional-sandwich} is the only point
where the scopes meet: it checks both conditions at the attained
\(n_*\) before applying the lower result.  Equations
\eqref{eq:step-015-exact-ceiling-removal} and
\eqref{eq:step-015-final-rate-bridge} give
\[
O\!\left(
\frac{M_{\oplus}(C)}{\varepsilon}
\log^2\!\frac{eM_{\oplus}(C)}{\varepsilon\delta}
\right)
\]
with universal hidden constant \(C_{\rm up}C_{\rm quota}\).  If the
candidate check fails, only the lower side is inactive.

Proposition~\ref{prop:step-015-one-factor-upper} proves the upper
baseline by object identity: one block, one quota, the same
quotient-first factor kernel, no padding, exact pullback risk, and the
stronger direct factor utility.
Proposition~\ref{prop:step-015-one-factor-lower} proves the lower
baseline by object and budget identity: factor weight one, all \(n\)
rows routed to it, budget at least \(4n\), zero overflow, and the
unrestricted VC/Littlestone dichotomy converted to the strict PAC
event.

The simultaneous universal constants are \(C_{\rm up}=65536\),
\(C_{\rm quota}=\max\{1,K_Y+1/20\}\), and \(c_{\rm low}>0\), independent
of the class, factors, candidate, distribution, and privacy parameters.
The resulting two-sided theorem remains conditional on the canonical
finite product and countable evaluation coding, and the lower
sample-complexity comparison remains conditional on its candidate
check.  It therefore does not assert a characterization for all
finite-Littlestone classes.
\end{proof}

\subsection{Proof of the Main Theorem}
\label{app:main-theorem-proof}

\begin{proof}[Proof of Theorem~\ref{thm:main}]
Proposition~\ref{prop:step-015-global-upper}, with the constants in
\eqref{eq:step-015-public-constants}, proves that the specified routed
rule is a measurable Markov kernel, is private on all labeled inputs,
and satisfies \eqref{eq:main-upper-pac} at every stated sample size.
The same proposition proves the first inequality in
\eqref{eq:main-upper-rate}, and its application of
Proposition~\ref{prop:step-006-public-quota-bridge} proves the second.
That proof uses only Assumptions~\ref{assump:canonical-product},
\ref{assump:vc-one-factors},
\ref{assump:countably-coded-evaluation}, and
\ref{assump:global-privacy-range}; in particular, it is valid for every
\(0<\delta<1\) in the theorem.

At a fixed candidate satisfying
Assumption~\ref{assump:candidate-delta-budget},
Proposition~\ref{prop:step-015-candidate-lower} gives
\eqref{eq:main-candidate-lower}.  Its deterministic witness is exactly
\eqref{eq:main-lower-witness}, with block masses
\(D(X_i)=\omega_i\), and its learner quantifier is the unrestricted
measurable one from Section~\ref{sec:preliminaries}.

Proposition~\ref{prop:step-015-global-upper} also makes
\(n_*=\operatorname{SC}_{\varepsilon,\delta}(C)\) finite and attained.
Under \eqref{eq:main-nstar-global-cap} and
\eqref{eq:main-nstar-factor-caps},
Proposition~\ref{prop:step-015-conditional-sandwich} applies at that
same integer and gives all three inequalities in
\eqref{eq:main-sandwich}.  Its exact ceiling calculation
\eqref{eq:step-015-exact-ceiling-removal} and rate bridge
\eqref{eq:step-015-final-rate-bridge} establish the displayed universal
constant \(C_{\rm up}C_{\rm quota}\).  The proposition also proves that
failure of either candidate check leaves the upper clause intact and
activates no lower conclusion.

Finally, Proposition~\ref{prop:step-015-one-factor-upper} proves the
object identities, absence of padding, exact risk pullback, and stronger
utility \eqref{eq:main-one-factor-upper} when \(k=1\).
Proposition~\ref{prop:step-015-one-factor-lower} proves
\(\omega_1=1\), the exact budget, zero overflow, the rate
\eqref{eq:main-one-factor-lower}, and the strict failure event for the
same unrestricted learner.  These four cited propositions retain the
probability, horizon, exact-risk, and constant-dependence modes stated
at the end of the theorem, completing every clause.
\end{proof}
